% Options for packages loaded elsewhere
% Options for packages loaded elsewhere
\PassOptionsToPackage{unicode}{hyperref}
\PassOptionsToPackage{hyphens}{url}
\PassOptionsToPackage{dvipsnames,svgnames,x11names}{xcolor}
%
\documentclass[
  english,
  letterpaper,
  DIV=11,
  numbers=noendperiod]{scrreprt}
\usepackage{xcolor}
\usepackage{amsmath,amssymb}
\setcounter{secnumdepth}{5}
\usepackage{iftex}
\ifPDFTeX
  \usepackage[T1]{fontenc}
  \usepackage[utf8]{inputenc}
  \usepackage{textcomp} % provide euro and other symbols
\else % if luatex or xetex
  \usepackage{unicode-math} % this also loads fontspec
  \defaultfontfeatures{Scale=MatchLowercase}
  \defaultfontfeatures[\rmfamily]{Ligatures=TeX,Scale=1}
\fi
\usepackage{lmodern}
\ifPDFTeX\else
  % xetex/luatex font selection
\fi
% Use upquote if available, for straight quotes in verbatim environments
\IfFileExists{upquote.sty}{\usepackage{upquote}}{}
\IfFileExists{microtype.sty}{% use microtype if available
  \usepackage[]{microtype}
  \UseMicrotypeSet[protrusion]{basicmath} % disable protrusion for tt fonts
}{}
\makeatletter
\@ifundefined{KOMAClassName}{% if non-KOMA class
  \IfFileExists{parskip.sty}{%
    \usepackage{parskip}
  }{% else
    \setlength{\parindent}{0pt}
    \setlength{\parskip}{6pt plus 2pt minus 1pt}}
}{% if KOMA class
  \KOMAoptions{parskip=half}}
\makeatother
% Make \paragraph and \subparagraph free-standing
\makeatletter
\ifx\paragraph\undefined\else
  \let\oldparagraph\paragraph
  \renewcommand{\paragraph}{
    \@ifstar
      \xxxParagraphStar
      \xxxParagraphNoStar
  }
  \newcommand{\xxxParagraphStar}[1]{\oldparagraph*{#1}\mbox{}}
  \newcommand{\xxxParagraphNoStar}[1]{\oldparagraph{#1}\mbox{}}
\fi
\ifx\subparagraph\undefined\else
  \let\oldsubparagraph\subparagraph
  \renewcommand{\subparagraph}{
    \@ifstar
      \xxxSubParagraphStar
      \xxxSubParagraphNoStar
  }
  \newcommand{\xxxSubParagraphStar}[1]{\oldsubparagraph*{#1}\mbox{}}
  \newcommand{\xxxSubParagraphNoStar}[1]{\oldsubparagraph{#1}\mbox{}}
\fi
\makeatother


\usepackage{longtable,booktabs,array}
\usepackage{calc} % for calculating minipage widths
% Correct order of tables after \paragraph or \subparagraph
\usepackage{etoolbox}
\makeatletter
\patchcmd\longtable{\par}{\if@noskipsec\mbox{}\fi\par}{}{}
\makeatother
% Allow footnotes in longtable head/foot
\IfFileExists{footnotehyper.sty}{\usepackage{footnotehyper}}{\usepackage{footnote}}
\makesavenoteenv{longtable}
\usepackage{graphicx}
\makeatletter
\newsavebox\pandoc@box
\newcommand*\pandocbounded[1]{% scales image to fit in text height/width
  \sbox\pandoc@box{#1}%
  \Gscale@div\@tempa{\textheight}{\dimexpr\ht\pandoc@box+\dp\pandoc@box\relax}%
  \Gscale@div\@tempb{\linewidth}{\wd\pandoc@box}%
  \ifdim\@tempb\p@<\@tempa\p@\let\@tempa\@tempb\fi% select the smaller of both
  \ifdim\@tempa\p@<\p@\scalebox{\@tempa}{\usebox\pandoc@box}%
  \else\usebox{\pandoc@box}%
  \fi%
}
% Set default figure placement to htbp
\def\fps@figure{htbp}
\makeatother



\ifLuaTeX
\usepackage[bidi=basic]{babel}
\else
\usepackage[bidi=default]{babel}
\fi
% get rid of language-specific shorthands (see #6817):
\let\LanguageShortHands\languageshorthands
\def\languageshorthands#1{}
\ifLuaTeX
  \usepackage[english]{selnolig} % disable illegal ligatures
\fi


\setlength{\emergencystretch}{3em} % prevent overfull lines

\providecommand{\tightlist}{%
  \setlength{\itemsep}{0pt}\setlength{\parskip}{0pt}}



 


\usepackage[most]{tcolorbox}
\newtcolorbox{info}{
  breakable,
  title=Info
}
\KOMAoption{captions}{tableheading}
\makeatletter
\@ifpackageloaded{tcolorbox}{}{\usepackage[skins,breakable]{tcolorbox}}
\@ifpackageloaded{fontawesome5}{}{\usepackage{fontawesome5}}
\definecolor{quarto-callout-color}{HTML}{909090}
\definecolor{quarto-callout-note-color}{HTML}{0758E5}
\definecolor{quarto-callout-important-color}{HTML}{CC1914}
\definecolor{quarto-callout-warning-color}{HTML}{EB9113}
\definecolor{quarto-callout-tip-color}{HTML}{00A047}
\definecolor{quarto-callout-caution-color}{HTML}{FC5300}
\definecolor{quarto-callout-color-frame}{HTML}{acacac}
\definecolor{quarto-callout-note-color-frame}{HTML}{4582ec}
\definecolor{quarto-callout-important-color-frame}{HTML}{d9534f}
\definecolor{quarto-callout-warning-color-frame}{HTML}{f0ad4e}
\definecolor{quarto-callout-tip-color-frame}{HTML}{02b875}
\definecolor{quarto-callout-caution-color-frame}{HTML}{fd7e14}
\makeatother
\makeatletter
\@ifpackageloaded{bookmark}{}{\usepackage{bookmark}}
\makeatother
\makeatletter
\@ifpackageloaded{caption}{}{\usepackage{caption}}
\AtBeginDocument{%
\ifdefined\contentsname
  \renewcommand*\contentsname{Table of contents}
\else
  \newcommand\contentsname{Table of contents}
\fi
\ifdefined\listfigurename
  \renewcommand*\listfigurename{List of Figures}
\else
  \newcommand\listfigurename{List of Figures}
\fi
\ifdefined\listtablename
  \renewcommand*\listtablename{List of Tables}
\else
  \newcommand\listtablename{List of Tables}
\fi
\ifdefined\figurename
  \renewcommand*\figurename{Figure}
\else
  \newcommand\figurename{Figure}
\fi
\ifdefined\tablename
  \renewcommand*\tablename{Table}
\else
  \newcommand\tablename{Table}
\fi
}
\@ifpackageloaded{float}{}{\usepackage{float}}
\floatstyle{ruled}
\@ifundefined{c@chapter}{\newfloat{codelisting}{h}{lop}}{\newfloat{codelisting}{h}{lop}[chapter]}
\floatname{codelisting}{Listing}
\newcommand*\listoflistings{\listof{codelisting}{List of Listings}}
\makeatother
\makeatletter
\makeatother
\makeatletter
\@ifpackageloaded{caption}{}{\usepackage{caption}}
\@ifpackageloaded{subcaption}{}{\usepackage{subcaption}}
\makeatother
\usepackage{bookmark}
\IfFileExists{xurl.sty}{\usepackage{xurl}}{} % add URL line breaks if available
\urlstyle{same}
\hypersetup{
  pdftitle={RM2 ODL Qualitative Project Resources},
  pdfauthor={Ashley Robertson \& Wil Toivo},
  pdflang={en},
  colorlinks=true,
  linkcolor={blue},
  filecolor={Maroon},
  citecolor={Blue},
  urlcolor={Blue},
  pdfcreator={LaTeX via pandoc}}


\title{RM2 ODL Qualitative Project Resources}
\author{Ashley Robertson \& Wil Toivo}
\date{17 April 2026}
\begin{document}
\maketitle

\renewcommand*\contentsname{Table of contents}
{
\hypersetup{linkcolor=}
\setcounter{tocdepth}{2}
\tableofcontents
}

\bookmarksetup{startatroot}

\chapter*{Overview}\label{overview}
\addcontentsline{toc}{chapter}{Overview}

\markboth{Overview}{Overview}

In this book, you will find the resources you need to complete the
qualitative project for your Research Methods 2 course.

We will work through the book step-by-step, linking to relevant sections
on Moodle. However, as the resources are made available from the start,
you will also have some flexibility to work at a pace that suits you.

You will become familiar with the book in Semester 1, as we will ask you
to look at certain sections to accompany the lecture content. However,
the book will become most relevant in Semester 2, when you complete your
larger assignments.

You will find links to each chapter at the left hand side of the page,
and then there will be links to each heading within a chapter at the
right hand side.

\part{Assessment Information}

\chapter{Data Skills (Individual)}\label{ch1_data_skills}

In this Chapter, we have included information about the Data Skills
Assessments in RM2 (ANOVA and Regression). You can use the menu at the
right hand side of this page to jump to the different sections.

\section{General Information}\label{general-information}

\begin{itemize}
\item
  Deadlines for this assessment can be found under `Deadlines' within
  the `Course Information' section on the RM2 Moodle.
\item
  In total this assessment is worth 10\% of your final course grade.
  Each data skills worksheet is worth 5\% and you will be awarded a mark
  out of 22.
\item
  The submission links (and where appropriate, files) will be made
  available at least one week before the deadline.
\end{itemize}

\section{Important information about R submissions (please read and
watch the
video!)}\label{important-information-about-r-submissions-please-read-and-watch-the-video}

Please note the following for the R submissions:

\begin{itemize}
\item
  you must use the R markdown file provided. If you submit a file that
  you made yourself, the submission will automatically receive an H.
\item
  you must not alter the R chunk itself (e.g.~adding/removing backticks,
  removing or altering t0x etc.) in the R markdown file. If you do, you
  will automatically receive an H for that question and it may result in
  an H for the whole assessment.
\item
  you must submit the completed .Rmd file. If you submit a blank .Rmd
  file, your submission will automatically receive an H.
\item
  you must not alter the name of the object that the task is asking you
  to create
\end{itemize}

Please see
\href{https://uofglasgow.zoom.us/rec/play/kGPffSDx2RT741ZEXDG28KaYFxRDyriV1cO-Tmy0VDcWEDn_2wcCeX853_1t8T9jzrOzEfxMFcgu6yCt.zBN1JhinbBQSRjzG?eagerLoadZvaPages=&accessLevel=meeting&canPlayFromShare=true&from=share_recording_detail&startTime=1725965232000&componentName=rec-play&originRequestUrl=https\%3A\%2F\%2Fuofglasgow.zoom.us\%2Frec\%2Fshare\%2FC7w5EqVNPod9aLKY4SwMtW67PhxNwBaS1GXjvRrOwvVffKGB3HqNler6at6ESWzP.eiXr0xmkpri2n-VX\%3FstartTime\%3D1725965232000}{this
video} for further details.

\section{Type of
Assessment/Structure}\label{type-of-assessmentstructure}

\begin{itemize}
\item
  You will be provided with an .Rmd file that you should complete and
  submit, like the homeworks you did in RM1. Some questions are
  dependent on each other, but most are independent.
\item
  Both worksheets have 19 questions each. Some questions are worth 1
  point and some are worth 2 points. You can see how much each task is
  worth on the worksheet.
\item
  The points on each worksheet will total to 22.
\item
  You will be presented with three types of questions: error mode, code
  writing and multiple choice
\item
  For \textbf{error mode}, instead of asking you to write the code
  yourself, we're going to provide you with code that has errors in it.
  This code will either not run and it will produce an error message, or
  it will run, but it won't produce what it is supposed to. In both
  cases you have to fix the code to produce the intended output. The
  purpose of these tasks is to resolve errors and learn to debug, so
  please don't ask on Teams about the errors you get on the error mode
  tasks.
\item
  For \textbf{code writing}, you will be given instructions on code that
  you should write. Provide the code that gives exactly the intended
  output (and only that). For some questions, you will be provided a
  data frame or a plot in the instructions file, and you will be asked
  to recreate that in the task. You should replace the NULL in the code
  chunk with your code.
\item
  For \textbf{multiple choice}, you will be asked questions about the
  data or interpreting your statistical output. Please replace the NULL
  in the code chunk with the number of the option you think is correct.
  Please only use a single number, do not use words, and do not put the
  number in quotation marks.
\end{itemize}

\section{Assessment Support}\label{assessment-support}

\begin{itemize}
\item
  Most of the tasks are based on code and functions you have seen in the
  data skills book
\item
  Each worksheet has a couple of questions on things you may not have
  seen before. This is either to use a new function, or to combine
  functions in a new way. These tasks are there to test your
  generalisation skills with R functions.
\item
  You can find instructions for how to complete the .Rmd files in the
  data skills book
\item
  For the multiple choice questions about interpreting statistical
  output, you may also need to revise some of the key concepts from your
  lectures.
\item
  Further information about feedback can be found in
  \hyperref[data_skills_feedback]{the Feedback Section of this chapter}
\end{itemize}

\section{How to do well in this
Assessment}\label{how-to-do-well-in-this-assessment}

\begin{itemize}
\item
  Use the data skills activities to guide you.
\item
  Read the question carefully and ensure that you provide exactly what
  is asked for (e.g., code or a single value), and only that.
\item
  Follow the instructions for how to fill out the .Rmd file carefully
  making sure you do not change anything in the file you shouldn't.
\item
  Ensure you are up to date with the R activities, as the data skills
  activities will assess your knowledge of these.
\item
  Before you submit, make sure that your .Rmd file knits -- if it does
  it means that what you have written is legal code (this is not to say
  that you will have written the correct code, it simply means you
  definitely haven't written any code that doesn't work).
\end{itemize}

\section{Common Mistakes}\label{common-mistakes}

\begin{itemize}
\item
  Changing the .Rmd file other than to provide answers and your GUID
  (e.g., deleting backticks, changing code chunk names, not using the
  file provided).
\item
  Failure to follow instructions (e.g., writing code when a single value
  was requested).
\item
  Including any illegal code in your .Rmd file, e.g., install.packages
  (you should never write code that would change something on someone
  else's machine, it causes issues and it's impolite!).
\item
  Changing object names
\end{itemize}

\section{How is the assessment related to the lectures for this
course?}\label{how-is-the-assessment-related-to-the-lectures-for-this-course}

\begin{itemize}
\item
  The data skills submissions assess your ability to wrangle and
  visualise data in an open and reproducible way, as discussed in the
  lectures.
\item
  Some tasks test your ability to interpret statistical test output.
\end{itemize}

\section{Why am I being assessed like
this?}\label{why-am-i-being-assessed-like-this}

\begin{itemize}
\item
  The worksheets assess your ability to wrangle and analyse data in a
  reproducible way. These are important skills for psychological
  researchers to develop.
\item
  Your skills in data skills will progress throughout your degree and
  the worksheets ensure that you are maintaining a good rate of progress
  so that you are prepared for your dissertation.
\end{itemize}

\section{How does this relate to previous work I have
completed?}\label{how-does-this-relate-to-previous-work-i-have-completed}

\begin{itemize}
\item
  The feedback from your RM1 data skills worksheets will be useful for
  these submissions.
\item
  In addition, the guidance and solutions in the data skills book,
  posting on Teams and attending office hours will help you complete
  this assessment.
\end{itemize}

\section{Academic intergrity}\label{academic-intergrity}

Please note that when submitting your work for assessment we accept it
on the understanding that it is your own effort and work and unique to
the set assignment.

To support you in understanding what plagiarism is and in avoiding it,
please read the following resources that the University provides:

\begin{itemize}
\tightlist
\item
  \href{https://www.glasgowunisrc.org/advice/}{SRC Advice and Support}
\item
  \href{https://www.gla.ac.uk/myglasgow/apg/policies/uniregs/regulations2025-26/feesandgeneral/studentsupportandconductmatters/reg33/}{Code
  of Student Conduct} and
  \href{https://www.gla.ac.uk/myglasgow/apg/policies/uniregs/regulations2025-26/feesandgeneral/studentsupportandconductmatters/reg32/}{Plagiarism
  \& Academic Integrity Code}
\item
  \href{https://moodle.gla.ac.uk/course/view.php?id=13038}{Avoiding
  plagiarism and engage in good academic practice} (a Moodle course you
  can self-enrol in)
\item
  \href{https://www.gla.ac.uk/myglasgow/sld/}{Student support for AI,
  plagiarism and digital skills}
\end{itemize}

\textbf{In summary:}

All work submitted by students for assessment is accepted on the
understanding that it is the student's own effort. This means students'
work should not contain:

\begin{itemize}
\tightlist
\item
  plagiarised content; or
\item
  content that has been produced by another person, website, software or
  Artificial intelligence (AI) tool (except where AI use is explicitly
  permitted); or
\item
  content that has been prepared jointly with any other person (except
  where this is explicitly permitted); or
\item
  content that has already been submitted for assessment by the student
  at this or any other institution, known as self-plagiarism.
\end{itemize}

\textbf{Statement on groupwork:} We encourage students to form a study
group and peer feedback groups. However, this assignment is not a group
work assignment, so your work must be your own individual contribution.
If you make a study group or a peer review group, avoid sharing final
drafts or near final drafts of your work.

\textbf{University statement on AI:} The University of Glasgow
recognises the value of generative artificial intelligence (AI) tools in
academic and professional workplaces.The university has a responsibility
to ensure that students acquire the necessary knowledge, skills, and
other competencies associated within their discipline. The Student
Learning Development service provides general guidance and support for
students on the use of generative AI. Each item of assessment in your
courses will have specific guidance about the use of AI. Where
generative AI restrictions are in place, they have been carefully
designed to maximise your learning opportunity whilst discouraging
reliance on generative AI in a way that undermines your learning or
development of good professional practice and graduate attributes.

\textbf{Statement on use of generative AI:} The current assessment is
summative, meaning that it contributes to your course grade. \textbf{The
purpose of this assessment is to develop and practice coding skills and
data analysis. You can use AI as a learning assistant to explain code,
to understand error messages and help documentation, and as a rubber
duck to help articulate problems}. If you want to use translation
software (i.e., code in another language first) be cautious as the
syntax produced by generative AI is not always in keeping with the
current coding language and can result in incorrect answers.

\textbf{Avoid using AI to write your code as you risk not fully
understanding the analysis and so you fail to build a foundation for
future coding/analysis tasks. Avoid using AI to write comments for your
code, as commenting helps you to articulate what you did and also helps
you to understand the analysis and results more deeply.}

There is no expectation that you will use generative AI, and we have no
evidence that its use will confer an advantage for this assessment. If
you do use generative AI, you MUST acknowledge use in-text via citations
and referencing and in an appendix with a declaration of AI use as
appropriate. If you choose to use it, we recommend that you use the
Microsoft Edge Browser with Copilot and sign in with your university
account using the multi-factor authentication to ensure that your work
is private and secure. \textbf{Please keep a log of your use of AI as we
may ask to see this.}

For this assignment, we will consider it a misuse of generative AI if
you do not acknowledge using it. Declare all uses of AI, including
initial exploration of the subject, literature searching, writing and
editing, corrections for grammar and spelling, as well as any other
tasks from the course. Be aware that AI may not represent the best
response for this task, and you need to take responsibility for
everything that is submitted.

\section{Feedback}\label{data_skills_feedback}

\subsection{How is this assessment
graded?}\label{how-is-this-assessment-graded}

\begin{itemize}
\item
  Each worksheet will have 19 questions. Each question will be worth 1
  or 2 points, with all questions adding up to 22 points for each
  worksheet.
\item
  The computer-assisted marking code that we use checks that the objects
  in each code chunk are identical to the solution. It is therefore
  crucial that you follow the instructions carefully and do not change
  the names of the objects or the code chunks. These instructions are
  clearly stated in the submission files and you will lose marks if you
  provide the correct code but change the names.
\end{itemize}

\subsection{How will the feedback from assessment help me in the
future?}\label{how-will-the-feedback-from-assessment-help-me-in-the-future}

\begin{itemize}
\item
  This assessment will help you in any future work that requires working
  with quantitative data such as your dissertation project or future
  postgraduate courses.
\item
  Additionally, R can be used for tasks such as conducting text analyses
  and building websites and therefore is an extremely useful
  transferable skill.
\end{itemize}

\subsection{What feedback will I receive for this
assessment?}\label{what-feedback-will-i-receive-for-this-assessment}

\begin{itemize}
\item
  You will get a feedback sheet for each data skills worksheet that
  tells you your score for each question and the correct solution.
\item
  There will be additional generic feedback provided for questions with
  common mistakes and to explain the answers if necessary beyond the
  solution.
\item
  Finally, there will be individual feedback for any incorrect answers
  as necessary.
\end{itemize}

\subsection{Who assessed my work?}\label{who-assessed-my-work}

\begin{itemize}
\tightlist
\item
  The worksheets will be graded using computer-assisted marking. In the
  first instance, the marking will be done automatically using R and
  then any incorrect answers will be checked manually.
\end{itemize}

\subsection{Can I get more feedback?}\label{can-i-get-more-feedback}

\begin{itemize}
\item
  Yes! We encourage you discuss your assessment (regardless of the grade
  you received) with Wilhelmiina Toivo, who is in charge of marking this
  assessment.
\item
  You can do so by attending her office hours, or contacting her
  directly to arrange an appointment if you cannot make the office
  hours.
\end{itemize}

\subsection{What if I don't agree with my feedback or
grade?}\label{what-if-i-dont-agree-with-my-feedback-or-grade}

\begin{itemize}
\item
  Your first point of contact should be to arrange an additional
  feedback meeting with the course lead.
\item
  Following this, if you still have concerns you should consult
  \href{https://www.glasgowunisrc.org/advice/academic/appeals/}{the
  guidance from the SRC} which provides a clear explanation of the
  University appeals procedures. There are only three grounds for
  appeal:

  \begin{itemize}
  \tightlist
  \item
    Unfair or defective procedure
  \item
    Failure to take into account medical or other adverse personal
    circumstances
  \item
    There are relevant medical or other adverse personal circumstances
    which for good reason have not previously been presented.
  \end{itemize}
\end{itemize}

You cannot appeal against academic judgement.

\chapter{Multiple Choice Questions (Individual)}\label{ch2_mcqs}

In this Chapter, we have included information about the Multiple Choice
Quiz in RM2. You can use the menu at the right hand side of this page to
jump to the different sections.

\section{General Information}\label{general-information-1}

\begin{itemize}
\item
  The deadline for this assessment can be found under `Deadlines' within
  the `Course Information' section on the RM2 Moodle.
\item
  This assessment is worth 10\% of your final course grade.
\end{itemize}

\textbf{Please note:} The MCQ will open up one week before the deadline.
Please ensure you have reviewed the Week 1-14 lecture content before
taking the MCQ.

\section{Type of
Assessment/Structure}\label{type-of-assessmentstructure-1}

\begin{itemize}
\item
  All the information needed to answer the Multiple choice questions
  will have been covered in the research methods lectures between Weeks
  1 and 14.

  \begin{itemize}
  \tightlist
  \item
    The MCQ will assess different types of content:

    \begin{itemize}
    \tightlist
    \item
      Knowledge of quantitative content covered in the lectures
    \item
      Understanding of different quantitative research designs\\
    \item
      Knowledge of qualitative content covered in the lectures
    \end{itemize}
  \end{itemize}
\item
  The MCQ will be provided through a submission link on Moodle. The
  submission link will be made available one week before the deadline.
\item
  There are 22 questions and each MCQ will have four options - the
  answer will be one of these.
\item
  Each question will be presented on a single page. There are many
  reasons for this, including that having multiple questions on a page
  increases moodle loading times, slowing it down.
\item
  You will have 40 minutes to complete the MCQ. If you are registered
  with disability services for extra time it should show longer than
  this. If the quiz does not display additional completion time when you
  click on the link, please contact Admin or Ashley/Wil to let them
  know.
\end{itemize}

\section{Assessment Support}\label{assessment-support-1}

\begin{itemize}
\item
  All the information needed to answer the research methods questions
  will have been covered in the research methods lectures.
\item
  The course leads provide practice questions for each of the lectures
  as well as a practice MCQ. You can find these in the `Formative
  Activities' section on Moodle.
\item
  Further information about feedback can be found in
  \hyperref[MCQ_feedback]{the Feedback section}
\end{itemize}

\section{How to do well in this
assessment}\label{how-to-do-well-in-this-assessment-1}

\begin{itemize}
\item
  Read each question carefully.
\item
  Ensure you are up to date on all lectures and have completed all the
  activities for each week.
\item
  Ensure that you leave yourself enough time to complete the MCQ in a
  single sitting.
\item
  Complete the essential reading for each lecture.
\item
  Engage in opportunities to practice MCQ questions by completing the
  formative mini quizzes.
\end{itemize}

\section{Common Mistakes}\label{common-mistakes-1}

\begin{itemize}
\item
  Leaving the lecture content and reading to build up (i.e.~not keeping
  up with it as you go along).
\item
  Not reading the question carefully.
\item
  Leaving some questions unanswered.
\end{itemize}

\section{How is the assessment related to the lectures for this
course?}\label{how-is-the-assessment-related-to-the-lectures-for-this-course-1}

\begin{itemize}
\item
  The MCQs test your understanding of statistical concepts and how to
  analyse and interpret data according to the methods presented in the
  lectures.
\item
  The MCQs test your understanding of research designs in quantitative
  research.
\item
  The MCQs also test your understanding of the qualitative content
  covered in the lectures.
\end{itemize}

\section{Why am I being assessed like
this?}\label{why-am-i-being-assessed-like-this-1}

\begin{itemize}
\item
  The MCQs assess your understanding of the course content, helping you
  to actively engage with the materials. These are important skills for
  psychological researchers to develop.
\item
  Engaging with the material on the course will help build on the strong
  foundation of Research Methods developed thus far, helping to set you
  up for your dissertation.
\end{itemize}

\section{How does this relate to previous work I have
completed?}\label{how-does-this-relate-to-previous-work-i-have-completed-1}

\begin{itemize}
\tightlist
\item
  Feedback from formative mini quizzes based on the lectures and
  essential reading will help you with this assignment
\end{itemize}

\section{Academic integrity}\label{academic-integrity}

Please note that when submitting your work for assessment we accept it
on the understanding that it is your own effort and work and unique to
the set assignment.

To support you in understanding what plagiarism is and in avoiding it,
please read the following resources that the University provides:

\begin{itemize}
\tightlist
\item
  \href{https://www.glasgowunisrc.org/advice/}{SRC Advice and Support}
\item
  \href{https://www.gla.ac.uk/myglasgow/apg/policies/uniregs/regulations2025-26/feesandgeneral/studentsupportandconductmatters/reg33/}{Code
  of Student Conduct} and
  \href{https://www.gla.ac.uk/myglasgow/apg/policies/uniregs/regulations2025-26/feesandgeneral/studentsupportandconductmatters/reg32/}{Plagiarism
  \& Academic Integrity Code}
\item
  \href{https://moodle.gla.ac.uk/course/view.php?id=13038}{Avoiding
  plagiarism and engage in good academic practice} (a Moodle course you
  can self-enrol in)
\item
  \href{https://www.gla.ac.uk/myglasgow/sld/}{Student support for AI,
  plagiarism and digital skills}
\end{itemize}

\textbf{In summary:}

All work submitted by students for assessment is accepted on the
understanding that it is the student's own effort. This means students'
work should not contain:

\begin{itemize}
\tightlist
\item
  plagiarised content; or
\item
  content that has been produced by another person, website, software or
  Artificial intelligence (AI) tool (except where AI use is explicitly
  permitted); or
\item
  content that has been prepared jointly with any other person (except
  where this is explicitly permitted); or
\item
  content that has already been submitted for assessment by the student
  at this or any other institution, known as self-plagiarism.
\end{itemize}

\textbf{Statement on groupwork:} We encourage students to form a study
group and peer feedback groups. However, this assignment is not a group
work assignment, so your work must be your own individual contribution.
If you make a study group or a peer review group, avoid sharing final
drafts or near final drafts of your work.

\textbf{University statement on AI:} The University of Glasgow
recognises the value of generative artificial intelligence (AI) tools in
academic and professional workplaces.The university has a responsibility
to ensure that students acquire the necessary knowledge, skills, and
other competencies associated within their discipline. The Student
Learning Development service provides general guidance and support for
students on the use of generative AI. Each item of assessment in your
courses will have specific guidance about the use of AI. Where
generative AI restrictions are in place, they have been carefully
designed to maximise your learning opportunity whilst discouraging
reliance on generative AI in a way that undermines your learning or
development of good professional practice and graduate attributes.

\textbf{Statement on use of generative AI:} The current assessment is
summative, meaning that it contributes to your course grade. \textbf{The
purpose of this assessment is to provide an opportunity to test your
understanding of the lecture content in the course. You can use AI as a
learning assistant to help understand lecture content and to generate
practice MCQs.} If you want to use translation software (i.e., you write
the assignment in another language first) be cautious as the vocabulary
and syntax produced by generative AI is not generally in keeping with
the current language use and vocabulary in our field and can result in
subtle misunderstanding in communication.

\textbf{Avoid using AI to answer the questions because your
understanding is a skill you need to develop.}

There is no expectation that you will use generative AI, and we have no
evidence that its use will confer an advantage for this assessment. If
you do use generative AI, you MUST acknowledge use in-text via citations
and referencing and in an appendix with a declaration of AI use as
appropriate. If you choose to use it, we recommend that you use the
Microsoft Edge Browser with Copilot and sign in with your university
account using the multi-factor authentication to ensure that your work
is private and secure. \textbf{Please keep a log of your use of AI as we
may ask to see this.}

For this assignment, we will consider it a misuse of generative AI if
you do not acknowledge using it. Declare all uses of AI, including
initial exploration of the subject, literature searching, writing and
editing, corrections for grammar and spelling, as well as any other
tasks from the course. Be aware that AI may not represent the best
response for this task, and you need to take responsibility for
everything that is submitted.

\section{Feedback}\label{MCQ_feedback}

\subsection{How is this assessment
graded?}\label{how-is-this-assessment-graded-1}

\begin{itemize}
\item
  The MCQ will consist of 22 questions, each with four response options.
  Each question is worth one point.
\item
  The marking will be automatically computed on Moodle, and grades
  downloaded by staff.
\end{itemize}

\subsection{How will feedback from this assessment help me in the
future?}\label{how-will-feedback-from-this-assessment-help-me-in-the-future}

\begin{itemize}
\tightlist
\item
  This assessment will help you consolidate your knowledge of the
  statistical and qualitative content covered within the lecture
  materials. This will be useful in any future work that requires
  working with quantitative or qualitative data such as your
  dissertation project or future postgraduate courses.
\end{itemize}

\subsection{What type of feedback will I receive for this
assessment?}\label{what-type-of-feedback-will-i-receive-for-this-assessment}

\begin{itemize}
\tightlist
\item
  After the marking has been completed, you will find out a) whether you
  answered each question correctly or not and b) what the correct answer
  was for each question
\end{itemize}

\subsection{Who assessed my work?}\label{who-assessed-my-work-1}

\begin{itemize}
\tightlist
\item
  The worksheets will be graded using computer-assisted marking. The
  marking will be done automatically using Moodle.
\end{itemize}

\chapter{Qualitative Portfolio (Group)}\label{ch3_AIS_portfolio}

In this Chapter, we have included information about the Qualitative
Portfolio in RM2. You can use the menu at the right hand side of this
page to jump to the different sections.

\section{General Information}\label{general-information-2}

\begin{itemize}
\item
  The deadline for this assessment can be found under `Deadlines' within
  the `Course Information' section on the RM2 Moodle.
\item
  This is a group assessment, with each group typically consisting of
  around 4-6 students.
\item
  The qualitative portfolio consists of three activities that should be
  completed and submitted in your groups.
\item
  This assessment is worth 30\% of your final course grade and will be a
  group mark.
\item
  One member of your group should submit your portfolio to Moodle by the
  deadline.
\item
  The submission link will open no later than 5 working days before the
  deadline and will be found in the Assignment Submission section of
  Moodle.
\end{itemize}

\section{Word Count and Formatting}\label{word-count-and-formatting}

\begin{itemize}
\item
  The maximum word count for the qualitative portfolio is 1000 words.
  The word limit includes all written aspects, including any additional
  headings and in-text citations for Activities 1B, 2 and 3. However, it
  does not include the qualitative questions for Activity 1A and your
  reference list. There is no 10\% rule, 1000 words is a strict upper
  limit.
\item
  Please use the template provided for the portfolio.
\item
  Your work should be presented in Times New Roman, 12-point font,
  double-spaced with 1-inch (2.54cm) margins.
\item
  Your submission should be written and formatted according to APA 7th
  guidelines.
\end{itemize}

\section{Type of
Assessment/Structure}\label{type-of-assessmentstructure-2}

\begin{itemize}
\item
  This assessment will be a group and collaborative submission. Please
  note that there is a single mark for each submission, which will
  consist of three separate sections. All group members will receive the
  same mark for the submission. Sections will not, under any
  circumstances, be considered separately. Therefore it is vital that
  \textbf{all members of the group read and agree on the whole
  submission beforehand.}
\item
  The portfolio will consist of a) developing questions for qualitative
  research, and reflecting on these, b) an ethics exercise, and c)
  identifying issues in a recording of a focus group.
\item
  For this assessment, you will be required to:

  \begin{itemize}
  \tightlist
  \item
    follow the guidance for each of the three sections, ensuring that
    you complete the tasks as directed
  \item
    reflect on the issues presented in the materials, considering how to
    mitigate them
  \item
    link to the literature base and ethical guidelines, helping to
    support your points.
  \end{itemize}
\end{itemize}

\section{Intended Learning Outcomes (ILOs)}\label{portfolio_ILOs}

Below, we have the ILOs for this assessment. Please note that some of
the ILOs are over-arching (i.e.~apply to all three sections) whereas
others are specifically for a single section. To increase clarity, we
have stated which section/s apply for each ILO.

\begin{enumerate}
\def\labelenumi{\arabic{enumi}.}
\item
  Quality of the Knowledge and Research

  \begin{itemize}
  \tightlist
  \item
    Develop six questions that would be likely to generate open
    discussion and rich responses from participants (\#1: Qual
    questions)
  \item
    Correctly identify ONE ethical issue, explain its importance, and
    provide guidance on how to mitigate it (\#2: Ethics)
  \item
    Correctly identify TWO or THREE issues present in the focus group,
    explain why they are problematic, and discuss possible improvements
    (\#3: Focus group)
  \end{itemize}
\item
  Quality of the Evaluation

  \begin{itemize}
  \tightlist
  \item
    As appropriate for each section, use academic evidence and/or
    reference to ethical guidelines to support your evaluation (\#1:
    Qual questions, \#2 Ethics, \#3: Focus group)
  \item
    As appropriate for each section, demonstrate understanding of the
    practical issues of running qualitative research through evaluation
    and reflection (\#1: Qual questions, \#2 Ethics, \#3: Focus group)
  \item
    As appropriate for each section, synthesise academic evidence from
    multiple perspectives to demonstrate critical thinking and to
    support your position (\#1 Qual questions, \#2 Ethics, \#3 Focus
    group)
  \end{itemize}
\item
  Quality of the Academic Communication

  \begin{itemize}
  \tightlist
  \item
    Write clearly and succinctly with appropriate use of paragraphs,
    spelling and grammar, following APA 7th guidelines for all citations
    and references (\#1: Qual questions, \#2 Ethics, \#3: Focus group)
  \end{itemize}
\end{enumerate}

\emph{Please note: In the first bullet point under ILO 2 (Quality of the
Evaluation), we are not asking for you to relate to ethical guidelines
for each section. We are asking you to link to academic evidence/ethical
guidelines (as appropriate) for each of the three sections.}

\section{Assessment Support}\label{assessment-support-2}

\begin{itemize}
\item
  Guidance on how to complete the group portfolio will be provided as
  part of the course. These will be released on Moodle, but you will
  also be able to find all key resources in
  \hyperref[ch10_group_project]{the Group Project Chapter}.
\item
  Further information about feedback can be found in
  \hyperref[portfolio_feedback]{the Feedback section}
\item
  Additional writing and study advice, including 1-to-1 guidance from
  \href{https://www.gla.ac.uk/myglasgow/sld/}{Student Learning
  Development}
\item
  FAQs from previous students can be found in \hyperref[ch13_faqs]{the
  FAQs Chapter}
\end{itemize}

\section{How to do well in this
assessment}\label{how-to-do-well-in-this-assessment-2}

\begin{itemize}
\item
  Meet each of the intended learning outcomes - use these as a checklist
  for your work.
\item
  Consult the guidelines provided closely, as these will provide key
  information about what is required in this assessment.
\item
  Identify the key points in the literature that are relevant to each
  task, and relate them to your study.
\item
  \textbf{Task 1: Qual questions} Develop qualitative questions that are
  open, clear, concise and appropriate. Take into account whether you
  are designing them for an interview, focus group or survey and use
  relevant evidence in your reflection.
\item
  \textbf{Task 2 Ethics} Remember to link your identified ethics issue
  to the BPS code of Ethics.
\item
  \textbf{Task 3: Focus Group} Make sure to use evidence to support your
  discussion on your chosen issues. Follow the guidelines on how many
  issues to discuss.
\item
  Allow time to proof-read your work before submission.
\item
  Work as a group and make sure everybody reads and agrees on the final
  submission.
\end{itemize}

\section{Common Mistakes}\label{common-mistakes-2}

\begin{itemize}
\item
  Developing questions that are closed, inappropriate, rambling or ask
  multiple different questions at once.
\item
  Not linking to the evidence base in your written response.
\item
  Not following the guidance provided for the assessment.
\item
  Not following the guidance on how many ethical/focus group issues to
  discuss.
\item
  Providing reflections that are shallow.
\item
  Writing that is unclear and/or vague.
\item
  Failure to adhere to the overall word limit.
\item
  Not following the suggested word counts for each activity.
\item
  Splitting sections among members of the group and not then considering
  the submission as a coherent whole. It is important to remember that
  all members should approve the full submission beforehand.
\end{itemize}

\section{How is the assessment related to the lectures for this
course?}\label{how-is-the-assessment-related-to-the-lectures-for-this-course-2}

\begin{itemize}
\item
  The qualitative portfolio allows you the opportunity to apply many of
  the concepts that you will learn about in more detail in the lectures.
\item
  It will help give you insight to some of the practical issues related
  with conducting qualitative research, and considering how you might
  address these.
\end{itemize}

\section{Why am I being assessed like
this?}\label{why-am-i-being-assessed-like-this-2}

\begin{itemize}
\item
  Completing the qualitative portfolio is very important for a) better
  understanding how to conduct qualitative research practically, b)
  giving an understanding of the processes that underlie the data you
  will use for your qualitative report. These are key steps to your
  project, and help form a strong foundation.
\item
  The qualitative portfolio is a group submission to reflect the fact
  that in most research, these decisions will be made as a team and it
  allows you to pool your collective knowledge to design the best study
  possible. Group work is an essential part of requirements for the
  British Psychological Society, and is a key part of this course.
\end{itemize}

\section{How does this relate to previous work I have
completed?}\label{how-does-this-relate-to-previous-work-i-have-completed-2}

\begin{itemize}
\item
  You can gain informal feedback by posting on Teams, attending student
  office hours and attending the course leads Q\&A session.
\item
  Feedback on any previous written assignment will help with academic
  communication and using evidence to support your arguments.
\end{itemize}

Feedback on your Stage 1 and Stage 2 reports for RM1 will help you with
how you link to the evidence base in your evaluation.

\section{Academic integrity}\label{academic-integrity-1}

Please note that when submitting your work for assessment we accept it
on the understanding that it is your own effort and work and unique to
the set assignment.

To support you in understanding what plagiarism is and in avoiding it,
please read the following resources that the University provides:

\begin{itemize}
\tightlist
\item
  \href{https://www.glasgowunisrc.org/advice/}{SRC Advice and Support}
\item
  \href{https://www.gla.ac.uk/myglasgow/apg/policies/uniregs/regulations2025-26/feesandgeneral/studentsupportandconductmatters/reg33/}{Code
  of Student Conduct} and
  \href{https://www.gla.ac.uk/myglasgow/apg/policies/uniregs/regulations2025-26/feesandgeneral/studentsupportandconductmatters/reg32/}{Plagiarism
  \& Academic Integrity Code}
\item
  \href{https://moodle.gla.ac.uk/course/view.php?id=13038}{Avoiding
  plagiarism and engage in good academic practice} (a Moodle course you
  can self-enrol in)
\item
  \href{https://www.gla.ac.uk/myglasgow/sld/}{Student support for AI,
  plagiarism and digital skills}
\end{itemize}

\textbf{In summary:}

All work submitted by students for assessment is accepted on the
understanding that it is the student's own effort. This means students'
work should not contain:

\begin{itemize}
\tightlist
\item
  plagiarised content; or
\item
  content that has been produced by another person, website, software or
  Artificial intelligence (AI) tool (except where AI use is explicitly
  permitted); or
\item
  content that has been prepared jointly with any other person (except
  where this is explicitly permitted); or
\item
  content that has already been submitted for assessment by the student
  at this or any other institution, known as self-plagiarism.
\end{itemize}

\textbf{Statement on groupwork:} This is a group assignment and one
person from the group should be nominated to submit the final version of
the assignment on Moodle. Your group's work should not be exactly the
same as that of another group in the class, however, as you worked
closely from common templates, we know that there may be some
unavoidable similarities.

\textbf{University statement on AI:} The University of Glasgow
recognises the value of generative artificial intelligence (AI) tools in
academic and professional workplaces.The university has a responsibility
to ensure that students acquire the necessary knowledge, skills, and
other competencies associated within their discipline. The Student
Learning Development service provides general guidance and support for
students on the use of generative AI. Each item of assessment in your
courses will have specific guidance about the use of AI. Where
generative AI restrictions are in place, they have been carefully
designed to maximise your learning opportunity whilst discouraging
reliance on generative AI in a way that undermines your learning or
development of good professional practice and graduate attributes.

\textbf{Statement on use of generative AI:} The current assessment is
summative, meaning that it contributes to your course grade. \textbf{The
purpose of this assessment is to provide an opportunity to develop your
writing and communication. You can use AI as a learning assistant to
help understand lecture content, to give feedback on your work, and to
help in finding spelling and grammatical errors.} If you want to use
translation software (i.e., you write the assignment in another language
first) be cautious as the vocabulary and syntax produced by generative
AI is not generally in keeping with the current language use and
vocabulary in our field and can result in subtle misunderstanding in
communication.

\textbf{Avoid using AI to draft your assignment or structure your work,
as these are specialist skills that you need to practice. Avoid using AI
to make your citations or reference section, as there is a risk that it
will fabricate information.}

There is no expectation that you will use generative AI, and we have no
evidence that its use will confer an advantage for this assessment. If
you do use generative AI, you MUST acknowledge use in-text via citations
and referencing and in an appendix with a declaration of AI use as
appropriate. If you choose to use it, we recommend that you use the
Microsoft Edge Browser with Copilot and sign in with your university
account using the multi-factor authentication to ensure that your work
is private and secure. \textbf{Please keep a log of your use of AI as we
may ask to see this.}

For this assignment, we will consider it a misuse of generative AI if
you do not acknowledge using it. Declare all uses of AI, including
initial exploration of the subject, literature searching, writing and
editing, corrections for grammar and spelling, as well as any other
tasks from the course. Be aware that AI may not represent the best
response for this task, and you need to take responsibility for
everything that is submitted.

\section{Feedback}\label{portfolio_feedback}

\subsection{What type of feedback will I receive for this
assessment?}\label{what-type-of-feedback-will-i-receive-for-this-assessment-1}

\begin{itemize}
\item
  You will receive written on-script comments, as well as three
  actionable feed-forward suggestions for work in the future (e.g.~your
  qual report, your dissertation).
\item
  You will also receive a rating on each of the expanded ILOs.
\end{itemize}

\subsection{Can I get more feedback?}\label{can-i-get-more-feedback-1}

\begin{itemize}
\tightlist
\item
  If you would like to discuss your mark and feedback you should contact
  your marker, however, \textbf{we ask that you wait 24 hours after the
  release of the grades} before you do so to give you time to fully
  reflect on the feedback given.
\end{itemize}

You can arrange a meeting with your marker to clarify your grades and
feedback but it can also be the case that you understand everything that
was written and you just want a bit more feedback, or you'd like to chat
about the essay generally. Please do make use of the opportunity to talk
with your marker because it really does make a difference.

To help the discussion, when you e-mail your marker, you must complete
\href{https://gla-my.sharepoint.com/:w:/g/personal/ashley_robertson_glasgow_ac_uk/EanCqf_RV7pMumQt1Bm0N4YB3RlzoIjP9kmbe0qbYj9rYQ?e=MDWxPP}{the
feedback reflection form} which will ask you to consider the below:

\begin{itemize}
\tightlist
\item
  Confirm that you read the following feedback:

  \begin{itemize}
  \tightlist
  \item
    All on-script comments
  \item
    Any feed-forward comments
  \item
    The ratings on the ILOs
  \end{itemize}
\item
  Whether there was any feedback you did not understand or agree with.
\item
  Whether you think your feedback aligned with your grade. If not, we
  ask for an explanation.
\item
  Whether there is any aspect of your work you would like more feedback
  on.
\end{itemize}

\subsection{How is this assessment
graded?}\label{how-is-this-assessment-graded-2}

\begin{itemize}
\item
  We use the
  \href{https://www.gla.ac.uk/media/Media_1200302_smxx.pdf}{Schedule A
  marking criteria}, marking each assessment on a 22-point scale.
\item
  We use the \hyperref[portfolio_ILOs]{Intended Learning Outcomes
  presented above} to mark this assessment.
\end{itemize}

\subsection{How will feedback from this assessment help me in the
future?}\label{how-will-feedback-from-this-assessment-help-me-in-the-future-1}

\begin{itemize}
\item
  Primarily, the feedback (and feedforward) obtained will support the
  write-up of your RM2 report. It will also be relevant for your
  dissertation in Year 3.
\item
  Additionally, the feedback will help in any future research work you
  conduct that requires qualitative data collection, qualitative
  analysis, and evidence-based justification.
\end{itemize}

\subsection{Who assessed my work?}\label{who-assessed-my-work-2}

\begin{itemize}
\item
  The first marker for this assessment will be a member of the RM2 staff
  team.
\item
  Following University's policy, your assessment may be second marked by
  another member of the RM2 team. A range of work from across all
  markers will be second marked, to ensure that we are applying
  appropriate standards in assessment and that they are being applied
  consistently across the cohort of students being assessed.
\end{itemize}

\subsection{What if I don't agree with my feedback or
grade?}\label{what-if-i-dont-agree-with-my-feedback-or-grade-1}

\begin{itemize}
\item
  Your first point of contact should be to arrange an additional
  feedback meeting with the marker of the report (the name will have
  been provided in your feedback). It is most likely to be either Ashley
  or Wil who will have marked this assessment. This meeting should be to
  gain additional feedback from the marker, rather than to contest the
  grade.
\item
  Following this, if you still have concerns, you should consult
  \href{https://www.glasgowunisrc.org/advice/academic/appeals/}{the
  guidance from the SRC} which provides a clear explanation of the
  University appeals procedures. There are only three grounds for
  appeal:

  \begin{itemize}
  \tightlist
  \item
    Unfair or Defective Procedure
  \item
    Failure to take into account medical or other adverse personal
    circumstances
  \item
    There are relevant medical or other adverse personal circumstances
    which for good reason have not previously been presented.
  \end{itemize}
\end{itemize}

It is not possible to appeal against academic judgement.

\chapter{Qualitative Report (Individual)}\label{ch4_AIS_report}

In this Chapter, we have included information about the Qualitative
Report in RM2. You can use the menu at the right hand side of this page
to jump to the different sections.

\section{General Information}\label{general-information-3}

\begin{itemize}
\item
  The deadline for this assessment can be found under `Deadlines' within
  the `Course Information' section on the RM2 Moodle.
\item
  There are different topics for this assessment, and you should choose
  ONE of these. These will be released by course leads at the beginning
  of Semester 2.
\item
  This is an individual submission. We ask for the following documents
  to be submitted as \textbf{separate} documents:

  \begin{itemize}
  \tightlist
  \item
    the qualitative report (including references and appendices)
  \item
    the step-by step analysis template
  \end{itemize}
\item
  This assessment is worth 50\% of your final course grade.
\item
  The submission link will open no later than 5 working days before the
  deadline and will be found in the Assignment Submission section of the
  RM2 Moodle.
\end{itemize}

\section{Word Count and Formatting}\label{word-count-and-formatting-1}

\begin{itemize}
\item
  The maximum word count for the report is 3000 words. This includes all
  aspects of the report with the exception of the reference list and
  appendices. There is no 10\% rule, 3000 words is a strict upper limit.
\item
  Your work should be presented in Times New Roman, 12-point font,
  double-spaced with 1-inch (2.54cm) margins.
\item
  All aspects of your report should be written and formatted according
  to APA guidelines.
\end{itemize}

\section{Intended Learning Outcomes}\label{report_ILOs}

\begin{enumerate}
\def\labelenumi{\arabic{enumi}.}
\item
  Quality of the Knowledge and Research

  \begin{itemize}
  \tightlist
  \item
    Provide a comprehensive review of literature, which synthesises the
    main ideas and contextualises key issues.
  \item
    Clearly justify the aims of the project
  \item
    Clearly and correctly report the methods of the study
  \end{itemize}
\item
  Quality of the Evaluation

  \begin{itemize}
  \tightlist
  \item
    Evaluate your role and positionality as a researcher in shaping the
    research process.
  \item
    Clearly present themes and interpret them within a narrative. Ensure
    these are supported by verbatim evidence from the transcript.
  \item
    Evaluate your study and how your results fit into the wider
    literature and social/psychological context.
  \item
    Identify relevant methodological limitations, implications, and
    future directions in the context of previous research.
  \end{itemize}
\item
  Quality of the Academic Communication

  \begin{itemize}
  \tightlist
  \item
    Write clearly and succinctly with appropriate use of paragraphs,
    spelling and grammar. Ensure that all parts of the report have a
    logical structure. Ensure that the report follows APA style
    consistently.
  \item
    Clearly support ideas and statements with appropriate citations.
  \item
    Clearly and transparently record the analytical process. Upload the
    separate submission evidencing the stages of your Thematic Analysis
    process.
  \end{itemize}
\end{enumerate}

\section{Assessment Support}\label{assessment-support-3}

\begin{itemize}
\item
  Guidance on how to complete the qualitative report will be provided as
  part of the course. These will be released on Moodle, but you will
  also be able to find all key resources in
  \hyperref[ch12_qualitative_report]{the Qualitative Report Chapter}.
\item
  Further information about feedback can be found in
  \hyperref[report_feedback]{the Feedback section}
\item
  Additional writing and study advice, including 1-to-1 guidance from
  \href{https://www.gla.ac.uk/myglasgow/sld/}{Student Learning
  Development}
\item
  FAQs from previous students can be found in \hyperref[ch13_faqs]{the
  FAQs Chapter}
\end{itemize}

\section{How to do well in this
assessment}\label{how-to-do-well-in-this-assessment-3}

\begin{itemize}
\item
  Meet each of the intended learning outcomes - use these as a checklist
  for your work.
\item
  Write clearly, concisely, with an academic tone and logical structure.
\item
  Provide an evidence-based rationale for your research question(s), AND
  an evidence base for the use of qualitative methods.
\item
  Narrow the focus of the literature review to a specific topic (for
  example, not a general review of all the literature on disabled
  people's experiences).
\item
  Be clear about who the participants in your sample is and why this is
  relevant in relation to previous research.
\item
  Engage with qualitative research! Read some of the articles in the
  reading list, or qualitative research in your topic of interest.
\item
  Be ethical - do not share sensitive information or anything that might
  identify your participants. The quotes and transcript that you share
  must be appropriate for the audience (i.e.~markers, examiners, fellow
  students).
\item
  Keep the method section brief. Saying that, refer to materials
  (e.g.~focus group schedules) that are placed in the appendices.
\item
  In the analysis, focus on two main themes OR one main theme and two
  sub-themes. Describe each theme clearly and explain the themes using
  quotes as evidence.
\item
  In the analysis the narrative goes beyond pure description and
  identifies underlying themes (e.g., contrasting different viewpoints,
  conflicting emotions, identity).
\item
  Evaluation of findings in relation to previous research demonstrates a
  clear link between the introduction and discussion.
\item
  Produce a professional document by proof reading the entire report.
\item
  Provide a concise conclusion and an overall summary that covers each
  section of the report in the form of an abstract.
\item
  Adhere to APA conventions for referencing and formatting.
\end{itemize}

\section{Common Mistakes}\label{common-mistakes-3}

\begin{itemize}
\item
  Vague or general statements as to why the project is utilising
  qualitative methods.
\item
  Broad review of the topic area which lacks focus.
\item
  Basing the rationale for the report upon personal opinion or ``common
  knowledge'' rather than peer-reviewed evidence.
\item
  Missing detail and/or unnecessary detail in the method section.
\item
  Fragmented analysis with no coherent narrative.
\item
  Broad or overly descriptive analysis.
\item
  A lack of evaluation of the literature in relation to the project
  findings, and/or a discussion of the limitations that is not supported
  by evidence.
\item
  Evaluation of findings using quantitative, rather than qualitative,
  terminology (e.g., hypotheses validity, generalisation).
\item
  Writing that contains grammatical errors, a lack of clarity, or an
  informal tone.
\item
  Failure to adhere to the word limit.
\end{itemize}

\section{How is the assessment related to the lectures for the
course?}\label{how-is-the-assessment-related-to-the-lectures-for-the-course}

\begin{itemize}
\item
  The qualitative project assesses your ability to apply your knowledge
  of qualitative research methods from the lectures and using this in a
  practical context. This will allow you to develop skills in
  qualitative research design and analytical techniques.
\item
  Depending on the exact topic you choose, the project may also be
  related to content from our lectures across the rest of the Conversion
  course.
\end{itemize}

\section{Why am I being assessed like
this?}\label{why-am-i-being-assessed-like-this-3}

\begin{itemize}
\item
  The assessment reflects the structure of a qualitative research
  article and is an example of the format used in this area to report
  qualitative research and is relevant dissertations using qualitative
  methods.
\item
  The qualitative project will help you with future reports you write
  (whether they are qualitative or not), including your dissertation.
\end{itemize}

\section{How does this relate to previous work I have
completed?}\label{how-does-this-relate-to-previous-work-i-have-completed-3}

\begin{itemize}
\item
  Research Methods 2 trains you to design, conduct and write up research
  reports. Although the use of qualitative research method in this
  assessment is new, you have already completed work that requires
  evidence-based writing, critical analysis, and clear and concise
  communication.
\item
  You have already completed a report in RM1, and -- although this was
  quantitative -- it's similar in overall structure, and therefore the
  feedback received from this will help support this assessment.
\item
  Feedback from the group portfolio assessment will also help you with
  the report in terms of evaluation, academic communication and writing
  skills.
\item
  Feedback on any written assignment will help with academic
  communication.
\end{itemize}

\section{Academic integrity}\label{academic-integrity-2}

Please note that when submitting your work for assessment we accept it
on the understanding that it is your own effort and work and unique to
the set assignment.

To support you in understanding what plagiarism is and in avoiding it,
please read the following resources that the University provides:

\begin{itemize}
\tightlist
\item
  \href{https://www.glasgowunisrc.org/advice/}{SRC Advice and Support}
\item
  \href{https://www.gla.ac.uk/myglasgow/apg/policies/uniregs/regulations2025-26/feesandgeneral/studentsupportandconductmatters/reg33/}{Code
  of Student Conduct} and
  \href{https://www.gla.ac.uk/myglasgow/apg/policies/uniregs/regulations2025-26/feesandgeneral/studentsupportandconductmatters/reg32/}{Plagiarism
  \& Academic Integrity Code}
\item
  \href{https://moodle.gla.ac.uk/course/view.php?id=13038}{Avoiding
  plagiarism and engage in good academic practice} (a Moodle course you
  can self-enrol in)
\item
  \href{https://www.gla.ac.uk/myglasgow/sld/}{Student support for AI,
  plagiarism and digital skills}
\end{itemize}

\textbf{In summary:}

All work submitted by students for assessment is accepted on the
understanding that it is the student's own effort. This means students'
work should not contain:

\begin{itemize}
\tightlist
\item
  plagiarised content; or
\item
  content that has been produced by another person, website, software or
  Artificial intelligence (AI) tool (except where AI use is explicitly
  permitted); or
\item
  content that has been prepared jointly with any other person (except
  where this is explicitly permitted); or
\item
  content that has already been submitted for assessment by the student
  at this or any other institution, known as self-plagiarism.
\end{itemize}

\textbf{Statement on groupwork:} This report is not a group work
assignment, so your work must be your own individual contribution.
However, as you worked closely in a small team and from common
templates, we know that there may be some unavoidable similarities
between team members in the methods and results, but it should never be
identical or close to identical.

\textbf{University statement on AI:} The University of Glasgow
recognises the value of generative artificial intelligence (AI) tools in
academic and professional workplaces.The university has a responsibility
to ensure that students acquire the necessary knowledge, skills, and
other competencies associated within their discipline. The Student
Learning Development service provides general guidance and support for
students on the use of generative AI. Each item of assessment in your
courses will have specific guidance about the use of AI. Where
generative AI restrictions are in place, they have been carefully
designed to maximise your learning opportunity whilst discouraging
reliance on generative AI in a way that undermines your learning or
development of good professional practice and graduate attributes.

\textbf{Statement on use of generative AI:} The current assessment is
summative, meaning that it contributes to your course grade. \textbf{The
purpose of this assessment is to provide an opportunity to complete a
qualitative analysis and write it up as a research report. You can use
AI as a learning assistant to help understand lecture content, to give
feedback on your work, and to help in finding spelling and grammatical
errors.} If you want to use translation software (i.e., you write the
assignment in another language first) be cautious as the vocabulary and
syntax produced by generative AI is not generally in keeping with the
current language use and vocabulary in our field and can result in
subtle misunderstanding in communication.

\textbf{Avoid using AI to draft your assignment or structure your work,
as these are specialist skills that you need to practice. Avoid using AI
to make your citations or reference section, as there is a risk that it
will fabricate information. Avoid using AI to do coding or generating
themes in your qualitative analysis, for three reasons: 1) there are
ethical issues as participants in the research datasets have not
approved their data being loaded to GenAI, 2) the output generated tends
to be of poor quality on its own, and 3) you need to practice these
skills so you can develop them.}

There is no expectation that you will use generative AI, and we have no
evidence that its use will confer an advantage for this assessment. If
you do use generative AI, you MUST acknowledge use in-text via citations
and referencing and in an appendix with a declaration of AI use as
appropriate. If you choose to use it, we recommend that you use the
Microsoft Edge Browser with Copilot and sign in with your university
account using the multi-factor authentication to ensure that your work
is private and secure. \textbf{Please keep a log of your use of AI as we
may ask to see this.}

For this assignment, we will consider it a misuse of generative AI if
you do not acknowledge using it. Declare all uses of AI, including
initial exploration of the subject, literature searching, writing and
editing, corrections for grammar and spelling, as well as any other
tasks from the course. Be aware that AI may not represent the best
response for this task, and you need to take responsibility for
everything that is submitted.

\section{Feedback}\label{report_feedback}

\subsection{What type of feedback will I receive for this
assessment?}\label{what-type-of-feedback-will-i-receive-for-this-assessment-2}

\begin{itemize}
\item
  You will receive written on-script comments, as well as three
  actionable feed-forward suggestions for your dissertation.
\item
  You will also receive a rating on each of the expanded ILOs.
\end{itemize}

\subsection{Can I get more feedback?}\label{can-i-get-more-feedback-2}

\begin{itemize}
\tightlist
\item
  If you would like to discuss your mark and feedback you should contact
  your marker, however, \textbf{we ask that you wait 24 hours after the
  release of the grades} before you do so to give you time to fully
  reflect on the feedback given.
\end{itemize}

You can arrange a meeting with your marker to clarify your grades and
feedback but it can also be the case that you understand everything that
was written and you just want a bit more feedback, or you'd like to chat
about the essay generally. Please do make use of the opportunity to talk
with your marker because it really does make a difference.

To help the discussion, when you e-mail your marker, you must complete
\href{https://gla-my.sharepoint.com/:w:/g/personal/ashley_robertson_glasgow_ac_uk/EanCqf_RV7pMumQt1Bm0N4YB3RlzoIjP9kmbe0qbYj9rYQ?e=MDWxPP}{the
feedback reflection form} which will ask you to consider the below:

\begin{itemize}
\tightlist
\item
  Confirm that you read the following feedback:

  \begin{itemize}
  \tightlist
  \item
    All on-script comments
  \item
    Any feed-forward comments
  \item
    The ratings on the ILOs
  \end{itemize}
\item
  Whether there was any feedback you did not understand or agree with.
\item
  Whether you think your feedback aligned with your grade. If not, we
  ask for an explanation.
\item
  Whether there is any aspect of your work you would like more feedback
  on.
\end{itemize}

\subsection{How is this assessment
graded?}\label{how-is-this-assessment-graded-3}

\begin{itemize}
\item
  We use the
  \href{https://www.gla.ac.uk/media/Media_1200302_smxx.pdf}{Schedule A
  marking criteria}, marking each assessment on a 22-point scale.
\item
  We use \hyperref[report_ILOs]{the Intended Learning Outcomes presented
  above} to mark this assessment.
\end{itemize}

\subsection{How will feedback from this assessment help me in the
future?}\label{how-will-feedback-from-this-assessment-help-me-in-the-future-2}

\begin{itemize}
\item
  Primarily, the feedback (and feedforward) obtained will support the
  write-up of your dissertation.
\item
  Additionally, the feedback will help in any future research work you
  conduct that requires qualitative data collection, qualitative
  analysis, and evidence-based justification.
\end{itemize}

\subsection{Who assessed my work?}\label{who-assessed-my-work-3}

\begin{itemize}
\item
  The first marker for this assessment will be a member of the RM2 staff
  team.
\item
  Following University's policy, your assessment may be second marked by
  another member of the RM2 team. A range of work from across all
  markers will be second marked, to ensure that we are applying
  appropriate standards in assessment and that they are being applied
  consistently across the cohort of students being assessed.
\end{itemize}

\subsection{What if I don't agree with my feedback or
grade?}\label{what-if-i-dont-agree-with-my-feedback-or-grade-2}

\begin{itemize}
\item
  Your first point of contact should be to arrange an additional
  feedback meeting with the marker of the report (the name will have
  been provided in your feedback). This meeting should be to gain
  additional feedback from the marker, rather than to contest the grade.
  If this does not bring a resolution, please reach out to the course
  lead/s in the first instance.
\item
  Following this, if you still have concerns, you should consult
  \href{https://www.glasgowunisrc.org/advice/academic/appeals/}{the
  guidance from the SRC} which provides a clear explanation of the
  University appeals procedures. There are only three grounds for
  appeal:

  \begin{itemize}
  \tightlist
  \item
    Unfair or Defective Procedure
  \item
    Failure to take into account medical or other adverse personal
    circumstances
  \item
    There are relevant medical or other adverse personal circumstances
    which for good reason have not previously been presented.
  \end{itemize}
\end{itemize}

It is not possible to appeal against academic judgement.

\part{Qualitative Project Information}

\chapter{Qualitative designs}\label{ch5_qual_designs}

In this chapter we are comparing and contrasting quantitative and
qualitative designs to help you with the gear shift.

These are two very different ways to approach data collection and
consequently, the type of data they yield is very different. It is
important to recognise that not one method is ``superior'' or easier
than the other - the paradigm you choose depends on the type of
questions you want to answer.

Complete the activities below to familiarise yourself with the key
differences between the two methodologies.

\section{Activity 1: Differences between qualitative and quantitative
designs}\label{activity-1-differences-between-qualitative-and-quantitative-designs}

\begin{longtable}[]{@{}
  >{\raggedright\arraybackslash}p{(\linewidth - 2\tabcolsep) * \real{0.5319}}
  >{\raggedright\arraybackslash}p{(\linewidth - 2\tabcolsep) * \real{0.4681}}@{}}
\toprule\noalign{}
\begin{minipage}[b]{\linewidth}\raggedright
Quantitative
\end{minipage} & \begin{minipage}[b]{\linewidth}\raggedright
Qualitative
\end{minipage} \\
\midrule\noalign{}
\endhead
\bottomrule\noalign{}
\endlastfoot
Numbers used as data & Words (written/spoken) and images used as
data. \\
Seeks to identify relationships between variables to explain or predict
with the aim of generalising findings to a wider population & Seeks to
understand and interpret more local meanings. Recognises data as
gathered in a context. \emph{Sometimes} produces knowledge that
contributes to more general understandings \\
Generates ``shallow'' but broad data: not a lot fo complex detail
obtained from each participant but lots of participants take part to
generate statistical power & Generates narrow but rich data - ``thick
descriptions'' - detailed and complex accounts from each participant,
not many participants \\
Seeks consensus, norms or general patterns, often aims to reduce
diversity of responses to an average response & Tends to seek patterns
but accommodates and explores difference and divergence within the
data \\
Tends to be theory-testing and deductive & Tends to be theory generating
and \emph{inductive} (working up from the data) \\
Values detachment and impartiality (objectivity) & Values personal
involvement and partiality (subjectivity, reflexivity) \\
Has a fixed method (harder to change focus when data collection has
begun) & Method is less fixed (can accommodate a shift in focus in the
same study) \\
Can be completed quickly & Tends to take longer to complete because it
is interpretative and there is no formula \\
\end{longtable}

Look at the table above and write down answers for the following
questions:

\begin{itemize}
\tightlist
\item
  What might a `typical' quantitative study look like? What kind of data
  might be collected? What would be the aim of the study? What might
  your sample look like?
\item
  What might a'typical' qualitative study look like? What kind of data
  might be collected? What would be the aim of the study? What might
  your sample look like?
\item
  What type of research question might be best answered with a
  quantitative design?
\item
  What type of research question might be best answered with a
  qualitative design?
\end{itemize}

\section{Activity 2: Research designs
quiz}\label{activity-2-research-designs-quiz}

Choose whether a quantitative or qualitative design is most appropriate
for each of the following studies. Why did you choose each option?

\begin{enumerate}
\def\labelenumi{\arabic{enumi}.}
\tightlist
\item
  Angela is conducting a study exploring the experiences of people
  currently experiencing mental ill health. She is interested in a)
  better understanding what people with mental ill health experience
  day-to-day, and b) exploring the support mechanisms that are in place
\end{enumerate}

This study is more likely to be:

\begin{itemize}
\tightlist
\item
  \begin{enumerate}
  \def\labelenumi{(\Alph{enumi})}
  \tightlist
  \item
    Qualitative\\
  \end{enumerate}
\item
  \begin{enumerate}
  \def\labelenumi{(\Alph{enumi})}
  \setcounter{enumi}{1}
  \tightlist
  \item
    Quantitative
  \end{enumerate}
\end{itemize}

\begin{tcolorbox}[enhanced jigsaw, breakable, bottomtitle=1mm, leftrule=.75mm, opacityback=0, colbacktitle=quarto-callout-caution-color!10!white, opacitybacktitle=0.6, toprule=.15mm, colback=white, left=2mm, coltitle=black, bottomrule=.15mm, titlerule=0mm, title={Explain the answer}, toptitle=1mm, colframe=quarto-callout-caution-color-frame, arc=.35mm, rightrule=.15mm]

This study would be best conducted as a qualitative study - the
researcher is interested in exploring the experiences of people

\end{tcolorbox}

\begin{enumerate}
\def\labelenumi{\arabic{enumi}.}
\setcounter{enumi}{1}
\tightlist
\item
  John is conducting a study where he is measuring whether the amount of
  sleep that children get is related to their cognitive performance in
  class. He asks the children to wear a sleep tracker to bed, so he can
  see when they are sleeping and when they are awake.
\end{enumerate}

This study is more likely to be:

\begin{itemize}
\tightlist
\item
  \begin{enumerate}
  \def\labelenumi{(\Alph{enumi})}
  \tightlist
  \item
    Qualitative\\
  \end{enumerate}
\item
  \begin{enumerate}
  \def\labelenumi{(\Alph{enumi})}
  \setcounter{enumi}{1}
  \tightlist
  \item
    Quantitative
  \end{enumerate}
\end{itemize}

\begin{tcolorbox}[enhanced jigsaw, breakable, bottomtitle=1mm, leftrule=.75mm, opacityback=0, colbacktitle=quarto-callout-caution-color!10!white, opacitybacktitle=0.6, toprule=.15mm, colback=white, left=2mm, coltitle=black, bottomrule=.15mm, titlerule=0mm, title={Explain the answer}, toptitle=1mm, colframe=quarto-callout-caution-color-frame, arc=.35mm, rightrule=.15mm]

This study would best be conducted as a quantitative study - the data
will be quantified in numbers

\end{tcolorbox}

\begin{enumerate}
\def\labelenumi{\arabic{enumi}.}
\setcounter{enumi}{2}
\tightlist
\item
  Sally is conducting a study where she is interested in the effects of
  violence on TV on children.
\end{enumerate}

This study is more likely to be:

\begin{itemize}
\tightlist
\item
  \begin{enumerate}
  \def\labelenumi{(\Alph{enumi})}
  \tightlist
  \item
    Qualitative\\
  \end{enumerate}
\item
  \begin{enumerate}
  \def\labelenumi{(\Alph{enumi})}
  \setcounter{enumi}{1}
  \tightlist
  \item
    Quantitative
  \end{enumerate}
\end{itemize}

\begin{tcolorbox}[enhanced jigsaw, breakable, bottomtitle=1mm, leftrule=.75mm, opacityback=0, colbacktitle=quarto-callout-caution-color!10!white, opacitybacktitle=0.6, toprule=.15mm, colback=white, left=2mm, coltitle=black, bottomrule=.15mm, titlerule=0mm, title={Explain the answer}, toptitle=1mm, colframe=quarto-callout-caution-color-frame, arc=.35mm, rightrule=.15mm]

On balance, we've chosen qualitative for this, but this study could be
either quantitative or qualitative depending on how it's set up - if the
researcher is interested in a causal relationship and measures the
effects numerically, the study would be quantitative. However, if the
researcher for example does interviews with parents and wants to find
out their experiences and perceptions more in depth, the study could
also be qualitative.

\end{tcolorbox}

\begin{enumerate}
\def\labelenumi{\arabic{enumi}.}
\setcounter{enumi}{3}
\tightlist
\item
  James is conducting a study where he wants to better understand
  motivations for smoking among teenagers.
\end{enumerate}

This study is more likely to be:

\begin{itemize}
\tightlist
\item
  \begin{enumerate}
  \def\labelenumi{(\Alph{enumi})}
  \tightlist
  \item
    Qualitative\\
  \end{enumerate}
\item
  \begin{enumerate}
  \def\labelenumi{(\Alph{enumi})}
  \setcounter{enumi}{1}
  \tightlist
  \item
    Quantitative
  \end{enumerate}
\end{itemize}

\begin{tcolorbox}[enhanced jigsaw, breakable, bottomtitle=1mm, leftrule=.75mm, opacityback=0, colbacktitle=quarto-callout-caution-color!10!white, opacitybacktitle=0.6, toprule=.15mm, colback=white, left=2mm, coltitle=black, bottomrule=.15mm, titlerule=0mm, title={Explain the answer}, toptitle=1mm, colframe=quarto-callout-caution-color-frame, arc=.35mm, rightrule=.15mm]

This study would probably best be conducted as a qualitative study - the
researcher is looking at the motivations of the participants. However,
it could also possibly be a quantitative study, if the motivations are
measured and quantified on a numeric scale instead of text data.

\end{tcolorbox}

\begin{enumerate}
\def\labelenumi{\arabic{enumi}.}
\setcounter{enumi}{4}
\tightlist
\item
  Elaine is conducting a study exploring whether the amount of studying
  done in preparation for an exam predicts exam performance.
\end{enumerate}

This study is more likely to be:

\begin{itemize}
\tightlist
\item
  \begin{enumerate}
  \def\labelenumi{(\Alph{enumi})}
  \tightlist
  \item
    Qualitative\\
  \end{enumerate}
\item
  \begin{enumerate}
  \def\labelenumi{(\Alph{enumi})}
  \setcounter{enumi}{1}
  \tightlist
  \item
    Quantitative
  \end{enumerate}
\end{itemize}

\begin{tcolorbox}[enhanced jigsaw, breakable, bottomtitle=1mm, leftrule=.75mm, opacityback=0, colbacktitle=quarto-callout-caution-color!10!white, opacitybacktitle=0.6, toprule=.15mm, colback=white, left=2mm, coltitle=black, bottomrule=.15mm, titlerule=0mm, title={Explain the answer}, toptitle=1mm, colframe=quarto-callout-caution-color-frame, arc=.35mm, rightrule=.15mm]

This study would be quantitative - the researcher is interested in a
causal relationship between the variables, which you cannot achieve with
a qualitative design.

\end{tcolorbox}

\chapter{Qualitative questions}\label{ch6_qual_questions}

The purpose of this chapter is to get you thinking about how to
formulate good questions for different qualitative research methods, and
what things to avoid. Do the activities below on your own and then
discuss your answers with your group.

\section{Focus Groups}\label{focus-groups}

When you are conducting qualitative research, thinking about the
questions you are going to ask your participants is really important to
ensure your data will be rich and of good quality. You want to write
questions that are clear, open-ended and prompt the participants to
discuss/share their experiences. Here are ten tips for writing good
focus group questions:

\begin{enumerate}
\def\labelenumi{\arabic{enumi}.}
\item
  Ensure that the questions on your Focus Group Schedule are relevant
  and allow you to explore and possibly answer your research questions.
\item
  Remember you are not doing quantitative research! The beauty of
  qualitative research is that your participants can respond in a
  variety of ways and tell you all about their experiences. In order to
  allow them to do this, don't ask them questions, such as: Do people
  like A more than B? Does one thing cause another?
\item
  Make sure that your questions are open-ended, such as: What are your
  experiences of\ldots{} , Can you think of a situation when\ldots? This
  will keep conversation going. Don't ask questions that prompt a simple
  yes/no answer.
\item
  Don't write questions that are leading e.g.~Do you agree that\ldots.
  You are investigating something you don't know the answer to. Even if
  you think you might know the answer your questions should not reflect
  this.
\item
  Think about ethics! Are the questions you are asking people ethically
  problematic? If they are, it doesn't always mean you should not ask
  them, you just have to be more cautious when you do - it's important
  that we respect our participants and give them the opportunity to
  decide how much they are willing to share!
\item
  Think about time. The focus group is normally around 30-40 mins long.
  It takes quite a bit of time to get through a focus group schedule,
  especially if you have to ask additional questions. We would recommend
  you have 5 or maybe 6 questions/discussion points on your schedule.
\item
  It is useful to put a star next to/highlight the really important
  questions you want to ask. This means that if you do start to run out
  of time then you can miss some of the other questions and go straight
  to these questions.
\item
  Think about the order of the questions on the schedule. You want to
  create a discussion that is as natural as possible, and this means one
  question should follow from the next. If you structure the questions
  well at this point you may find that in the focus group you have to
  say very little in order to cover everything.
\item
  Have an opening question that makes people feel relaxed and gets them
  talking. Something that people can discuss easily and something they
  don't need to think about too much.
\item
  Additional prompts underneath the questions can be helpful for the
  person running the focus group, especially when you are new to
  qualitative research. See below for some examples from a study Ashley
  was involved with a few years ago. So, think of some probing questions
  that might allow people to discuss the question/topic in more detail.
  You might not have to use these, but they are there if you do!
\end{enumerate}

\section{Interviews}\label{interviews}

Many of the rules for good focus group questions also apply to
interviews (e.g.~writing open-ended questions that are not leading,
having prompts ready etc.). However, in an interview, the researcher is
typically one-on-one with a participant, which means that the discussion
may in some cases be a bit less structured than a focus group schedule
and may need to include more probes.

This is one section from a semi-structured interview schedule Ashley
conducted, and is provided purely to give an example of how you might
follow up your questions with prompts or probes. Please note that these
were sensitive interviews that required NHS ethical approval, and
therefore will likely be very different to any research you will be
doing during the ODL course in terms of content/tone.

\begin{longtable}[]{@{}
  >{\raggedright\arraybackslash}p{(\linewidth - 2\tabcolsep) * \real{0.7015}}
  >{\raggedright\arraybackslash}p{(\linewidth - 2\tabcolsep) * \real{0.2985}}@{}}
\toprule\noalign{}
\begin{minipage}[b]{\linewidth}\raggedright
Questions about receiving anxiety interventions
\end{minipage} & \begin{minipage}[b]{\linewidth}\raggedright
PROBES
\end{minipage} \\
\midrule\noalign{}
\endhead
\bottomrule\noalign{}
\endlastfoot
What help have you received for your anxiety? & What help did you
receive in the past? Are you currently receiving help? Who helps/helped
you? \\
Do you feel as if the help you received (or are receiving) has made any
difference to your anxiety? It may have made a positive difference, a
negative difference or no difference at all. & Anything else? \\
What was the best thing about the help you received for your anxiety? &
How did it help? Anything else? \\
What was the worst thing about the help you received for your anxiety? &
Can you expand on this? Anything else? \\
Some interventions will be adapted for the specific needs of the client.
Did the help you received seem to be modified in any way? If so, what
changes were made, in your opinion. & If there were changes, were they
helpful? What helped the most? If any changes at all could have been
made, what would your suggestions be? \\
If someone was to develop an intervention to help adults with ASD with
their anxiety, what would help the most, in your opinion? & What would
be the best type of intervention? What changes should be made? What
information could you give that would help the person giving the
intervention? \\
\end{longtable}

\section{Qualitative surveys}\label{qualitative-surveys}

Many of the principles discussed above also apply to qualitative survey
questions - you want your questions to be clear, qualitative, and
open-ended. However, due to the different response format and the fact
that you are not there to clarify or provide prompts, there are some
additional things to take into account when designing qualitative
surveys. Here are some tips for writing good qualitative survey
questions:

\begin{enumerate}
\def\labelenumi{\arabic{enumi}.}
\item
  You don't get instant feedback from your participants the same way you
  do in focus groups or interviews - hence, it's really important to
  ensure your questions are clear before starting data collection.
  Pre-testing the questions with people who are not involved in the
  project is a good way to do this.
\item
  Clarity of the questions is particularly important - make sure there
  is no ambiguity and keep your questions short.
\item
  Avoid questions which have multiple parts (as participants may just
  pick up on one) - if your question has multiple parts, it might be
  better to split it into two
\item
  Think of the order of your questions - if you start with a broad,
  over-arching question, participants might only answer this question
  and write ``see above'' in others, which reduced the richness of your
  data
\item
  You can provide examples in brackets after the question that work as
  ``prompts'' and help the participant answer the question. However, be
  careful not to steer them too much or assume what their answer will
  be. In a recent qualitative survey that Wil conducted, this was one of
  the questions: ``What do you think language barriers in teaching are,
  and what are your experiences (and/or those of your colleagues) of
  them in teaching international students? (Here you may want to reflect
  on your experience of teaching, tutoring, supervising and interacting
  with international students.)''
\item
  Include a final open question that allows the participant to share
  anything they would like to add. For example: ``Is there anything else
  you would like to add about \ldots{} ?''
\item
  Don't include too many topic questions - this will reduce the quality
  of your data as participants will get tired and more likely to drop
  out (you may have some demographic questions which are shorter and
  that's fine).
\end{enumerate}

\section{Activities}\label{activities}

\subsection{Activity 1: Which of these questions would you change and
why?}\label{activity-1-which-of-these-questions-would-you-change-and-why}

Have a look at these fictional focus group questions and write down
which ones you would change and why.

Research question: How does students' sense of belonging relate to
participation in team sports?

\begin{enumerate}
\def\labelenumi{\arabic{enumi}.}
\tightlist
\item
  How many hours of group sport do you participate in each week?
\item
  What does being a part of a team mean to you?
\item
  What does belonging mean to you?
\item
  How has your sense of belonging improved since joining your sports
  team?
\item
  I am going to go around the group one-by-one and I want everybody to
  tell me about a time when they felt that they haven't belonged to a
  team, and how that impacted them personally?
\item
  Tell me on a scale from 1-5, how much do you feel like you belong to a
  sports team
\item
  Do you like being part of a sports team?
\end{enumerate}

\begin{tcolorbox}[enhanced jigsaw, breakable, bottomtitle=1mm, leftrule=.75mm, opacityback=0, colbacktitle=quarto-callout-caution-color!10!white, opacitybacktitle=0.6, toprule=.15mm, colback=white, left=2mm, coltitle=black, bottomrule=.15mm, titlerule=0mm, title={Suggested Answers}, toptitle=1mm, colframe=quarto-callout-caution-color-frame, arc=.35mm, rightrule=.15mm]

\begin{itemize}
\tightlist
\item
  \textbf{Question 1.} This question will result in a numeric answer
  that does not produce meaningful qualitative data. Be careful not to
  accidentally ask quantitative questions!
\item
  \textbf{Question 2. and 3.} These two questions are open-ended and
  qualitative, and will likely provide rich answers. One thing to think
  about is what question to start with and where to include prompts -
  number 3 is quite big and vague, and participants may struggle to get
  started with this, especially if it's early on in the FG.
\item
  \textbf{Question 4.} This is a leading question - it implies that
  their sense of belonging has indeed \emph{improved} when this may not
  be the case. A better way to phrase this would be: \textbf{``Has your
  sense of belonging changed since joining a team? If so, how? If not,
  why not?''}
\item
  \textbf{Question 5.} This question puts participants on the spot,
  which is something you never want to do, especially with a sensitive
  question like this. Instead, the question should be given to the
  group, so participants can answer if they want to, and disclose as
  much or as little they wish to.
\item
  \textbf{Question 6.} This is a quantitative question and does not
  provide meaningful qualitative data.
\item
  \textbf{Question 7.} This is a yes/no-question. You should avoid
  phrasing questions like this, as they don't prompt rich answers.
  Instead, you may want to ask something like: \textbf{``What has been
  the best part about joining a team and why?} or \textbf{What aspects
  of being part of a team do you like? Why?} If you ask a
  yes/no-question, make sure to open it with some prompts (see example
  above for question 4).
\end{itemize}

\end{tcolorbox}

\section{Activity 2: Re-wording
questions}\label{activity-2-re-wording-questions}

Have a look at these poorly written fictional focus group questions. Try
to re-write them to improve them for clarity and the quality of data
they would give you.

Research question: Exploring international students' experiences of
their sense of belonging and homesickness

\begin{enumerate}
\def\labelenumi{\arabic{enumi}.}
\tightlist
\item
  Have you ever felt homesick while studying abroad?
\item
  How has this negatively affected you?
\item
  How often do you feel like you don't belong?
\item
  Does homesickness ever make you feel like you don't belong?
\item
  What do you do to feel better?
\end{enumerate}

When you are happy with your own answers, go through the suggested
answers below.

\begin{tcolorbox}[enhanced jigsaw, breakable, bottomtitle=1mm, leftrule=.75mm, opacityback=0, colbacktitle=quarto-callout-caution-color!10!white, opacitybacktitle=0.6, toprule=.15mm, colback=white, left=2mm, coltitle=black, bottomrule=.15mm, titlerule=0mm, title={Suggested Answers}, toptitle=1mm, colframe=quarto-callout-caution-color-frame, arc=.35mm, rightrule=.15mm]

\begin{itemize}
\tightlist
\item
  Question 1: This question prompts a yes/no answer - a better way to
  phrase this would be: \textbf{``What have your experiences of
  homesickness been while studying abroad?''}
\item
  Question 2: This question can be a bit sensitive - a better way to
  phrase this would be to say \textbf{``How do you think this has
  affected you?''} so the participant doesn't feel like they need to
  share a negative experience if they don't want to. Alternatively you
  could ask: \textbf{``What about this has been challenging to you and
  how have you managed that?''}
\item
  Question 3: This question doesn't produce meaningful qualitative data
  - if they answer ``Mostly on weekends'', you don't have much to
  analyse. A better way of phrasing the same idea it would be:
  \textbf{``How have your experiences of belongingness been?''}
\item
  Question 4: This question also prompts a yes/no answer. This is not
  necessarily a very good question in the first place because you're not
  directly looking at a causal relationship the same way you would in a
  quantitative study, and it becomes easily a leading question assuming
  there is a relationship between the two. A better question would be:
  \textbf{``How do you think your experiences of belongingness and
  homesickness compare or connect to each other?''}
\item
  Question 5: This question is a bit ambiguous and it's not really clear
  what you are asking - feel better about what? This also assumes that
  homesickness is always a negative state, which it might not be to the
  participants. A better way to phrase this would be: \textbf{``How do
  you like to manage your feelings of homesickness?''} or \textbf{``What
  do you do when you are feeling homesick?''}
\end{itemize}

\end{tcolorbox}

\section{OPTIONAL: Research skills session - Developing
questions}\label{optional-research-skills-session---developing-questions}

Ashley has recorded \textbf{optional} research skills sessions on
different aspects of qualitative research. These sessions take a little
bit deeper dive into some of the topics and may be particularly useful
if you are planning on doing a qualitative dissertation next year.

\href{https://echo360.org.uk/media/56e4abb3-e58c-4560-897b-18d2eeb91833/public?autoplay=false&automute=false&startTimeMillis=0}{Watch
the video here}

\subsection{Overall structure}\label{overall-structure}

It's important to think about the overall structure of an interview or
focus group schedule. A good structure to use is: A warm up question,
the main questions you want to ask, and an exit question. It is likely
particularly important to use a warm up/ice breaker question for a focus
group, especially if participants don't know each other particularly
well.

\textbf{Warm up Question}

The aim is to allow participants to introduce themselves, ensure
everyone understands the procedure, and to encourage people to feel
comfortable talking. Previous projects have used photos, cartoons, and
media, or asked participants to generate meanings of the terms used in
the discussion (e.g., ``poor sleep'',''well-being'') to get people
talking. \emph{What is your idea for an ice breaker that might lead
naturally into the main discussion?}

\emph{What makes a good/poor icebreaker?}

Thinking about the points above: \emph{What do you think of the
following ice-breakers?}

\begin{itemize}
\tightlist
\item
  \textbf{A FG on lived experience during COVID-19 of international
  students:} What is your favourite meal from your home country?
\item
  \textbf{A FG on coping strategies in students:} What do you think is
  the best way to reduce stress for students?
\item
  \textbf{A FG on feelings of belonging in a work context:} What do you
  think about the graphic below? Does it resonate in any way?
\end{itemize}

\begin{figure}[H]

{\centering \includegraphics[width=1\linewidth,height=\textheight,keepaspectratio]{images/belonging.png}

}

\caption{Belonging icebreaker}

\end{figure}%

\textbf{Main Questions}

Bear in mind that the discussion may wander off track or focus on one
aspect. Consider which part of the research question you think is most
important to get people to answer. Start with this and be prepared to
come back to it at the end if you need to. If you have more sensitive
questions in your research schedule, they should be included later on,
so people answer them when they are more comfortable.

\textbf{What should the questions be like?} Avoid asking questions that
ask two or three things in one sentence (e.g., about behaviour and
impact on mood). Break it down into simpler questions. Limit the number
of direct questions that elicit yes/no answers. These can be useful to
gather information but are less helpful in generating a discussion. Add
in probes/prompts to elicit further information if required.

Try to:

\begin{enumerate}
\def\labelenumi{\arabic{enumi}.}
\item
  Frame questions in an open way that will facilitate discussion
  (e.g.~`do you use facebook?' is closed (can be answered as yes/no)
  whereas `can you tell me about your experiences with social media?' is
  a bit more open - you could then follow up with prompts to
  specifically ask about types if you want to get that information
\item
  Frame questions in a way that is not leading. For example, `how did
  lockdown negatively impact on your ability to study?' assumes that it
  did, when this may not be the case. A better question would be `'Do
  you think that lockdown restrictions impacted on students' experiences
  of studying? If so, how?''
\item
  If you have sensitive questions, try to frame questions in a way that
  is not personal (i.e.~you give people a chance to answer related to
  themselves, but they can also choose not to). This apporach also
  invites people who may not have lived experience to respond,
  e.g.~``How do you take care of yourself when you are anxious'' would
  be better phrased as ``How might students support themselves when they
  are feeling anxious?'' This is more important for focus groups, as
  people might feel put on thespot and go beyond their comfort levels.
\item
  Not ask lots of different things within the one question, e.g.~``Can
  you tell us about your experiences of studying in Semester 1, how it
  made you feel and how that affects your studying in Semester 2?'' You
  could ask about experiences of studying generally, and then could
  follow up with specific sub-questions asking them about semester 1 and
  2, if that is something that was important to you.
\end{enumerate}

\textbf{Closing question}

The aim here is to give people an opportunity to add anything that might
have been missed in the discussion/to wrap things up.

Examples:

\begin{itemize}
\tightlist
\item
  Thank you very much for all the great points you had today!Before we
  wrap things up, you have any final points that you would like to make,
  that perhaps wasn't covered in the interview?
\item
  Today we covered X, Y, and Z in the focus group. Is there anyone who
  has any final thoughts about what we covered here today that they'd
  like to share?
\item
  That's us finished with the main part of the focus group. We'd like to
  end by giving you all an opportunity to add anything that you think is
  important but has not been covered in the focus group so far. If you'd
  like to add something, please press raise hand on teams and we can
  make sure we go round those people before we finish up!
\end{itemize}

\subsection{Braun \& Clarke (2013)
activity}\label{braun-clarke-2013-activity}

Here, we have an adapted version of an interview schedule that is
included in the Braun \& Clarke (2013) book, on the topic of Partner
Relationships. This is purposefully poor. What problems can you spot?

\begin{enumerate}
\def\labelenumi{\arabic{enumi}.}
\tightlist
\item
  If you are not already married, are you planning to get married within
  the next few years?
\item
  I love going to the cinema with my partner, do you?
\item
  Do you have a good time with your partner?
\item
  Would you go to Relate if you had problems in your relationship?
\item
  When and how did you meet your partner?
\item
  Do you earn more money that your partner?
\end{enumerate}

\begin{tcolorbox}[enhanced jigsaw, breakable, bottomtitle=1mm, leftrule=.75mm, opacityback=0, colbacktitle=quarto-callout-caution-color!10!white, opacitybacktitle=0.6, toprule=.15mm, colback=white, left=2mm, coltitle=black, bottomrule=.15mm, titlerule=0mm, title={Show me the solution}, toptitle=1mm, colframe=quarto-callout-caution-color-frame, arc=.35mm, rightrule=.15mm]

\textbf{In general} Overall, the interview schedule lacks coherence --
it covers a wide range of different topics, and it's not clear what the
focus of the interview is. The ordering of the questions is also very
poor (very sensitive questions early in the interview; moving straight
from questions about positive aspects of the relationship to
relationship difficulties). The researchers should decide on the focus
of their research, and their research question, and design the interview
schedule with these in mind. The schedule should also be reordered so it
starts with less sensitive questions (these should be moved to later in
the interview) and the questions flow logically. The interview should
also end with a general `clean-up' question (e.g., `Is there anything
else you would like to add?'). (NB. Our answers to this exercise assume
that the research is being conducted in Britain)

\textbf{Question 1}

\emph{Problems} First, this is a closed question which does not
encourage participants to provide rich and detailed answers. Second,
unless the researcher already has quite a bit of information about the
participant, it is not an ideal opening interview question because it
makes a number of problematic and potentially alienating assumptions,
which may damage rapport. For example, unless the study is specifically
focused on heterosexual couples (`partner relationships' suggests a
inclusive focus), or the researcher already knows how the participants
identify their sexuality, assuming that people are married, or asking
whether they are planning to marry, when marriage is not legally
available to same-sex couples, reinforces heteronormativity.
Furthermore, unless the research is specifically focused on marriage,
asking whether people are `planning to get married' assumes a particular
model of relationships whereby marriage (and civil partnership) is
treated as the `gold standard' of relationships, the ideal end-point of
a couple's relational journey, and not everyone views marriage in this
way (some couples are critical of marriage and choose actively not to
get married).

\emph{Solutions} The question needs to be reworded so that it is
inclusive of same-sex and heterosexual relationships and does not `buy
into' assumptions about marriage as gold standard of relationships. The
question should also be reworded as an open question and probably moved
to later in the interview.

\textbf{Question 2}

\emph{Problems} Again, a closed question and extremely leading -- it
invites the participant to agree with the researcher's preferences! This
is not the best way of accessing participants' perspectives.

\emph{Solutions} Reword the question so that it is open and not leading
(do not use the researcher's preferences in this way).

\textbf{Question 3}

\emph{Problems} Again closed, and rather vague (what constitutes a `good
time'?).

\emph{Solutions} Reword the question so that it is open and more
concrete.

\textbf{Question 4}

\emph{Problems} Again closed, and assumes that the participants know
what Relate is (a UK charity that that offers relationship counselling
and sex therapy to heterosexual and same-sex couples experiencing
relationship difficulties). Making assumptions about what the
participants know (and don't know) can alienate participants (making
them feel stupid for not knowing) and damage rapport. Moreover, Relate
is not the only provider of relationship counselling in the UK.

\emph{Solutions} Replace this with an open question and one that does
not make assumptions about the participants' knowledge (e.g.,
`relationship counselling' is a good substitute for Relate).

\textbf{Question 5}

\emph{Problems} Again, a double-barrelled question and also one that
would work better earlier in the interview.

\emph{Solutions} Reword the question to remove the double-barrelled
element (decide what is more important to know when {[}how long ago{]}
they met their partner or how they met them?). This would work better as
the opening interview question (it makes little sense to ask questions
1-6 before asking this question) because it is a relatively `gentle'
question and possibly one that participants would be expecting to answer
in an interview on partner relationships. It provides participants with
an opportunity to `tell the story' of their current relationship at the
start of the interview.

\textbf{Question 6}

\emph{Problems} Again a closed question and also a very sensitive one!

\emph{Solutions} Again, this needs to be replaced with a more open
question that invites detailed responses. The researcher should reflect
on what they want to know and how this question relates to their
research question.

\end{tcolorbox}

\subsection{Further things to consider with qualitative
surveys}\label{further-things-to-consider-with-qualitative-surveys}

\begin{enumerate}
\def\labelenumi{\arabic{enumi}.}
\tightlist
\item
  No opportunity for clarity, so wording has to be clear
\item
  Similar points to above in terms of not asking too many different
  things, having a sensible structure
\item
  No rapport, unlikely to be an `ice breaker' (although might have an
  intro question)
\item
  People might be more comfortable answering sensitive questions as they
  can be anonymous
\item
  Try not to have lots of questions as your participants will likely
  lose patience - what are the key things you want to know?
\item
  Can use branching e.g.~in Qualtrics
\end{enumerate}

\chapter{Data collection}\label{ch7_data_collection}

This chapter runs through some of the practical issues you need to think
about when collecting data for a qualitative study. Do the activities
below on your own and then discuss your answers with your group.

\section{Considering ethics}\label{considering-ethics}

\textbf{Informed consent}

Before and data collection takes place, you must provide your
participants with an information sheet and consent form. You can find
more information about this in the \hyperref[ch9_ethics]{Ethics, Rigour
and Reflexivity chapter}

You must give everyone participating the Information sheet and consent
form before the focus group/interview takes place. Completed consent
forms should be returned to the researcher and you should verbally
confirm consent before you start any data collection. In a survey,
participants should complete consent form before seeing any survey
questions.

\textbf{Confidentiality in Focus groups}

Remember that everyone (both those participating and running) the focus
group have a responsibility to maintain confidentiality of the
discussion.

Discuss some `ground rules' at the start with the group, and also go
through the Information Sheet and consent form at the start, to ensure
everyone is aware of their responsibilities. You may also want to
include an ice breaker to make your participants feel more comfortable.

\section{Focus group structure}\label{focus-group-structure}

There are a number of stages to a focus group, and it's not always
exactly the same, but here is an example structure for a focus group:

\begin{itemize}
\tightlist
\item
  \textbf{Stage 1: Setting the scene.} The topic of the project should
  be introduced by the facilitator, and everyone agrees to take part. It
  must be made clear to the people in your focus group a) that they do
  not have to answer any questions that they do not feel comfortable
  with, b) highlight that the session will be recorded (if conducting
  your study online, also make sure to let your participants know they
  have a choice to have their camera on or off, but that the mic must be
  used), c) what will happen with the recordings (i.e.~that they will be
  kept safe on University-approved servers and the transcript
  anonymised), d) that everyone has a responsibility to maintain
  confidentiality of the topics discussed in the group and e) that they
  are free to leave the session at any time. Run through the Information
  Sheet and Consent Form here, and discuss any ground rules you have. I
  would also raise the point here that there are a few questions to get
  through, and therefore you might have to step in to help manage time.
\item
  \textbf{Stage 2: Recording.} Remember to start recording and let your
  participants know when you do it. This is very important.
\item
  \textbf{Stage 3: Introductions.} Go around the table and ask everyone
  to introduce themselves. Start with yourself and the note-taker first.
  Everyone should state their first name clearly.
\item
  \textbf{Stage 4: Ice-breaker.} Ask a warm-up question or ice-breaker.
  The aim is to allow participants to introduce themselves, ensure
  everyone understands the procedure, and to encourage people to feel
  comfortable talking. Previous projects have used photos, cartoons, and
  media, or asked participants to generate meanings of the terms used in
  the discussion (e.g., ``poor sleep'',''well-being'') to get people
  talking.
\item
  \textbf{Stage 5: Discussion.} This is where you ask your main
  questions, and follow-up questions. It is a good idea to prepare
  around five or six main questions, and perhaps highlight three or four
  that are most important to focus on. The discussion may follow a
  different order to that on the focus group schedule - this is fine.
\item
  \textbf{Stage 6: Facilitating the discussion.} The facilitator should
  be prepared to help steer conversation back on track, and balance
  discussion so that everyone gets a chance to contribute. Have a few
  specific prompts that relate to each question and/or some general
  prompts to keep the discussion flowing. What are you going to do if
  no-one participates? Or if one person is contributing? For
  example\ldots{} ``has anyone else had a similar experience? Does
  anyone have a different view?'' At the same time be ethical and avoid
  targeting individual people. Give participants the opportunity to
  describe their interpretations and take ownership. E.g., ``How
  would\ldots make you feel'', ``Have you ever noticed that\ldots{}''
  ``Think back to when\ldots{}'', ``What was most important to you
  about\ldots?''
\item
  \textbf{Stage 7: Ending the discussion.} It's good to end the
  discussion by going round the group and asking each participant if
  they have anything they want to add to the discussion/sum up their
  thoughts on what was discussed.Remember to allow focus group members
  the chance to confirm/withdraw consent for their anonymised
  contributions to be used.
\end{itemize}

\section{Facilitation and
note-taking}\label{facilitation-and-note-taking}

\textbf{Facilitation/moderation}

The facilitator (sometimes also known as the moderator) will be the
person taking the lead during the focus group. They will be the one
asking the questions of the participants, and the one facilitating the
discussion. As you are aiming to mimic a naturalistic conversation, you
have to build rapport. Ensure you allow everyone time to introduce
themselves, and contribute to the ice-breaking phase by saying a bit
about yourself.

A good facilitator shares the perspective of the people they are talking
to throughout the focus group. The moderator should not be listening out
for good quotes and should not be thinking about what to analyse. The
role of the facilitator is to listen to the participants, and understand
their points of view.

\textbf{Note-taking}

It is a good idea to have a note-taker there. Hopefully there will be no
IT issues that mean that the recording doesn't work! Therefore, the
note-taker should remind the facilitator about recording if required,
and note down any impressions/initial thoughts around themes/good
quotes. Also, you could note down how well the questions worked - were
there any that participants didn't understand? It should be clear from
the recording who is contributing when; it is useful to note down here
and there who is speaking at various points, to ensure it matches up
later on.

\section{Facilitating focus groups: How to manage the
group?}\label{facilitating-focus-groups-how-to-manage-the-group}

Focus groups are designed to be an informal discussion in a moderated
setting. This allows for the data collection to be flexible and occur in
a collective, more natural way - however, moderating a discussion in a
group setting can also be challenging. Below we have some top tips for
how to manage focus groups.

\begin{enumerate}
\def\labelenumi{\arabic{enumi}.}
\item
  Make sure your participants are comfortable: establish ground rules
  and expectations at the start, ensure that you are interested in their
  perceptions and experiences and that there are no right or wrong
  answers, start with an ice-breaker to warm your participants up and
  make sure nobody is put on the spot and participants don't feel
  pressured to answer anything they do not wish to answer.
\item
  Manage the conversation effectively: make sure all participants have
  the opportunity to contribute if they wish and one participant doesn't
  dominate the whole conversation, ensure that conversation doesn't
  stray off topic or get stuck on one topic for too long. Give prompts
  if the participants struggle to get started on a question.
\item
  Stay neutral and \emph{listen}: be careful not to provide your own
  opinion or views (this is surprisingly easy to do accidentally!) -
  this is about your participants and their experiences. Focus on
  listening to the participants and clarifying/prompting further if
  needed, give your participants time to gather their thoughts (learn to
  be comfortable with silent moments and don't fill them in with your
  opinions).
\end{enumerate}

View below some situations you may encounter when running a focus group:
before checking the answer, have a think how you would manage these
scenarios and discuss the answers with your group.

\begin{enumerate}
\def\labelenumi{\arabic{enumi}.}
\tightlist
\item
  One/some participants are very quiet and one participant dominates the
  conversation
\end{enumerate}

\begin{tcolorbox}[enhanced jigsaw, breakable, bottomtitle=1mm, leftrule=.75mm, opacityback=0, colbacktitle=quarto-callout-caution-color!10!white, opacitybacktitle=0.6, toprule=.15mm, colback=white, left=2mm, coltitle=black, bottomrule=.15mm, titlerule=0mm, title={Suggested Solutions}, toptitle=1mm, colframe=quarto-callout-caution-color-frame, arc=.35mm, rightrule=.15mm]

\begin{itemize}
\tightlist
\item
  The ice-breaker/easy opening questions are perhaps most important for
  quieter members of the group, to help make them feel more comfortable.

  \begin{itemize}
  \tightlist
  \item
    If you are running the focus group online, it might be a good idea
    to use the Raise hand feature - this would allow quieter members to
    show that they want to contribute, but without trying to gauge when
    to jump in the conversation. You can then take points in order of
    when they raised their hands.
  \item
    If some people have contributed to part of the discussion and others
    haven't, open it up to the rest by asking if anyone who hasn't
    responded yet has anything to add. However, don't put people on the
    spot by asking them directly.
  \item
    Make sure to distribute turns equally
  \end{itemize}
\end{itemize}

\end{tcolorbox}

\begin{enumerate}
\def\labelenumi{\arabic{enumi}.}
\setcounter{enumi}{1}
\tightlist
\item
  The discussion is stuck on one topic
\end{enumerate}

\begin{tcolorbox}[enhanced jigsaw, breakable, bottomtitle=1mm, leftrule=.75mm, opacityback=0, colbacktitle=quarto-callout-caution-color!10!white, opacitybacktitle=0.6, toprule=.15mm, colback=white, left=2mm, coltitle=black, bottomrule=.15mm, titlerule=0mm, title={Suggested Solutions}, toptitle=1mm, colframe=quarto-callout-caution-color-frame, arc=.35mm, rightrule=.15mm]

\begin{itemize}
\tightlist
\item
  Remind the members of the group that there are a number of different
  questions to explore, and that you are mindful of time, so it might be
  best to leave that topic for now, and you can return if there's time
  at the end.
\end{itemize}

\end{tcolorbox}

\begin{enumerate}
\def\labelenumi{\arabic{enumi}.}
\setcounter{enumi}{2}
\tightlist
\item
  The discussion goes off topic
\end{enumerate}

\begin{tcolorbox}[enhanced jigsaw, breakable, bottomtitle=1mm, leftrule=.75mm, opacityback=0, colbacktitle=quarto-callout-caution-color!10!white, opacitybacktitle=0.6, toprule=.15mm, colback=white, left=2mm, coltitle=black, bottomrule=.15mm, titlerule=0mm, title={Suggested Solutions}, toptitle=1mm, colframe=quarto-callout-caution-color-frame, arc=.35mm, rightrule=.15mm]

\begin{itemize}
\tightlist
\item
  The facilitator has to be mindful of this and be ready to steer the
  conversation back on topic if needed.
\item
  Remind the members of the group that there are a number of different
  questions to explore, and that you are mindful of time.
\item
  Gently remind participants of the question that was asked. If running
  the FG online, you may wish to put your questions on the chat or
  present them with screen sharing so the participants can see the
  question.
\end{itemize}

\end{tcolorbox}

\begin{enumerate}
\def\labelenumi{\arabic{enumi}.}
\setcounter{enumi}{3}
\tightlist
\item
  Participants are not comfortable answering the questions in the
  language the focus group is conducted in
\end{enumerate}

\begin{tcolorbox}[enhanced jigsaw, breakable, bottomtitle=1mm, leftrule=.75mm, opacityback=0, colbacktitle=quarto-callout-caution-color!10!white, opacitybacktitle=0.6, toprule=.15mm, colback=white, left=2mm, coltitle=black, bottomrule=.15mm, titlerule=0mm, title={Suggested Solutions}, toptitle=1mm, colframe=quarto-callout-caution-color-frame, arc=.35mm, rightrule=.15mm]

\begin{itemize}
\tightlist
\item
  Send questions in advance to your participants so they can prepare
\item
  Consider if the FG could be conducted in a different language
  (however, be mindful that this creates additional work with
  transcription and translations and requires the facilitator to be
  fluent in the language)
\item
  Give participants time to think and answer, and focus on making
  everybody feel comfortable at the start of the FG.
\end{itemize}

\end{tcolorbox}

\begin{enumerate}
\def\labelenumi{\arabic{enumi}.}
\setcounter{enumi}{4}
\tightlist
\item
  One participant looks uncomfortable and does not contribute
\end{enumerate}

\begin{tcolorbox}[enhanced jigsaw, breakable, bottomtitle=1mm, leftrule=.75mm, opacityback=0, colbacktitle=quarto-callout-caution-color!10!white, opacitybacktitle=0.6, toprule=.15mm, colback=white, left=2mm, coltitle=black, bottomrule=.15mm, titlerule=0mm, title={Suggested Solutions}, toptitle=1mm, colframe=quarto-callout-caution-color-frame, arc=.35mm, rightrule=.15mm]

\begin{itemize}
\tightlist
\item
  Try to create opportunities for everybody to speak, but remember that
  you should not point fingers at individual participants
\item
  You cannot force anybody to speak, and sometimes participants might
  feel uneasy.
\item
  Remember that they don't need to answer any of the questions if they
  don't want to - it is unethical to force your participants to answer.
\end{itemize}

\end{tcolorbox}

\begin{enumerate}
\def\labelenumi{\arabic{enumi}.}
\setcounter{enumi}{5}
\tightlist
\item
  Nobody answers the question you asked
\end{enumerate}

\begin{tcolorbox}[enhanced jigsaw, breakable, bottomtitle=1mm, leftrule=.75mm, opacityback=0, colbacktitle=quarto-callout-caution-color!10!white, opacitybacktitle=0.6, toprule=.15mm, colback=white, left=2mm, coltitle=black, bottomrule=.15mm, titlerule=0mm, title={Suggested Solutions}, toptitle=1mm, colframe=quarto-callout-caution-color-frame, arc=.35mm, rightrule=.15mm]

\begin{itemize}
\tightlist
\item
  Remember that this is the first time the participants may see the
  question, and they may need time to gather their thoughts
\item
  Consider sending questions to your participants in advance so they
  know what to expect.
\item
  Don't be uncomfortable with a bit of awkward silence (particularly if
  running your FG online) and don't fill the space with your opinions -
  let the participants take the lead.
\item
  Prepare prompts/probes that you can ask to help the participants get
  started.
\item
  You can ask if anybody would like to go first (don't put an individual
  participant on the spot)
\end{itemize}

\end{tcolorbox}

\begin{enumerate}
\def\labelenumi{\arabic{enumi}.}
\setcounter{enumi}{6}
\tightlist
\item
  You are not sure what the participant meant/the participant answers
  very shortly
\end{enumerate}

\begin{tcolorbox}[enhanced jigsaw, breakable, bottomtitle=1mm, leftrule=.75mm, opacityback=0, colbacktitle=quarto-callout-caution-color!10!white, opacitybacktitle=0.6, toprule=.15mm, colback=white, left=2mm, coltitle=black, bottomrule=.15mm, titlerule=0mm, title={Suggested Solutions}, toptitle=1mm, colframe=quarto-callout-caution-color-frame, arc=.35mm, rightrule=.15mm]

\begin{itemize}
\tightlist
\item
  You can ask if the participant would like to clarify the aspect of the
  answer you are not sure about.
\item
  You can ask the participant to elaborate further (they may not want
  to, which is fine!). If you do this, make sure that you are specific
  about what you want them to elaborate on (for example, ``that's really
  interesting, would you want to tell me a little bit more about why you
  think XXX?)
\item
  If the participant does not know how to elaborate, you can ask if
  anybody else wants to add anything. Sometimes this will prompt the
  participant to build on their answer as a response to what the other
  participants are saying.
\end{itemize}

\end{tcolorbox}

\section{Facilitating interviews}\label{facilitating-interviews}

Interviews are very different to focus groups, as the data collection is
typically set up in a one-on-one situation. This means the types of
questions you should ask, the types of answer you are likely to get, and
the practical issues are quite different to focus groups. Here are some
top tips for how to conduct an interview.

\begin{itemize}
\tightlist
\item
  Make an effort to develop rapport and make the participant feel
  comfortable.
\item
  If you are doing a sensitive topic, pay attention to how your
  participant is presenting and ask them if they need a break or if they
  would like to continue.
\item
  Focus on the person rather than taking notes - are you giving them
  encouragement to keep speaking (e.g.~nods etc.)?
\item
  If doing a semi-structured interview, pay attention to whether the
  participant brings something up that could/should be followed up with
  an addional question
\item
  Not all interviews will be the same length - some people are very
  concise in how they answer whereas others talk a lot more so the
  overall interview length will vary
\item
  Try to end on a positive question if you can - it can help to lift the
  mood for the participant and leave them in a good place
\end{itemize}

\section{GDPR-compliance and data
management}\label{gdpr-compliance-and-data-management}

GDPR: general data protection regulation. It is the most up to date data
protection law and your data collection and storage must comply to the
GDPR principles.

See the University guidelines on:

\begin{itemize}
\tightlist
\item
  \href{https://www.gla.ac.uk/myglasgow/it/informationsecurity/confidentialdata/}{handling
  confidential data}\\
\item
  \href{https://www.gla.ac.uk/myglasgow/dpfoioffice/gdpr/}{GDPR}
\end{itemize}

Make sure that:

\begin{itemize}
\tightlist
\item
  You have a data management plan in place before data collection (how
  and where are you storing the data, who will have access to it, how
  long do you need to store it for, what steps of confidentiality and
  anonymisation are you taking?)
\item
  Participants are made aware of your data management plan prior to any
  data collection in the information sheet and consent form
\item
  Data should be retained until after completion of the project on a
  password-protected encrypted drive and a secure database accessible
  only to the research team
\item
  Only use University-approved systems to collect, manage and store data
  - \textbf{do not use external providers (such as Google or Dropbox)},
  as they may not be GDPR compliant
\item
  If you share the data, your participants should be informed about this
  prior to taking part, and you should only share the anonymised data
  files
\item
  \textbf{Do not share data via social media, Google drive or personal
  email which is not secure.}
\end{itemize}

\section{OPTIONAL: Research skill session - data
collection}\label{optional-research-skill-session---data-collection}

Ashley has recorded \textbf{optional} research skills sessions on
different aspects of qualitative research. These sessions take a little
bit deeper dive into some of the topics and may be particularly useful
if you are planning on doing a qualitative dissertation next year.

\href{https://echo360.org.uk/media/1f84181c-364e-47e3-a61e-415e4ec5a368/public?autoplay=false&automute=false&startTimeMillis=0}{Watch
the video here}

\subsection{General tips}\label{general-tips}

\begin{itemize}
\tightlist
\item
  Be organised - get ethics in early etc., keep track of potential
  participants
\item
  Think through all the different steps involved as you will need to
  refer to this in your ethics application - how do you plan to recruit,
  collect the data, store the data? Before you start collecting data, it
  would be good to have a checklist with everything you need to consider
\item
  Only collect the demographics that you need - if you don't need to
  know how many siblings they have or what time they go to bed, then
  don't ask for that information
\item
  Be sensitive to who your participants are - you might ask different
  questions in a focus group (where people are in front of others) than
  you would in an online questionnaire (where it is easier to withdraw
  if uncomfortable).
\item
  Think about how you keep that data safe/anonymous
\item
  Know in advance how many participants you need
\item
  Know how you plan to analyse the data
\item
  Consider what you can do to make the experience as comfortable as
  possible for participants (e.g.~sending out questions in advance is
  particularly helpful for neurodivergent people/people who do not speak
  English as a first language).
\item
  For interviews and focus groups - reassert consent before you start \&
  always remember to press record! Also it's normal to feel a bt nervous
  - try to enjoy the process!
\item
  For all, pilot the questions you ask - make sure that people take from
  them what you think
\item
  When recruiting, think about who your target group is and it's good to
  justify why you are looking for a particular demographic. For example,
  do you definitely require someone to have a particular diagnosis?
\end{itemize}

\subsection{Practicalities of running focus
groups}\label{practicalities-of-running-focus-groups}

\begin{itemize}
\tightlist
\item
  Establish ground rules - esp.~confidentiality
\item
  Set expectations - ask participants to respond to each other rather
  than each person individually responding to the question (if possible,
  this is not always the case)
\item
  Moderator has to manage the group as best they can
\item
  Wrapping things up.
\item
  Can be helpful to have a notetaker to note down anything
  interesting/if doing it in person, to note down order that people
  speak in. Allows the moderator to focus on discussion
\end{itemize}

\textbf{Tricky situations}

\begin{itemize}
\tightlist
\item
  How do you develop rapport/make people feel more comfortable?
\item
  How do you manage tricky dynamics within a group?
\item
  What if no-one speaks/one person dominates?
\item
  What if people go off topic?
\item
  How to keep people to time?
\item
  Giving everyone who wants it an opportunity to speak - if online could
  ask people to raise their hand
\end{itemize}

\subsection{Practicalities of running
interviews}\label{practicalities-of-running-interviews}

Some of the things that we discussed in focus groups are relevant here,
but specifically when thinking about interviews:

\begin{itemize}
\tightlist
\item
  Try to develop rapport and make the participant feel comfortable
\item
  Try to focus on the person rather than taking notes - are you giving
  them encouragement to keep speaking (e.g.~nods etc.)?
\item
  If doing a semi-structured interview, try to notice if the participant
  brings something up that could/should be followed up with an
  additional question
\item
  It's usual for some interviews to be shorter than others - some people
  are very concise in how they answer whereas others talk a lot more
  (and this may be on topic or not)
\item
  Be aware of how someone is presenting - this is less relevant for
  coursework and more for future projects - if you are doing quite a
  sensitive topic, it's good practice to reassert consent and to ask
  people if they need a break for a little while/if they are okay to
  continue.
\item
  Try to end on a positive question if you can - it can help to lift the
  mood for the participant and leave them in a good place
\end{itemize}

\subsection{Practicalities of running qualitative
surveys}\label{practicalities-of-running-qualitative-surveys}

\begin{itemize}
\tightlist
\item
  Think about whether there would be any possibility for the questions
  to be misinterpreted - you won't have a chance to correct people if
  they interpret something wrong (this is where doing a pilot can be
  useful)
\item
  Try not to have too many questions - what are the key things you want
  to know? People will be less likely to respond with long, in depth
  answers if there are a lot of questions to answer.
\item
  It might be good to give an indication of how many questions they have
  to go through.
\item
  Easier to keep data anonymous than it is with other forms of data
  collection above - however still be sensitive to locations, names,
  personal experiences as these will have to be anonymised.
\end{itemize}

\chapter{Data analysis}\label{ch8_data_analysis}

This chapter has activities to support you with Thematic Analysis. This
chapter, alongside with Ashley's lectures, will talk you through the six
phases required for TA; we advise using Braun \& Clarke (2021) to
structure your data analysis.
\href{https://go.exlibris.link/Z9qxllzH}{This can be accessed through
the library}.

\textbf{Overview of the steps of TA}

\begin{enumerate}
\def\labelenumi{\arabic{enumi}.}
\tightlist
\item
  Familiarisation (p.42 in the Braun \& Clarke (2021) book)
\item
  Doing coding (p.59)
\item
  Generating initial themes (p.78)
\item
  Developing and Reviewing themes (p.97)
\item
  Refining, Defining and Naming themes (p.108)
\item
  Writing up your analysis (p.118)
\end{enumerate}

\section{Coding (Phases 1-2)}\label{coding-phases-1-2}

Here, we'll focus on Phases 1 and 2 of the process. Familiarisation
starts on p.42 and Coding starts on p.59 of the Braun \& Clarke (2021)
book. We also covered how to do Phases 1 and 2 of the analysis in the
Thematic Analysis lecture in Semester 1.

\subsection{Phase 1: Familiarisation
(p.42)}\label{phase-1-familiarisation-p.42}

\begin{itemize}
\tightlist
\item
  Transcribing (or correcting/anonymising your transcript) helps you to
  become more familiar with your data. Although you will not do this in
  this project, you may do so in your dissertation).
\item
  Read over the transcript, but do it in an `active' way. This means
  reading while searching for meanings and patterns
\item
  Note down anything of interest that comes up, either in a particular
  part of a transcript or thinking about the overall dataset
\end{itemize}

\begin{tcolorbox}[enhanced jigsaw, breakable, bottomtitle=1mm, leftrule=.75mm, opacityback=0, colbacktitle=quarto-callout-warning-color!10!white, opacitybacktitle=0.6, toprule=.15mm, colback=white, left=2mm, coltitle=black, bottomrule=.15mm, titlerule=0mm, title={Practical Steps}, toptitle=1mm, colframe=quarto-callout-warning-color-frame, arc=.35mm, rightrule=.15mm]

Read through your transcript and take notes and/or make comments in the
margins. Generate a list of ideas: what is interesting about the
discussion? Are there any contradictions that you spot? What patterns do
you observe?

\end{tcolorbox}

\subsection{Phase 2: Generating initial codes
(p.59)}\label{phase-2-generating-initial-codes-p.59}

\begin{itemize}
\tightlist
\item
  Codes identify a feature of the data that seems interesting
\item
  The process of coding involves going through the transcript and
  thinking about the data, and picking out what seems
  interesting/important.
\item
  Try to apply short, specific, labels to each part of the transcript as
  you go through it.
\item
  Coded data is specific (themes, which come later, are broad) -you will
  likely end up with quite a lot of different codes.
\item
  Code for as many patterns as you can at this point - it won't all make
  it to the next stage, but it's good if you pick up as much as you can
  at this stage.
\end{itemize}

\begin{tcolorbox}[enhanced jigsaw, breakable, bottomtitle=1mm, leftrule=.75mm, opacityback=0, colbacktitle=quarto-callout-warning-color!10!white, opacitybacktitle=0.6, toprule=.15mm, colback=white, left=2mm, coltitle=black, bottomrule=.15mm, titlerule=0mm, title={Practical Steps}, toptitle=1mm, colframe=quarto-callout-warning-color-frame, arc=.35mm, rightrule=.15mm]

Insert labels beside quotes - what is interesting about the quote? What
does it tell us in relation to your research question?

\end{tcolorbox}

\subsection{How to actually do the coding
(logistically)?}\label{how-to-actually-do-the-coding-logistically}

There are many different ways to approach this. We ask you to use one of
two methods for your report (hand-coding or coding using Word). The
reasons for this are that we need you to evidence the full process of
your analysis, and this is included as part of the learning criteria for
this assignment.

Option 1: Hand-coding

\begin{figure}[H]

{\centering \includegraphics[width=1\linewidth,height=\textheight,keepaspectratio]{images/handcoded.jpg}

}

\caption{Example of a hand-coded transcript}

\end{figure}%

Option 2: Code using Word, by adding comments

\begin{figure}[H]

{\centering \includegraphics[width=1\linewidth,height=\textheight,keepaspectratio]{images/commentcoded.png}

}

\caption{Example of a transcript coded using comments on Microsoft Word}

\end{figure}%

However, if you do a qualitative project for your dissertation, other
options will be open to you, so please do ask your supervisor what might
be best!

\subsection{Activities}\label{activities-1}

\subsubsection{Activity 1: Example of
coding}\label{activity-1-example-of-coding}

Have a look at
\href{https://gla-my.sharepoint.com/:w:/g/personal/ashley_robertson_glasgow_ac_uk/EdF3fJaWwyJJuVkP-jVOVPcBlNGsfnv04rsWlavw_pyowA?e=QgIRAO}{this
example of coding (using Thematic Analysis) from Braun \& Clarke (2013)}

Think about the following:

Do you agree/disagree with the codes that the authors have chosen? Which
ones, and why? Are there any additional codes that you think might be
worthy of consideration?

\subsubsection{Activity 2: Try it
yourself!}\label{activity-2-try-it-yourself}

It's time to have a go at coding yourself now!

Read through
\href{https://gla-my.sharepoint.com/:w:/g/personal/ashley_robertson_glasgow_ac_uk/ERHhgbPHADBJuJHk7e2yMtYBb1ddJMCTckAvoXpXB0UYFA?e=UFwJqE}{this
transcript}, taking brief notes individually the first time you read it.
What things did you pick up?

Identify codes for this transcript and
\href{https://padlet.com/ashleyrobertson3/coding-activity-for-rm2-odl-76lzunt0l2p5x0ma}{post
them on the padlet}. Ashley will then give class-wide feedback on the
codes to the class.

\section{Developing themes (Phases
3-6)}\label{developing-themes-phases-3-6}

Here, we'll focus on Phases 3 to 6 of the process. You can find
information on each of these phases in the Braun \& Clarke (2021) book:
Generating Initial Themes (p.78), Developing and Reviewing Themes
(p.97), Refining, Defining and Naming Themes (p.108) and Writing up your
Analysis (p.118). We also covered how to do Phase 3 in the Thematic
Analysis lectures in Semester 1. For these activities, you are welcome
to do them individually or in pairs/small groups - whichever you prefer!

\subsection{Phase 3: Generating initial themes
(p.78)}\label{phase-3-generating-initial-themes-p.78}

In this stage, you should sort the codes into potential themes. Here,
you will start to analyse your codes, and consider how they might share
an idea or concept.

It might be helpful to visualise this in some way, and see how the codes
fit together. You may have codes that don't seem to fit anywhere, others
will fit neatly into themes, others might be in sub-themes

You end this step with a collection of candidate themes (and potentially
sub-themes)

\begin{tcolorbox}[enhanced jigsaw, breakable, bottomtitle=1mm, leftrule=.75mm, opacityback=0, colbacktitle=quarto-callout-warning-color!10!white, opacitybacktitle=0.6, toprule=.15mm, colback=white, left=2mm, coltitle=black, bottomrule=.15mm, titlerule=0mm, title={Practical Steps}, toptitle=1mm, colframe=quarto-callout-warning-color-frame, arc=.35mm, rightrule=.15mm]

\begin{itemize}
\tightlist
\item
  Start to group your initial codes together
\item
  Can you see a pattern of shared meaning across different codes?
\item
  What are the similarities across codes?
\item
  End this stage by collating all the coded quotes relevant to each
  theme.
\item
  Themes can be descriptive or interpretative (see Braun \& Clarke,
  2013, Ch 11).
\item
  Using tables or mind maps at this stage can be helpful
\end{itemize}

\end{tcolorbox}

\subsection{Phase 4: Developing and Reviewing themes
(p.97)}\label{phase-4-developing-and-reviewing-themes-p.97}

In Phase 4, you check your themes to see whether they make sense a) when
you look at the coded extracts for that theme and b) when you look at
the whole dataset.

\emph{Checking the coded extracts}

Reviewing at the level of the coded extracts. Read all the collated
extracts for a theme\ldots do they form a coherent pattern? Do they tell
an important story about shared meaning? Are you sure they are themes
rather than topic summaries?

If they seem like themes and the coded extracts fit well underneath
them, then move on to checking the themes against the dataset.

If they don't, is it the case that a) the theme itself is problematic
and needs reworked or b) the theme is fine, but some of the data doesn't
fit?

\emph{Checking themes against the whole dataset}

Look across the entire dataset and consider your theme(s) -- do they
accurately reflect the meanings in the dataset? Do they tell you about
the important patterns that you observed in the dataset?

\begin{tcolorbox}[enhanced jigsaw, breakable, bottomtitle=1mm, leftrule=.75mm, opacityback=0, colbacktitle=quarto-callout-warning-color!10!white, opacitybacktitle=0.6, toprule=.15mm, colback=white, left=2mm, coltitle=black, bottomrule=.15mm, titlerule=0mm, title={Practical Steps}, toptitle=1mm, colframe=quarto-callout-warning-color-frame, arc=.35mm, rightrule=.15mm]

\begin{itemize}
\tightlist
\item
  This involves checking that the themes `work' in relation to both the
  coded extracts and the full data-set.
\item
  Reflect on whether the themes tell a convincing and compelling story
  about the data, and begin to define the nature of each individual
  theme, and the relationship between the themes.
\item
  It may be necessary to collapse two themes together or to split a
  theme into two or more themes, or to discard the candidate themes
  altogether and begin again the process of theme development.
\item
  It may be that your theme is actually a topic summary and needs
  reworked.
\item
  Not every participant will make specific reference to every theme -
  that is ok - but think about how each person is represented in the
  analysis.
\item
  It can be helpful at this stage to group the relevant quotes together
  that represent each given theme
\end{itemize}

\end{tcolorbox}

\subsection{Phase 5: Refining, Defining and Naming Themes
(p.108)}\label{phase-5-refining-defining-and-naming-themes-p.108}

\begin{itemize}
\tightlist
\item
  During this phase, you should define and refine the themes.
\item
  Define and refine = identifying what this theme is about, what aspect
  of the data is captured by each theme?
\item
  For each theme, you need to conduct and write a detailed analysis
\item
  In your report, you will report one theme with sub-themes or two
  themes (no subthemes) -- do they explain different things without
  overlapping too much?
\item
  Try to choose names for themes that are not one word or likely to
  indicate topic summaries (see the Semester 1 Thematic Analysis
  Lectures if you want to revisit this)
\end{itemize}

\begin{tcolorbox}[enhanced jigsaw, breakable, bottomtitle=1mm, leftrule=.75mm, opacityback=0, colbacktitle=quarto-callout-warning-color!10!white, opacitybacktitle=0.6, toprule=.15mm, colback=white, left=2mm, coltitle=black, bottomrule=.15mm, titlerule=0mm, title={Practical Steps}, toptitle=1mm, colframe=quarto-callout-warning-color-frame, arc=.35mm, rightrule=.15mm]

\begin{itemize}
\tightlist
\item
  This involves being able to explain each theme: `What story does this
  theme tell?' and `how does this theme fit into the overall story about
  the data?'.
\item
  Identify the `essence' of each theme and construct a concise, punchy
  and informative name. Usually, you would aim for two themes for the
  qualitative report
\end{itemize}

\end{tcolorbox}

\subsection{Phase 6: Writing up your Analysis
(p.118)}\label{phase-6-writing-up-your-analysis-p.118}

\begin{itemize}
\tightlist
\item
  The analysis should provide a concise, coherent, logical, non-
  repetitive and interesting account
\item
  The write up must provide sufficient evidence of the theme(s) (this is
  where the quotes come in!)
\item
  The extracts you choose should clearly illustrate the theme
\end{itemize}

\begin{tcolorbox}[enhanced jigsaw, breakable, bottomtitle=1mm, leftrule=.75mm, opacityback=0, colbacktitle=quarto-callout-warning-color!10!white, opacitybacktitle=0.6, toprule=.15mm, colback=white, left=2mm, coltitle=black, bottomrule=.15mm, titlerule=0mm, title={Practical Steps}, toptitle=1mm, colframe=quarto-callout-warning-color-frame, arc=.35mm, rightrule=.15mm]

\begin{itemize}
\tightlist
\item
  For the RM2 qualitative report, you should report either two themes
  (no subthemes) or one theme with 2-3 sub-themes. This is to allow you
  to balance depth and breadth.
\item
  Try to include at least one or two `longer' quotes for each theme,
  which allow you to explore them in depth. You can supplement these
  with shorter quotes too.
\item
  If you are doing a qualitative dissertation, you will have scope to go
  further with your analysis; it is a good idea to discuss your data
  with your supervisor throughout the analytic process.
\end{itemize}

\end{tcolorbox}

\subsection{Activities}\label{activities-2}

\subsubsection{Activity 1: Try it
yourself!}\label{activity-1-try-it-yourself}

We've provided the same transcript as for phases 1-2, but it's a little
longer (to give you more scope to develop themes) -
\href{https://gla-my.sharepoint.com/:w:/g/personal/ashley_robertson_glasgow_ac_uk/EZaA8YkYkyVGkwn0cEDFLCABWr5rREIlMAyIaAPHu3juTg?e=QgyPdd}{it
is linked here}. You can use the codes you developed last time (adding
to them for the additional parts of the transcript), or you can recode
the data again if you prefer - either option is fine.

\begin{itemize}
\tightlist
\item
  Once you have your codes, try to organise them into themes.
\item
  Group together the codes that seem similar in some way - can you think
  of an overarching `theme' that describes them?
\item
  Are there any codes that don't `fit' anywhere? (that's okay if so!)
\item
  Can you identify any quotes that you would use to evidence the themes?
  Do these definitely match well onto the theme?
\end{itemize}

\subsubsection{Activity 2: Bad analysis}\label{activity-2-bad-analysis}

This is an additional activity and only if you have time/want to delve a
bit further.
\href{https://gla-my.sharepoint.com/:w:/g/personal/ashley_robertson_glasgow_ac_uk/EXFrzL2zbfBCjDUeP6Py3MUB3p_WS15GyVYXur83wjuYUQ?e=fTMITU}{Have
a look at the analysis}. This is an example from Braun \& Clarke's book.
Why do you think it might be considered a poor example of thematic
analysis?

After you've finished, you can look at
\href{https://gla-my.sharepoint.com/:w:/g/personal/ashley_robertson_glasgow_ac_uk/Eded8-fYVWpOrj5vFY92gyQB0ieRX7BGUhZztrdA5auVhw?e=Pysxhy}{the
points that Braun \& Clarke identified}. Did you identify any additional
ones? Were there any you missed?

\section{Developing a narrative}\label{developing-a-narrative}

When doing your thematic analysis, it is important to develop a
narrative and go beyond paraphrasing/summarising your data.

Thematic Analysis requires an active approach from an analyst. It is
your role to make sense of the data and to present this in a coherent,
understandable and accessible way to the reader. You will take these
`patterns of meaning' and make sense of/interpret them, and convey these
to the reader.

One point to note: Remember that \textbf{themes do not emerge} - this
suggests that they are lying in wait for you to discover them, and that
the same themes would come up regardless of the analyst. This is not the
case - your experiences and interpretations will affect the themes you
end up with.

When you present your data in the analysis section, you will have a)
your list of themes and b) your quotes. It is your job to combine these
with analytical commentary, in order to tell the story of your data.
It's not enough to present the themes and then back these up with quotes
- your job is to add in a narrative that brings the themes and the
evidence (i.e.~quotes) together.

\textbf{Semantic and latent analysis}

Thematic Analysis allows you to explore \emph{semantic} or \emph{latent}
themes (or both) in your analysis:

\emph{Semantic:} captures explicitly-expressed meaning and stays close
to the language of the participants/overt meaning of data, realist. This
tends to be somewhat deductive. Semantic approaches still go beyond
summarising/paraphrasing.

\emph{Latent:} focuses on a deeper, more implicit meaning of the data.
What are the underlying meaning, ideas, concepts? This tends to be a bit
more inductive.

\textbf{Things to be aware of in developing a narrative}

\begin{enumerate}
\def\labelenumi{\arabic{enumi}.}
\tightlist
\item
  Remember to introduce your theme to the reader -- what is it about?
  Don't start a theme with a quote. If you have sub- themes, still
  introduce the overall theme briefly before you discuss your first
  sub-theme
\item
  Explain to the reader why you've chosen to include each quote. Why is
  it important enough to be included? What does it tell us about your
  data?
\item
  Focus on depth rather than breadth -- it is important that we get a
  real sense of your theme. In the RM2 report you only have 3000 words
  in total, and so will have to choose what you discuss
\item
  Show the reader who each quote comes from - we should be able to match
  this up with the demographics table/list of interviews in the Methods
  section (which then gives us some context). You can do this by using
  pseudonyms and linking this to the table by using the same names there
  too.
\item
  Think about where the opportunities are in your data to go beyond
  description. Is there anything you can infer? If so, remember to be
  clear with the reader that this is your interpretation. Your reader
  should be able why you have interpreted the data this way from the
  quotes you've chosen.
\end{enumerate}

\subsection{Activities}\label{activities-3}

\subsubsection{Activity 1: Semantic
vs.~latent}\label{activity-1-semantic-vs.-latent}

To explore differences between semantic and latent coding, as well as
analysing at an illustrative or analytical level, please read
\href{https://go.exlibris.link/wq9gqMKY}{Braun and Clarke (2013)}
p207-210 and also p252-254.

\subsubsection{Activity 2: Analysing
vs.~summarising}\label{activity-2-analysing-vs.-summarising}

Read pages 8-12 from
\href{https://study.sagepub.com/sites/default/files/anderson_clarke_reflective_commentary.pdf}{this
thematic analysis} on skin-picking. Here, the authors report three
themes and then have a reflection at the bottom. Try to identify where
the authors have gone beyond paraphrasing - discuss in your groups/pairs
if you can.

\subsubsection{Activity 3: Evaluating analysis and discussion
sections}\label{activity-3-evaluating-analysis-and-discussion-sections}

Do
\href{https://gla-my.sharepoint.com/:w:/g/personal/ashley_robertson_glasgow_ac_uk/EUiKwKgdqr1CsLeOxy5O42kBGNLxmbdeTqANFRu0caPcBA?e=BpEpBb}{this
exercise} on evaluating analysis sections of published papers.

\section{OPTIONAL: Research skills session - Preparing for data
analysis}\label{reflexivity_skills_session}

Ashley has recorded \textbf{optional} research skills sessions on
different aspects of qualitative research. These sessions take a little
bit deeper dive into some of the topics and may be particularly useful
if you are planning on doing a qualitative dissertation next year.

\href{https://echo360.org.uk/media/d4df06eb-4cf1-4ae0-bf8b-cf280d9da7e0/public}{Watch
the video here}

\subsection{Thinking about
reflexivity}\label{thinking-about-reflexivity}

We ask you to think about your own position in the research as part of
your RM2 report (there is a section in the methods!).

This can be quite tricky to do, but is very important for qualitative
research. Some suggestions to get you started:

\begin{itemize}
\tightlist
\item
  Think about the topic you are doing for your research and your
  position in it. Where would you place yourself in relation to it? What
  are your own views on the topic? Do you feel quite strongly about it
  or not?
\item
  Write as you reflect. Writing -- be it using pen or pencil on paper,
  typing on a keyboard, dictating to some voice capture device, or some
  other mode -- is an important tool for developing reflexive depth. It
  is good to do this at various stages of the research process.
\item
  Start with your own experiences, understandings and views, but then
  also try to interrogate those (i.e.~why do you have these particular
  understandings, views etc.?)
\end{itemize}

Some questions that you might ask yourself:

\begin{itemize}
\tightlist
\item
  What do I expect to come up in the data? How does this relate to what
  we found?
\item
  What broader experiences have I had that might influence how I'm
  thinking about this?
\item
  What are my values (i.e.~what is important to me)? What about my
  beliefs? How might these have infuenced my thinking?
\item
  How might all of these be connected to my identities and the
  communities I am a part of?
\end{itemize}

\subsection{Becoming familiar with the
data}\label{becoming-familiar-with-the-data}

It's important to become familiar with the data in the process of doing
the analysis. Read through your data and make the following notes:

\begin{itemize}
\tightlist
\item
  Things of potential interest
\item
  Ideas that you might want to explore further when you are coding
\item
  Your responses to the data (e.g.~how do you feel when you read it?)
\end{itemize}

Try to unpick any assumptions you have that underpin your initial
reactions and observations:

\begin{itemize}
\tightlist
\item
  What was familiar?
\item
  What was unfamiliar/surprising?
\item
  Why are you reacting to the data in that way?
\end{itemize}

\subsection{Activity 1: Read this extract and take familiarisation
notes}\label{activity-1-read-this-extract-and-take-familiarisation-notes}

Read this extract and take some familiarisation notes for yourself
(these are just for you and won't be shared). We will then go through
what Braun \& Clarke say about it.

Some context:

This activity is based on a data extract from an interview study
exploring LGBT students' experiences of university life. The participant
is a gay man -- a mature student, who spoke English as a second
language. The extract comes from early in the interview and in response
to a question about whether he is `out' as a gay man at university --
whether he is openly gay.

In this activity (5 mins):

\begin{itemize}
\tightlist
\item
  First, reflect on your assumptions and expectations about this topic,
  and your positioning in relation to this topic (are you an insider
  researcher? Outsider?).
\item
  Second, read the data extract and make some familiarisation notes.
\item
  Braun and Clarke's familiarisation notes follow -- reflect on
  similarities and differences in your and their analytic observations
  and `take' on the data.
\end{itemize}

\begin{info}
Andreas: \ldots I sometimes try to erm not conceal it that's not the
right word but erm let's say I'm in a in a seminar and somebody- a a man
says to me `oh look at her' (Int: mm) I'm not going `oh actually I'm
gay' (Int: mm {[}laughter{]}) I'll just go like `oh yeah' (Int: mmhm)
you know I won't fall into the other one and say `oh yeah' (Int: yep)
`she looks really brilliant' (Int: yep) but I sorta then and after them
you hate myself for it because I I don't know how this person would
react because that person might then either not talk to me anymore or
erm might sort of yeah (Int: yep) or next time we met not not sit next
to me or that sort of thing (Int: yep) so I think these this back to
this question are you out yes but I think wherever you go you always
have to start afresh (Int: yep) this sort of li- lifelong process of
being courageous in a way or not \ldots{}
\end{info}

\begin{tcolorbox}[enhanced jigsaw, breakable, bottomtitle=1mm, leftrule=.75mm, opacityback=0, colbacktitle=quarto-callout-caution-color!10!white, opacitybacktitle=0.6, toprule=.15mm, colback=white, left=2mm, coltitle=black, bottomrule=.15mm, titlerule=0mm, title={Show me Braun \& Clarke's familiarisation notes}, toptitle=1mm, colframe=quarto-callout-caution-color-frame, arc=.35mm, rightrule=.15mm]

\begin{itemize}
\tightlist
\item
  Andreas reports a common experience of presumed heterosexuality:
  coming out is not an obvious option.
\item
  Social norms dictate a certain response; the presumption of
  heterosexuality appears dilemmatic, and he colludes in the
  presumption, but minimally (to avoid social awkwardness).
\end{itemize}

Looking a bit more deeply, Braun and Clarke speculated that Andreas
values honesty and being true to yourself, but he recognises a
socio-political context in which that is constrained, and walks a
`tightrope' trying to balance his values and the expectations of the
context.

\end{tcolorbox}

\subsection{Activity 2: A code or a
theme?}\label{activity-2-a-code-or-a-theme}

In this activity, we'll
\href{https://gla-my.sharepoint.com/:w:/g/personal/ashley_robertson_glasgow_ac_uk/EchyMtm_oPxJjdf87kLvpyMB3nyxnSof5oEP-0A-PhnEGg?e=YKqeo8}{go
through a handout that asks you whether something is a code or a theme}.

For each label, decide whether it is a code or a theme.

\begin{itemize}
\tightlist
\item
  Keep in mind that although there is no absolute distinction between
  codes and themes, codes tend to have a single facet and capture one
  insight or observation about the data, whereas themes are ideally
  multifaceted and capture several insights and observations. Themes
  should also capture a pattern of shared meaning -- this can be shared
  meaning on the data surface, or shared implicit or underlying meaning.
\item
  Reflect on why you have decided that each code/theme is a code or
  theme. What features of the data, the label/name, and brief
  explanation of the scope and focus of each code/theme influenced your
  judgement?
\end{itemize}

\begin{tcolorbox}[enhanced jigsaw, breakable, bottomtitle=1mm, leftrule=.75mm, opacityback=0, colbacktitle=quarto-callout-caution-color!10!white, opacitybacktitle=0.6, toprule=.15mm, colback=white, left=2mm, coltitle=black, bottomrule=.15mm, titlerule=0mm, title={Show me Braun \& Clarke's solution}, toptitle=1mm, colframe=quarto-callout-caution-color-frame, arc=.35mm, rightrule=.15mm]

\begin{figure}[H]

{\centering \includegraphics[width=1\linewidth,height=\textheight,keepaspectratio]{images/themeorcode.png}

}

\caption{Braun and Clarke's suggestions}

\end{figure}%

Three themes were presented in the analysis: ``Two themes capture
contradictory ideas: that men's body hair is natural, and that men's
body hair is unpleasant. A third theme introduces the concept of
`excess' hair, which allowed sense-making of this contradiction,
mandating men's grooming of `excessive' hair.'' (Terry \& Braun, 2016,
page 14)

\end{tcolorbox}

\subsection{Some common problems with thematic
analysis}\label{some-common-problems-with-thematic-analysis}

\begin{itemize}
\tightlist
\item
  Thematic analysis that fails to address the research question.
\item
  Unconvincing or under-developed analysis: 1) Too many or too few
  themes 2) Too much overlap between themes or themes are unrelated 3)
  Themes are vague or not internally consistent or coherent
\item
  Little or no analytic work has been carried out (for example, using
  data collection questions as `themes').
\item
  Topic summaries being presented as themes
\item
  Mismatch between data and analytic claims (e.g.~going beyond what your
  participants say)
\item
  Too few or too many data extracts (and little or no analytic
  commentary).
\item
  Paraphrasing rather than analysing and interpreting data
\item
  Arguing with the data -- pointing out the `errors' in a participant's
  account.
\end{itemize}

\subsection{Strategies for ensuring quality in your
analysis}\label{strategies-for-ensuring-quality-in-your-analysis}

\textbf{Reflexive journalling}

\begin{itemize}
\tightlist
\item
  Reflect on your experiences throughout the whole research process (not
  just at the end!)
\item
  Think about your emotional responses to the data
\item
  Try to reflect on your assumptions that you already have about the
  topic
\end{itemize}

\textbf{Talk about your data with others}

\begin{itemize}
\tightlist
\item
  Discuss your analysis with other people - it helps to clarify your
  analytic insights
\item
  Do phases 1-2 in pairs if you want
\item
  Discuss your findings in groups
\item
  Come along to office hours
\item
  Present your analysis (even if its preliminary), this can help
  solidify where we are going with it
\end{itemize}

\textbf{Give yourself time}

\begin{itemize}
\tightlist
\item
  Try not to leave the analysis until the last minute
\item
  Give yourself time to step away and breathe
\item
  Perhaps work on other sections of your report at the same time, so
  that you have a mental break from analysis
\end{itemize}

\textbf{Make sure themes are themes (and be careful when naming them)}

Connelly and Peltzer (2016, p.55): ``When a researcher designated a
1-word theme, such as `collaboration', what does that mean in
relationship to the experiences of the informants as interpreted by the
researcher? Using only 1 word as a theme, there is no way of knowing,
for example if the experiences of collaboration were positive or
negative, or whether collaboration is important to the nurses. One-word
themes do not convey what the researcher found out about collaboration''

\chapter{Ethics, rigour and reflexivity}\label{ch9_ethics}

This chapter will help you to consider ethics and reflexivity in
qualitative research.

\section{Ethics}\label{ethics}

This section will talk you through the ethical considerations you need
to conduct qualitative research. As you remember from RM1, every time we
conduct research with human participants, we have to think carefully
about the potential impact of our research and we must adhere to the
\href{https://cms.bps.org.uk/sites/default/files/2022-06/BPS\%20Code\%20of\%20Human\%20Research\%20Ethics\%20(1).pdf}{BPS
Code of Human Research Ethics}. Some of the ethical issues are common to
both quantitative and qualitative research (for example, informed
consent and the right to withdraw), but due to the nature of qualitative
data, there are certain aspects we need to be particularly mindful
about.

\subsection{General ethical
considerations}\label{general-ethical-considerations}

These are some of the things we \textbf{MUST} take into account when
conducting research on human participants:

\begin{itemize}
\tightlist
\item
  Possible risks to participants (physical and/or psychological)
\item
  The elimination or minimization of risks
\item
  Justifying risks in terms of the scientific value of the results
\item
  Ethical approval prior to the study
\item
  The right to withdraw
\item
  Informed consent
\item
  Debriefing and aftercare of participants (especially with deception)
\item
  Anonymity and appropriate data management
\end{itemize}

\subsection{Activity 1: Identify potential ethical
issues}\label{activity-1-identify-potential-ethical-issues}

Before reading on, try to think what ethical issues each of these
elements of qualitative research might have:

\begin{itemize}
\tightlist
\item
  Research questions and interview/focus group questions
\item
  Data collection
\item
  Recording of interviews/focus group
\item
  Anonymity
\item
  Confidentiality
\end{itemize}

\textbf{1. Research questions and interview/FG questions}

\emph{This might look like:} In qualitative research, we often ask about
someone's lived experience, or their motivations, or a belief they have.
This is quite different to quantitative research where data comes in as
numbers and in that sense is less personal.

\emph{Why might this be an issue:} Qualitative research tends to be on
topics where we ask people to be vulnerable and open themselves up to us
as researchers. If we don't frame this correctly, our interview schedule
might contain questions that are inappropriate and uncomfortable for the
participants.

\textbf{2. Data collection}

\emph{This might look like:} Data collection is much less likely to be
anonymous, as typically in-person or online focus groups or interviews
are held. Samples are small, and participants may already be known to
the researcher.

\emph{Why might this be an issue:} There are possible issues of
confidentiality, anonymity (these will be discussed further down) and
bias. Our preconceptions might affect the way we frame the questions,
how we interact with participants in the focus group, and how we
interpret the data.

\textbf{3. Recording of focus groups/interviews}

\emph{This might look like:} The majority of the time, focus groups and
interviews are recorded (either audio + video or audio only). This means
that the `raw' data - as it were - is identifiable. In contrast, data
for quant studies are often collected anonymously.

\emph{Why might this be an issue:} Data should \textbf{always} be kept
safe and secure. However, with qualitative data, there is a real risk
that individual people could be identified from a recording of them.
This means that dealing with the recordings in a safe and secure way are
of utmost importance, and that participants cannot be identified from
transcripts.

\textbf{4. Anonymity }

\emph{This might look like:} In quantitative research, we tend to report
summary statistics from groups, and therefore it'd be quite unusual to
be able to identify a single participant from the data. With
qualitative, it's a bit different. We collect people's words and
experiences.

\emph{Why might this be an issue:} It is easier to identify someone from
the words they use, and the situations and circumstances they describe.

\textbf{5. Confidentiality}

\emph{This might look like:} In a focus group or interview, people might
be very open with us, and tell us things about their life, and the lives
of others.

\emph{Why might this be an issue:} We have a duty of care to keep
anything that comes out in a focus group/interview confidential (unless
they disclose a crime or someone is judged to be at risk). It's also
vital that the rest of the focus group members keep the information
confidential too.

\textbf{How would you mitigate these issues?}

\subsection{Activity 2: Ethics in published
studies}\label{activity-2-ethics-in-published-studies}

Look at the methods sections of these studies and try to identify how
they have addressed ethical issues

\begin{itemize}
\tightlist
\item
  \href{https://www.sciencedirect.com/science/article/pii/S2352721819301457\#s0010}{Scott
  et al., 2019}
\item
  \href{https://academic.oup.com/bjsw/article/43/3/596/1636960?login=true}{Sutton
  \& Stack, 2013}
\item
  \href{https://journals.sagepub.com/doi/pdf/10.1068/p7833}{Robertson \&
  Simmons, 2015}
\end{itemize}

\subsection{Activity 3: Information sheet and constent
form}\label{activity-3-information-sheet-and-constent-form}

Take a look at the (slightly modified) Information sheet and Consent
form from a qualitative survey study that had ethical approval from the
University. Can you identify where and how the principles listed in the
General ethical considerations and Activity 1 are discussed?

\subsubsection{Information sheet}\label{information-sheet}

How do perceived language barriers affect the teachers' experience
teaching international students in higher education

\textbf{Invitation} My name is XX, and I am a student at the University
of Glasgow. This study is conducted as a part of my MSc Psychological
Studies Dissertation project. You are invited to take part in this study
because you are currently teaching international students at a UK-based
university. Before you decide to take part, it is important for you to
understand why the research is being done and what it will involve.
Please take time to read the following information carefully and discuss
it with others if you wish. If you would like further details about the
study before deciding whether you want to participate, please get in
touch with us via email (XX@student.gla.ac.uk).

\textbf{Brief Summary} We are looking to learn more about higher
education teachers' perceptions and experiences of language barriers. We
are recruiting people who are currently teaching and/or supervising
international students in a UK-based university. To understand your
experience, opinions and thoughts on this topic, we are conducting the
study through an online qualitative survey. Some issues which we are
interested in learning more about include: 1) the impact of language
barriers on teaching and supporting students who speak English as a
second language; 2) how teaching staff perceive language barriers
affecting them and their students; 3) possible ways to bridge the
language barriers at different levels.

\textbf{What's involved if I agree to take part?} We want to support you
in deciding whether this research study is suited to you - if upon
reading this information sheet, you have more questions, please do not
hesitate to contact us. If you choose to take part, we will ask you to
provide informed consent, after which you will be asked to complete the
survey, which consists of two parts: a brief demographic questionnaire,
and six open-ended main questions about your experiences of language
barriers in your teaching. The survey should take no longer than 40
minutes to complete, and you may take a break at any point if you wish.
Because the data will be analysed using qualitative methods, we ask you
to provide a comprehensive account of your experiences, but only to
provide as much or as little detail as you feel comfortable sharing.
Anonymised (see below) quotes from your responses may be used in the
dissertation and other future dissemination to illustrate the findings.
If you choose to withdraw from the study at any point, you may do so by
simply closing your browser.

\textbf{Confidentiality: Ways to protect your identity and data} We do
not want to reveal any personal information about yourself or your
colleagues as part of this research. Your responses are recorded
completely anonymously through Microsoft Forms and your name will not be
recorded anywhere. Prior to the data analysis, all responses will be
screened and further anonymised in case you have provided information
which may be identifiable (such as specific course names or
institutions). We will take the following steps to make sure that your
personal information and particular details of your experiences are not
shared.

\begin{itemize}
\tightlist
\item
  Your data will be assigned an anonymous participant ID automatically
  and password protected on a secure server
\item
  Changing any names, places or organisations you write in your answers.
\item
  All the research data will be securely stored to comply with the
  General Data Protection Regulation (2018) and only made accessible to
  future researchers in a form which protects individual, personal
  information.
\end{itemize}

\textbf{What are the possible disadvantages and risks of taking part? }
The materials do not involve anything that would affect you in a
negative way -- as explained above, you may share as much or as little
detail about your experiences as you feel comfortable doing. You are
also free to withdraw from the research by closing your browser at any
point without giving any reason why.

\textbf{Who is organising this research?} This research is being
conducted as part of an MSc Dissertation. The research is supervised by
Dr Wilhelmiina Toivo, Wilhelmiina.toivo@glasgow.ac.uk and the research
has been granted ethical approval from the University of Glasgow. If you
would like further details about the study before deciding whether you
want to participate, please get in touch with us via email. Many thanks,
XX,

\subsubsection{Consent form}\label{consent-form}

\textbf{Title of Project:} How do perceived language barriers affect the
teachers' experience teaching international students in higher
education?

\textbf{Researcher:} XX, (Supervisor Wilhelmiina Toivo)

If you agree to participate in this study, then please read the
following statements and indicate your consent by ticking the boxes for
each statement

\begin{itemize}
\tightlist
\item
  I have read the Information Form for participants and so understand
  the procedures and have been informed about what to expect.
\item
  I agree to participate in this study, which aims to investigate how
  perceived language barriers affect university teachers' experiences
  teaching and interacting with international students who speak English
  as their foreign language.
\item
  I understand that my participation in this study is voluntary, and
  that I can withdraw from the study, at any time and for any reason, by
  closing the browser.
\item
  I understand that my participation in this project is for the purposes
  of research and is in no way an evaluation of me as an individual.
\item
  I understand that any information given in the investigation will be
  made and kept anonymous and will remain confidential and no
  information that identifies me will be made publicly available.
\item
  I understand that anonymous quotes from my responses may be used in
  academic publications/talks to illustrate the findings.
\item
  I understand that I can contact the researcher for this project by
  e-mail to receive more information and/or a summary of the results.
\end{itemize}

\section{Reflexivity}\label{reflexivity}

Please note: Some information below is included in the video for
\hyperref[reflexivity_skills_session]{Research Skills for data analysis}
as well.

We ask you to think about your own position in the research as part of
your report (there is a section in the methods, and you need to include
a
\href{https://gla-my.sharepoint.com/:w:/g/personal/ashley_robertson_glasgow_ac_uk/IQC1PIxhBlndR74S_nfeKHNlAW5uQmGOpoRXCpLHC2fI_fI?e=aunShJ}{reflexivity
diary} as an appendix to your report!).

This can be quite tricky to do, but is very important for qualitative
research. Some suggestions to get you started:

\begin{itemize}
\tightlist
\item
  Think about the topic you are doing for your research and your
  position in it. Where would you place yourself in relation to it?
\item
  What are your own views on the topic? Do you feel quite strongly about
  it or not?
\end{itemize}

Write as you reflect. Writing -- be it using pen or pencil on paper,
typing on a keyboard, dictating to some voice capture device, or some
other mode -- is an important tool for developing reflexive depth. It is
good to do this at various stages of the research process.

Start with your own experiences, understandings and views, but then also
try to interrogate those (i.e.~why do you have these particular
understandings, views etc.?)

Some questions that you might ask yourself:

\begin{itemize}
\tightlist
\item
  What do I expect to come up in the data? How does this relate to what
  we found?
\item
  What broader experiences have I had that might influence how I'm
  thinking about this?
\item
  What are my values (i.e.~what is important to me)?
\item
  What about my beliefs? How might these have influenced my thinking?
\end{itemize}

Here we have some data from Braun \& Clarke - there are different
datasets to choose from. If you want to use the ones from Braun \&
Clarke, you can either choose the
\href{https://gla-my.sharepoint.com/:w:/g/personal/ashley_robertson_glasgow_ac_uk/EWq1HW7Xm_BKnobsU_GRpcIBlSvnQ4k6XhFqBkszClAwWg?e=q4lc97}{`childfree'}
dataset or the
\href{https://gla-my.sharepoint.com/:w:/g/personal/ashley_robertson_glasgow_ac_uk/EdfitbKwhodKo8Kmj9szNgwBNmt313A8dYknAiehHdy2lg?e=0vWOnu}{`healthy
eating'} dataset. However, we know that some topics might be difficult
for people, so please do choose a topic of your own preference if you
would rather - any qualitative dataset (e.g.from UK Data Service) or
news article would be fine.

\subsection{Activity 1}\label{activity-1}

Spend about fifteen minutes reflecting on the topic and where you would
place yourself in relation to it, both personally and more general
views. Write as you reflect. Writing -- be it using pen or pencil on
paper, typing on a keyboard, dictating to some voice capture device, or
some other mode -- is an important tool for developing reflexive depth.
Start with your own experiences, understandings and views, but then also
try to interrogate those.

Ask yourself questions like:

\begin{itemize}
\tightlist
\item
  What assumptions am I basing my thinking on?
\item
  What broader experiences and values might inform my thinking on this?
\item
  How might these be connected to myself as a person, my various
  identities, my personal experience in relation to the topic, and the
  communities I am part of?
\end{itemize}

\subsection{Activity 2}\label{activity-2}

Discuss your experiences of reflecting with your group. Did you find it
easier than expected? Was it harder? Was there anything that surprised
you about the process? Please only share what you feel comfortable with.

\subsection{Activity 3 (optional)}\label{activity-3-optional}

Ashley has some additional reflexivity probes and tips in her
\hyperref[reflexivity_skills_session]{research skills session}.

\chapter{Group project}\label{ch10_group_project}

\section{What is the group project?}\label{what-is-the-group-project}

In the group project, you will engage in three tasks that are designed
to help you build your practical research skills in qualitative
research. You will compile your answers to the tasks into a portfolio
that you submit to Moodle as a group. Please refer to Moodle (Semester 1
Week 14) for a video overview of the project, and make sure that you
have engaged with the \hyperref[ch3_AIS_portfolio]{Assessment
Information section of the group portfolio} before starting any work on
it so you know what to expect.

To help you balance the different assignments in RM2, we have compiled
\href{https://gla-my.sharepoint.com/:w:/g/personal/ashley_robertson_glasgow_ac_uk/EQM2Fzsi0FJAlDKTJXgi7-8Brm7ITIBw5ta3MaXSiDZfrg?e=FMLoM4}{a
suggested timeline} to help you is a suggested timeline for the
qualitative portfolio and report

\section{Getting to know each other in
groups}\label{getting-to-know-each-other-in-groups}

This section is designed for you to engage with your group early, get to
know each other, and develop some common guidelines for working together
before you embark on the group project. If you haven't done so yet,
please get in touch with your group now to complete the following
activities. We have also provided a checklist of things to complete as a
group before the start of semester 2.

\subsection{Getting organised in groups
checklist:}\label{getting-organised-in-groups-checklist}

\begin{itemize}
\tightlist
\item
  Get in touch with group and make a group chat on Teams
\item
  Introduce yourselves
\item
  Post group introduction on RM2 channel
\item
  Set up a regular check-in time and/or meeting with group
\item
  Discuss group agreement and decide together on items to include
\item
  Submit group agreement to Moodle (individually)
\item
  Read effective group work document and discuss your strengths as a
  group
\item
  Download the group work template and assign tasks
\end{itemize}

\subsection{Activity 1: Create a group chat and set up an initial
meeting}\label{activity-1-create-a-group-chat-and-set-up-an-initial-meeting}

Get in touch with your group and create a Teams chat. You \textbf{MUST}
use Teams for all communication about the group project. Do not use any
other platforms (such as Whatsapp, WeChat, Facebook etc.) as these can
lead to people being excluded. Set up a meeting in which you will
discuss Activity 3 from this chapter. After the initial meeting, set up
a regular check-in time. This doesn't necessarily need to be a meeting
every time, you can also do it over text. You will need to coordinate
between time zones/work calendars/other responsibilities

\emph{DO:} Communicate regularly and be clear about where and how you
are going to communicate about the group project. Be proactive if you
are having a busy time and let your group know.

\emph{DON'T:} Disappear and stop talking to your group if you're having
a hard time (if something happens and you don't feel comfortable sharing
with your group, you can let the course leads know). Don't agree to
meeting times that don't suit you.

\subsection{Activity 2: Introductions}\label{activity-2-introductions}

Introduce yourself to your group members on your Teams chat. In your
introduction, tell your group where you are based, which parts of the
degree you have enjoyed the most up to date and what is your favourite R
function. As a group, find one \emph{non-psychology/university related
thing that you all have in common} and share that on the RM2 Teams
channel in the thread. Bonus points for the most random common thing!

\subsection{Activity 3: Group
agreements}\label{activity-3-group-agreements}

We appreciate that group work can be challenging (and even more so if
you are spread all over the World!) and we want to support you to have
an enjoyable and educational group work experience. To that end, we
would like each group to create a group agreement that all members sign.
We also ask that you complete a task allocation and upload this with
your agreement (they are in the same template). The purpose of this is
for the group to come together and establish common ground and some
rules that each of the group members is committed to. It also means that
the group begins working with clearly allocated roles and tasks. This is
not to mean that you cannot change tasks later, but it is helpful to set
expectations from the outset.

You can download
\href{https://gla-my.sharepoint.com/:w:/g/personal/ashley_robertson_glasgow_ac_uk/EcRHw5yd3v1JmFwm1drRu-MB33Im0__d6jRPqKbdnjjDzg?e=s7wEoU}{the
group agreement template}.

Then look at the list of potential items to include below, and as a
group, discuss which items you want to include in addition to the
standard ones provided, and if there are any you would like to edit or
add outside of this list. You will be able to find a sample list of
tasks to include for the task allocation in Appendix 1 within the
template.

Work on a collaborative document on OneDrive to create your group work
agreement and task allocation, and once the group is happy with their
agreement and initial tasks, each group member should submit it to
Moodle individually (on the Formative Activities tile).
\textbf{\emph{Please submit your group agreement by the end of week 14
of semester 1.}} If you haven't submitted your agreements and task
allocation, we will be in touch to check up on your group.

\textbf{Group agreement items}

Here is a list of items you may want to discuss with your group and add
to the group agreement:

\begin{itemize}
\tightlist
\item
  We all agree to establish a common timeline for the project and
  discuss any deviations from the timeline together as a group
\item
  We all agree to communicate with each other openly and honestly about
  the project and get in touch with the course leads in the event of an
  unsolvable group conflict.
\item
  We all agree to respect each other's individual ways of working and
  discuss these openly as a group.
\item
  We all agree to respect diversity in our group. This includes cultural
  differences, neurodiversity, different work/life situations and may
  manifest in different ways of working/approaching tasks.
\item
  We agree to discuss our strengths as a group and divide tasks
  accordingly.
\item
  I agree to stay in touch with the group and communicate as established
  by the group. If I am struggling and I do not feel comfortable sharing
  that with the group, I will get in touch with the course/programme
  leads for support
\item
  We all agree to meet on a regular basis as established by the group.
  If I cannot make a meeting I agree to communicate this openly with my
  group and get in touch to catch up about things I have missed.
\item
  We agree to plan out the individual contribution of each team member
  in advance before group work starting and openly communicate any
  changes that may arise during the project.
\item
  We agree to establish the best method of communication for our group
  that's inclusive and takes into account everybody's individual needs.
\end{itemize}

\subsection{Activity 4: Read the Effective Group work
document}\label{activity-4-read-the-effective-group-work-document}

\href{https://gla-my.sharepoint.com/:w:/g/personal/ashley_robertson_glasgow_ac_uk/EZCuPb2m-MFMtAIGTAb2eM4B_ateTxqyR5X6YLceaL2Qcg?e=XeHk9h}{This
document can be used} as a checklist to do the group work for your
qualitative portfolio. Give it a read and then discuss it with your
group once you have your initial meeting.

\subsection{Activity 5: Assign tasks}\label{activity-5-assign-tasks}

Work back from the deadline to establish tasks and when they need to be
done. Discuss your strengths and what each person in the group would
like to do.

\section{Qualitative portfolio guidelines and
materials}\label{qualitative-portfolio-guidelines-and-materials}

You will complete three tasks: Qualitative questions (Task 1a and 1b),
Ethics (Task 2) and Terrible focus group (Task 3). Below, you will find
information and the required materials to complete each task. Please use
\href{https://gla-my.sharepoint.com/:w:/g/personal/ashley_robertson_glasgow_ac_uk/EeFbaG-CE0pEnWBvsAwGIXUBIR8wAoFdJYYhklHoUKkEjg?e=DTA61b}{this
template} to complete your submission.

\subsection{Task 1a: Qualitative
questions}\label{task-1a-qualitative-questions}

Below, you will find a brief overview of a study topic. In this task, we
would like you to imagine that you are collecting data for this
qualitative study. A variety of different data collection methods are
possible (e.g.~interview, focus group, open-ended survey questions).
Thinking about the topic and research question presented below, please
develop \textbf{SIX} questions that you would ask your participants to
gather data on the topic of interest. Your questions should be
open-ended, clear, concise and appropriate, and the order that they are
presented in should make sense. Your aim is to develop questions that
would be likely to generate open discussion and rich responses from
participants.

\textbf{Please note: This section (1a) does not contribute towards the
1000 word count for the assessment.}

\begin{tcolorbox}[enhanced jigsaw, breakable, bottomtitle=1mm, leftrule=.75mm, opacityback=0, colbacktitle=quarto-callout-note-color!10!white, opacitybacktitle=0.6, toprule=.15mm, colback=white, left=2mm, coltitle=black, bottomrule=.15mm, titlerule=0mm, title={Study overview}, toptitle=1mm, colframe=quarto-callout-note-color-frame, arc=.35mm, rightrule=.15mm]

Universities provide support for students, and this support can come in
many forms, from free counselling services to well-being advisors
\href{https://link.springer.com/article/10.1007/s10447-008-9062-0}{(Raunic
\& Xenos, 2008)}.

In one study, it was reported that only 3\% of students used their
university counselling service, and only 5\% of those who were
characterised as vulnerable based on questionnaire scores, used the
services \href{https://doi.org/10.1080/03069880600942624}{(Cooke et al.,
2006)}. This suggests that the majority of students did not make use of
services available.

There are many reasons why students may not choose to seek help,
including being unaware of services, low socio-economic status, being
from different cultural background, embarrassment about having to seek
help, and scepticism about the services
(\href{https://journals.lww.com/lww-medicalcare/fulltext/2007/07000/help_seeking_and_access_to_mental_health_care_in_a.3.aspx}{Eisenberg
et al., 2007};
\href{https://bmcpsychiatry.biomedcentral.com/articles/10.1186/1471-244X-10-113}{Gulliver
et al., 2010)}.

Another issue around seeking support is that many who do chose to get
help, have to wait long periods for appointments
\href{https://go.exlibris.link/htR5xvWt}{(Mowbray et al., 2006)}. In a
recent article reported in the Independent, a waiting time of 4 months
was cited for some universities. These delays to support, can lead to
emergence of mental health crises situations. Unfortunately, in this the
University of Glasgow had one of the longest wait times for counselling
\href{https://www.independent.co.uk/news/uk/politics/students-mental-health-support-waiting-times-counselling-university-care-diagnosis-treatment-liberal-democrats-norman-lamb-a8124111.html}{(Buchan,
2018)}.

\emph{Research question: What are student's views on how universities
can support students and their mental health needs?}

\end{tcolorbox}

\subsection{Task 1b: Qualitative questions
reflection}\label{task-1b-qualitative-questions-reflection}

For task 1B, please reflect on your experience of generating the
questions in task 1A. Do you think they would achieve what you want to
achieve? Why or why not? Was it easier or harder than you expected?
Please link to the evidence base in your answer (suggested word count:
250 words)

\subsection{Task 2: Ethics}\label{task-2-ethics}

With the Open Science movement, the use of Open, secondary datasets has
become a more popular way of conducting research. You've learnt about
the importance of Open data in terms of reproducibility and
replicability for quantitative research in RM1 - however, the benefits
and challenges of Open data are quite different in qualitative research
and there are a number of ethical issues associated with sharing
qualitative research data openly. For task 2 of your portfolio, we would
like you to \textbf{identify one ethical issue} associated with the
sharing and/or use of open qualitative data, and discuss why it is an
issue, and how it could be mitigated. In your response, you should cite
relevant literature on the use of open qualitative data, as well as the
\href{https://cms.bps.org.uk/sites/default/files/2022-06/BPS\%20Code\%20of\%20Human\%20Research\%20Ethics\%20(1).pdf}{BPS
code of Ethics} (suggested word count 250 words)

\subsection{Task 3: Focus group gone
wrong}\label{task-3-focus-group-gone-wrong}

Here is a short video from a fictional focus group that is conducted
very badly (\emph{please note this video is created only for the
purposes of this assessment, is not from a real focus group, and should
not be shared outside of this course!}).

Your task is to \textbf{identify two or three issues with how the focus
group is run} and discuss why they are problematic and outline possible
ways to improve them. In your response, you should cite relevant
literature to support your ideas.

\href{https://echo360.org.uk/media/f51c8477-94cd-4a0e-88a8-0a133dd8d21e/public}{Access
the Focus Group gone wrong recording}

\subsection{Where to get guidance for the portfolio
tasks?}\label{where-to-get-guidance-for-the-portfolio-tasks}

You should make sure you are up to date on Ashley's lectures on
qualitative research. Further, some of the chapters in this book have
activities and information which will be helpful in completing the
portfolio:

\begin{itemize}
\tightlist
\item
  \hyperref[ch6_qual_questions]{Chapter on Qualitative questions}
\item
  \hyperref[ch7_data_collection]{Chapter on Data collection}
\item
  \hyperref[ch9_ethics]{Chapter on Ethics, Rigour and Reflexivity}
\end{itemize}

\chapter{Academic writing}\label{ch11_writing}

This chapter has some activities to help with your academic writing.

\section{Evaluation in introductions}\label{eval_intro}

The activities in this section are designed to help you improve your
evaluation in the introduction of your RM2 report (and your
dissertation). To re-cap from RM1, the purpose of your introduction is
to provide a rationale for your study and answer the \textbf{WHY}
question; why is your study worth doing? You want to sell your idea to
the reader, using existing literature to show how you are building on
the evidence that is there.

Although `identifying the gap' is something that is often key in
quantitative research, building a rationale is not exactly the same in
qualitative research. This is because qualitative research tends to be
very much situated in the context of the particular study (i.e.~why is
this particular research topic or question important in this
context/with this group). Regardless of whether you include a gap or not
(you still can), you should convey to your reader \textbf{why} your
study matters (or why a gap you've identified is important to explore).

\textbf{Things to do and things to avoid}

\begin{enumerate}
\def\labelenumi{\arabic{enumi}.}
\tightlist
\item
  Avoid doing one study per paragraph -- instead, build your arguments
  using multiple sources
\item
  Avoid describing previous studies in a lot of detail -- instead, show
  how your study builds on them
\item
  Avoid critiques that your study cannot answer - instead, focus on
  aspects that your study will address
\item
  Just because something is never done before, doesn't make it worth
  doing; show with evidence why we should do the study.
\end{enumerate}

Here are paragraphs from an imaginary introduction to a qualitative
study investigating postgraduate students' confidence in their zombie
apocalypse survival.

For each, you can click on them to have a look in depth.

\begin{tcolorbox}[enhanced jigsaw, breakable, bottomtitle=1mm, leftrule=.75mm, opacityback=0, colbacktitle=quarto-callout-caution-color!10!white, opacitybacktitle=0.6, toprule=.15mm, colback=white, left=2mm, coltitle=black, bottomrule=.15mm, titlerule=0mm, title={Show me Example 1: Using one source, writing descriptively}, toptitle=1mm, colframe=quarto-callout-caution-color-frame, arc=.35mm, rightrule=.15mm]

\emph{A study by Lane and Karin (2011) looked at zombie apocalypse
survival confidence levels in undergraduate students. The sample
consisted of 25 participants (20 male and 5 female), recruited through
convenience sampling. The study employed involved a survey comprising
Likert-scale questions designed to measure participants' self-assessed
confidence in various survival skills. Additionally, participants were
asked about their exposure to zombie-related media. The data were
subjected to t-tests and regression analyses, and the authors' found
that the biggest predictor of zombie apocalypse confidence was exposure
to zombie-themed video games.}

This paragraph uses just one source, and describes the study in a lot of
detail. Instead, it would be better to synthesise evidence from multiple
studies and focus on what the findings mean, rather than being
descriptive about the studies. Something like:

\emph{A study by Lane and Karin (2011) found that the biggest predictor
of zombie apocalypse survival confidence levels in undergraduate
students was exposure to zombie-themed video games. Miller and Williams
(2023) had similar findings, but exposure to zombie films was also a
significant predictor of confidence. This suggests that simulation
through media is an effective way to increase self-efficacy towards an
unknown threat, which has also been found with alien attack research
(Ripley, 2017).}

\end{tcolorbox}

\begin{tcolorbox}[enhanced jigsaw, breakable, bottomtitle=1mm, leftrule=.75mm, opacityback=0, colbacktitle=quarto-callout-caution-color!10!white, opacitybacktitle=0.6, toprule=.15mm, colback=white, left=2mm, coltitle=black, bottomrule=.15mm, titlerule=0mm, title={Show me Example 2: Making criticisms that are outwith the scope of your
study}, toptitle=1mm, colframe=quarto-callout-caution-color-frame, arc=.35mm, rightrule=.15mm]

\emph{It should be noted that the majority of studies on zombie
apocalypse survival confidence have looked at samples that consist of
mostly men (Lane \& Karin, 2011 but see also Grimes \& Dixon, 2019), and
that previous samples have been small.}

This paragraph makes a criticism that the study cannot address - issues
with sample size are not something that we can address with qualitative
designs.

Instead, to create a stronger rationale, it would be better to focus
your critical evaluation on aspects which your study actually builds on:

\emph{While evidence on the behavioural predictors of zombie apocalypse
confidence is robust, the majority of research (see Anderson, 2014) has
focused on the exposure time to media and survival skills training.
Hence, we lack understanding of how participants perceive their
confidence beyond simulation experiences.}

\end{tcolorbox}

\begin{tcolorbox}[enhanced jigsaw, breakable, bottomtitle=1mm, leftrule=.75mm, opacityback=0, colbacktitle=quarto-callout-caution-color!10!white, opacitybacktitle=0.6, toprule=.15mm, colback=white, left=2mm, coltitle=black, bottomrule=.15mm, titlerule=0mm, title={Show me Example 3: Why is your study worth doing?}, toptitle=1mm, colframe=quarto-callout-caution-color-frame, arc=.35mm, rightrule=.15mm]

\emph{To date, there has been only one qualitative study on the topic --
Seok (2013) collected data through semi-structured interviews, which
were analysed using thematic analysis and two themes, ``Survival mode''
and ``Confidence through skills'' were identified. This study looked at
undergraduate students; there is no previous literature on postgraduate
students' perceptions of their apocalyptic confidence. Therefore, this
qualitative study investigates the perceived confidence of postgraduate
students in their abilities to survive a zombie apocalypse.}

Just because it's never been looked at doesn't automatically mean it
should be - show with evidence why and how you would expect your new
sample/method/something else about the study to differ from previous
research and why it's worth exploring:

\emph{Previous research has looked at undergraduate students; however,
it has been suggested that postgraduate students would be the group most
at risk during a zombie apocalypse (Williams, 2021), but also that
apocalyptic self-efficacy decreases with age (Anderson, 2004). This
suggests that postgraduate students are the most mentally and physically
vulnerable group during a zombie apocalypse, and hence it is important
to understand their perceptions of apocalypse confidence, and the lack
of thereof.}

\end{tcolorbox}

\subsection{Activity 1: RM1 feedback}\label{activity-1-rm1-feedback}

Go back to your feedback from RM1 and see what your marker has said
about the evaluation ILO and comments for your introduction. List two
things that you did well, and two things you want to improve for RM2. If
you would like to discuss the feedback and possible improvement points
further, you can pop to Wil's or Ashley's office hours.

Quite often what we see is students engaging with a lot of literature in
the intro (which is good!), but not really building a rationale or
explaining with evidence why the study is worth doing. We want you to
move beyond describing the research to critically evaluating it to show
how your study expands on existing knowledge.

\subsection{Activity 2: Top tips and rationales in published
papers}\label{activity-2-top-tips-and-rationales-in-published-papers}

\textbf{Tip 1: Show how your study builds on existing literature} When
you are building your rationale, try to use multiple studies to show
your reader what the current context of the literature is, and how your
study is expanding that knowledge. What new is your study bringing? How
is it expanding on existing studies?

If you look at these paragraphs from
\href{https://www.sciencedirect.com/science/article/pii/S1750946717301277}{Robertson
et al.~(2018)}, they are building the rationale for their study by
showing how it expands on existing literature:

\begin{tcolorbox}[enhanced jigsaw, breakable, bottomtitle=1mm, leftrule=.75mm, opacityback=0, colbacktitle=quarto-callout-note-color!10!white, opacitybacktitle=0.6, toprule=.15mm, colback=white, left=2mm, coltitle=black, bottomrule=.15mm, titlerule=0mm, title={Building a rationale}, toptitle=1mm, colframe=quarto-callout-note-color-frame, arc=.35mm, rightrule=.15mm]

``Kerns and Kendall (2012) published an updated review of the prevalence
of anxiety in individuals with ASD, and all studies fell within White et
al.~(2009)'s original range for anxiety symptoms (11--84\%). It should
be noted that, of the 24 studies included in the review, only two
reported data for individuals older than 19 years of age (Bakken et al.,
2010, Hofvander et al., 2009), which illustrates the need to also focus
on adults with ASD in future research (Howlin, 2000).''

``Secondly, as far as we are aware, there have been no qualitative
studies investigating anxiety in autism using semi-structured
interviews. Although focus groups and interviews are both good data
collection methods for qualitative data, there is evidence that
individual interviews tend to elicit a broader range of discussion
points (Guest, Namey, Taylor, Eley, \& McKenna, 2017). Again, this
further extends the previous investigation conducted by Trembath et
al.~(2012).''

\end{tcolorbox}

\textbf{Tip 2: Focus your criticism}

When you are critically evaluating studies in your introduction, make
your criticism focused on things that your study is addressing. This
helps you build the rationale. For example, for the RM2 report, don't
focus on the sample size and generalisability of existing studies in
your introduction, because these are not things that your study will
actually be able to address. Instead, focus on gaps/problems your study
can address.

If you look at this paragraph from
\href{https://journals.plos.org/plosone/article?id=10.1371/journal.pone.0210450}{Toivo
\& Scheepers (2019)}, they are identifying a gap/problem in existing
literature and showing how the new study is addressing this:

\begin{tcolorbox}[enhanced jigsaw, breakable, bottomtitle=1mm, leftrule=.75mm, opacityback=0, colbacktitle=quarto-callout-note-color!10!white, opacitybacktitle=0.6, toprule=.15mm, colback=white, left=2mm, coltitle=black, bottomrule=.15mm, titlerule=0mm, title={Focusing cricitism}, toptitle=1mm, colframe=quarto-callout-note-color-frame, arc=.35mm, rightrule=.15mm]

``A potential issue in the above SCR studies (as well as in some of the
cognitive studies discussed earlier) is that the stimuli were not very
tightly controlled in terms of lexical variables (length, frequency,
abstractness, etc.), or in terms of syntactic/pragmatic complexity when
multi-word phrases were used. This, again, introduces a number of
potential confounds (both within and across languages), making it
difficult to separate effects of emotional resonance from those related
to cognitive effort in word processing. In the present study, we aimed
at eliminating such potential confounds by matching our stimuli on a
number of lexical variables--as was done in the ERP studies discussed
above {[}22, 23{]}--and by avoiding the use of translation
equivalents.''

\end{tcolorbox}

\textbf{Tip 3: Draw arguments from related topics. }

Sometimes if your topic has not been studied much, it might be difficult
to create an evidence-based rationale for your research question. In
this case, you can draw logical parallels from related topics/domains.
Remember to explain to the reader why you would expect the findings to
generalise across the related topics.

If you look at this paragraph from
\href{https://journals.plos.org/plosone/article?id=10.1371/journal.pone.0204991}{Mahrholz
et al.~(2018)}, it acknowledges that evidence on what they are looking
at is scarce in voice perception, but the same effect has been found in
face perception:

\begin{tcolorbox}[enhanced jigsaw, breakable, bottomtitle=1mm, leftrule=.75mm, opacityback=0, colbacktitle=quarto-callout-note-color!10!white, opacitybacktitle=0.6, toprule=.15mm, colback=white, left=2mm, coltitle=black, bottomrule=.15mm, titlerule=0mm, title={Drawing arguments from other topics}, toptitle=1mm, colframe=quarto-callout-note-color-frame, arc=.35mm, rightrule=.15mm]

``Similarly, looking at reliability of personality traits across
presentation durations, Willis and Todorov {[}40{]} found that ratings
of trustworthiness, competence, likeability, aggressiveness, and
attractiveness for faces, showed moderate to strong positive
correlations after 100 ms, 500 ms, and 1000 ms, when compared to ratings
made without time constraints. Only participants' confidence in their
own judgements increased as a function of duration. Likewise, again
using photographs of faces, Bar et al.~{[}33{]} reported medium positive
correlations between ratings at 39 ms and 1700 ms. The authors indicated
that the lower threshold was sufficient for reliable assessments of
threat but not intelligence, supporting the theory that rapid first
impressions serve as a mean of self-preservation and help determine
appropriate approach-avoidance behaviour {[}1, 33, 38{]}. The idea being
that it should not require much information to decide whether a stranger
is friend or foe. Finally, Todorov and colleagues {[}39{]} obtained a
similar finding, again for faces, showing 33 ms of exposure to be
sufficient to distinguish between trustworthy- and untrustworthy-looking
stimuli. Whilst correlations with control ratings improved between 33 ms
and 100 ms, increased exposure duration did not significantly increase
the correlations.

In voice research, though there are limited studies that consider the
reliability of personality judgements across varying lengths of stimulus
types, similar findings have been shown as in face research.''

\end{tcolorbox}

\subsection{Activity 3: Active reading}\label{activity-3-active-reading}

When you are reading journal articles, make sure to take notes and think
of how each article can be used to build your rationale in the
introduction.
\href{https://gla-my.sharepoint.com/:w:/g/personal/wilhelmiina_toivo_glasgow_ac_uk/EXPNU5j4pQhKvRMYVx8dBvcBs85f4Z8q1V21wSokGOIbCA?e=eg7zd0}{Here
is a resource you can use to help structure your notes} about the
studies you are going to use for your introduction.

\section{Evaluation in your analysis section}\label{eval_analysis}

\subsection{Activity 1}\label{activity-1-1}

\href{https://gla-my.sharepoint.com/:w:/g/personal/ashley_robertson_glasgow_ac_uk/EUiKwKgdqr1CsLeOxy5O42kBGNLxmbdeTqANFRu0caPcBA?e=BpEpBb}{Complete
this activity} to improve your evaluation for your thematic analysis,
(as well as the discussion section). You may have already completed this
as a part of your Data analysis chapter, but if you haven't please do it
now. We ask you to compare two published papers against some of the
qualitative report ILOs.

\subsection{Activity 2: Going deeper in your
analysis}\label{activity-2-going-deeper-in-your-analysis}

One of the characteristics of a strong analysis section is that the
writer demonstrates an ability to go beyond paraphrasing/summarising.
This activity is designed to help you understand what might be involved
in doing so.

Read the following two extracts from an Analysis section and discuss
these questions:

\begin{enumerate}
\def\labelenumi{\arabic{enumi}.}
\tightlist
\item
  Which of the two examples is better? Why?
\item
  What have the authors done that makes the analysis better?
\end{enumerate}

\begin{tcolorbox}[enhanced jigsaw, breakable, bottomtitle=1mm, leftrule=.75mm, opacityback=0, colbacktitle=quarto-callout-note-color!10!white, opacitybacktitle=0.6, toprule=.15mm, colback=white, left=2mm, coltitle=black, bottomrule=.15mm, titlerule=0mm, title={Extract 1 - from a study about the impact of the COVID pandemic}, toptitle=1mm, colframe=quarto-callout-note-color-frame, arc=.35mm, rightrule=.15mm]

Sarah spoke of using her leisure time to plan and fantasise about her
wedding:

\emph{``it's enough of a fantasy\ldots{} it's meant to be stressful, but
it's really just been fun. It's like such a nice break from this humdrum
life we're living in right now'' (Sarah: lines 459-461).}

Organising a significant life event seems to be viewed as granting Sarah
an escape from her current monotony. While she refers to it as a
``fantasy'', potentially given uncertainty of the pandemic's
progression, she appears to perceive this as constructive and realistic
for her future self. Alternately, Mandy's leisure activity involved
impossible daydreams:

\emph{``I started daydreaming more ridiculous stuff\ldots{} I know it's
never going to come true 'cause it's not real, but it still entertains
me'' (Mandy: lines 443-444)}

While Mandy seems to be under no illusion that her daydreams are
feasible, daydreaming impossible scenarios appears to provide more of a
release than realistic situations that are subject to ever changing
government regulations (e.g.~travel bans) that may fall through.

\end{tcolorbox}

\begin{tcolorbox}[enhanced jigsaw, breakable, bottomtitle=1mm, leftrule=.75mm, opacityback=0, colbacktitle=quarto-callout-note-color!10!white, opacitybacktitle=0.6, toprule=.15mm, colback=white, left=2mm, coltitle=black, bottomrule=.15mm, titlerule=0mm, title={Extract 2 - from a study about social media use during the pandemic}, toptitle=1mm, colframe=quarto-callout-note-color-frame, arc=.35mm, rightrule=.15mm]

Some participants said it was difficult to fall asleep at night,
including those who knew they should fall asleep but could not
(\emph{``if I can't stay up and can't sleep''- Sarah}), or participants
subjectively thought that they did not need sleep or did not want to
fall asleep (\emph{``I am not really able to sleep at night, or like,
sleep isn't really needed.''} -Andrew)

Participants said that they used social media but that it was only
positive at certain points:

\emph{I feel like TikTok was fun in the first lockdown during the
summer,'cause of all this summer content I guess, but after that it
really went downhill. I don't know like, OK personally I deleted all my
apps about three weeks ago} (John)

John tells us here that he deleted all his apps, because they all went
`downhill'. Tik Tok was helpful during the first lockdown, but not after
it.

\end{tcolorbox}

\begin{tcolorbox}[enhanced jigsaw, breakable, bottomtitle=1mm, leftrule=.75mm, opacityback=0, colbacktitle=quarto-callout-caution-color!10!white, opacitybacktitle=0.6, toprule=.15mm, colback=white, left=2mm, coltitle=black, bottomrule=.15mm, titlerule=0mm, title={Show me the answer}, toptitle=1mm, colframe=quarto-callout-caution-color-frame, arc=.35mm, rightrule=.15mm]

Example 1 (COVID):\\
This is the better example, because it goes beyond paraphrasing. The
researcher is analysing the quotes and have engaged with the meaning
behind the words. Each quote is selected in a way that helps the author
build a narrative of what is happening in the data

Example 2 (sleep): This is the weaker example, because the author just
paraphrases what the participants say. This ends up not adding any real
value in terms of analysis, because the reader could just read the
quote. The narrative needs to analyse the data and show the reader why
they interpret something in a particular way.

\end{tcolorbox}

\section{Evaluation in discussions}\label{eval_discussion}

In this section, we aim to provide some guidance about how to go beyond
description in your writing. We will specifically focus on the
Discussion section for your RM2 qualitative report here, but you can
apply this to other types of writing. In particular, it should be
helpful for you as you proceed to your dissertation.

\subsection{What were your findings?}\label{what-were-your-findings}

In your discussion section, start off by telling us what you found in
your analysis. However, the key thing to think about here, in terms of
evaluation, is to relate your findings to the wider research. It will be
likely a good idea to consider on a theme by theme basis (if you are
doing themes with no sub-themes) or on a subtheme basis (if you have one
theme with 2-3 subthemes).

\textbf{Some things to consider:}

\begin{itemize}
\tightlist
\item
  What did you find? What was your specific contribution?
\item
  Thinking about your findings, what are the similarities and
  differences from other research findings/theory? Why might this be the
  case?
\item
  Situate your findings within the broader context - tell the reader
  about why your findings are important. This can help a reader
  understand the importance of the results.
\item
  Do your findings challenge existing research? If so, in what way?
\item
  Can you bring in any critical thinking?
\end{itemize}

\subsubsection{Activity 1}\label{activity-1-2}

Read the paragraphs below from
\href{https://journals.sagepub.com/doi/full/10.1177/1362361320908976?fbclid=IwAR0tMsMbmeATwp5RpoKA-yqUr3QzyHuBSuy00uMAvlqZepaS4Pv_48VgQjA}{Crompton
et al., (2020)} and think about what the authors are communicating, with
a view to considering the points above.

In particular, identify:

\begin{itemize}
\tightlist
\item
  How the authors summarised their results
\item
  The contribution of the themes to the evidence base
\item
  How the authors discuss their results in relation to existing research
\item
  How the authors relate their findings to theory
\item
  How the authors link the findings to `broader issues'
\end{itemize}

\begin{tcolorbox}[enhanced jigsaw, breakable, bottomtitle=1mm, leftrule=.75mm, opacityback=0, colbacktitle=quarto-callout-note-color!10!white, opacitybacktitle=0.6, toprule=.15mm, colback=white, left=2mm, coltitle=black, bottomrule=.15mm, titlerule=0mm, title={Study 1: Crompton et al., 2020}, toptitle=1mm, colframe=quarto-callout-note-color-frame, arc=.35mm, rightrule=.15mm]

This study aimed to examine the experiences of autistic adults spending
time with autistic and non-autistic family and friends using a thematic
analysis framework. Social relationships are an important, though often
complicated, part of autistic people's lives. Previous research has
tended to focus on autistic people's relationships with (assumed)
non-autistic friends and family. Here, we specifically contrasted
relationships across and within neurotypes. The analysis revealed three
themes: cross-neurotype understanding, minority status and belonging.
The themes help us understand why relationships between autistic and
non-autistic people might be so challenging, and how relationships
between autistic people are different.

The results align with previous research on the challenges that autistic
people face when interacting with non-autistic others, but highlight
that interactions with other autistic people are fundamentally
different. All participants reported that spending time with
non-autistic family and friends involved specific difficulties, which
were not experienced when interacting with other autistic friends and
family. This aligns with the double-empathy theory of autism which
suggests that autistic and non-autistic people have a mutual difficulty
in understanding and empathising with one another due to differences in
how each person understands and experiences the world, rather than
because of a communicative deficit on the part of the autistic person
(Milton, 2012). Neurotypical people have been shown to overestimate how
ego-centric their autistic family members are (Heasman \& Gillespie,
2018), and overestimate how helpful they are to autistic people (Heasman
\& Gillespie, 2019). Our findings suggest that this translates into
real-world difficulties in interactions with neurotypical friends and
family that may affect the mental health, well-being and self-esteem of
autistic people.

\end{tcolorbox}

\subsection{What are the implications of the
study?}\label{what-are-the-implications-of-the-study}

In addition to theoretical implications you would discuss in a
quantitative study, there is scope to discuss broader implications here.
Try not to go too far in these implications (e.g.~saying that the
NHS/University could make changes to how they run mental health
services\ldots this would require a much more extensive body of work).

\textbf{Some things to consider:}

\begin{itemize}
\tightlist
\item
  How could these findings be applied to real-life situations?
\item
  What impact might this have for society in the future (even if with
  the caveat that it would depend on future research)?
\item
  Can you back up your assertions with reference to the literature?
  Arguments here are much stronger if you are able to cite papers to
  support.
\end{itemize}

\subsubsection{Activity 2}\label{activity-2-1}

Read the paragraph below from
\href{https://www.sciencedirect.com/science/article/pii/S1750946717301277}{Robertson
et al.~(2018)} and think about how the authors have communicated the
implications of the research.

In particular, identify:

\begin{itemize}
\tightlist
\item
  What the implications of the research are
\item
  How the authors link the findings to practical suggestions
\item
  How the authors back up these suggestions to the existing literature
  base
\end{itemize}

\begin{tcolorbox}[enhanced jigsaw, breakable, bottomtitle=1mm, leftrule=.75mm, opacityback=0, colbacktitle=quarto-callout-note-color!10!white, opacitybacktitle=0.6, toprule=.15mm, colback=white, left=2mm, coltitle=black, bottomrule=.15mm, titlerule=0mm, title={Study 2: Robertson et al., 2018}, toptitle=1mm, colframe=quarto-callout-note-color-frame, arc=.35mm, rightrule=.15mm]

There are a number of potential clinical applications for the findings
detailed in this paper. They echo the modifications that are already
undertaken for anxiety interventions designed for children and
adolescents with ASD (Walters, Loades, \& Russell, 2016). For example,
it may be important to include support for individuals undertaking an
intervention for their anxiety. This has already been identified as
helpful in anxiety management interventions for children with ASD
(Sofronoff et al., 2005), but the evidence here suggests that this might
also be the case for adults. Furthermore, as adults reported feeling
that they enjoyed interacting with other autistic people, it would be
important to test whether running autism-only groups in terms of anxiety
interventions helps outcomes or other aspects related to the feasibility
of anxiety management programs (e.g.~retention of participants).

\end{tcolorbox}

\subsection{Limitations}\label{limitations}

All research will have limitations - there is never going to be a
perfect study. Research is about weighing up the options and choosing
the best design for the specific question we want to answer. Remember
that qualitative research is not automatically limited because it is not
quantitative. Therefore, be very careful not to use criteria we have for
quantitative studies (e.g.~generalisability, reliability, small sample
sizes) and apply them to qualitative studies. Doing so would be similar
to critiquing a behavioural study for not asking participants about
their experience of interacting with the stimuli ;)

\textbf{Some things to consider:}

\begin{itemize}
\tightlist
\item
  Do you think the quality, source, or types of data (or the analytic
  process) might have affected the methodological integrity in some way?
  If so, how?
\item
  Remember to link your limitations to the evidence base. This makes
  your argument much stronger than if you do not cite to back your
  points up.
\item
  If you want to mention generalisability, perhaps consider the concept
  of transferability instead, which is much more relevant to
  qualitative.
\item
  Try to aim for depth rather than breadth here - it's much stronger to
  have two points addressed comprehensively than have 12 limitations,
  each covered in a single sentence (you don't want to make your
  discussion section a limitations-fest!).
\end{itemize}

\subsubsection{Activity 3}\label{activity-3}

Read the paragraph below from
\href{https://journals.plos.org/plosone/article?id=10.1371/journal.pone.0210450}{Toivo
\& Scheepers (2019)}. Although this is not a qualitative paper, the
following paragraph exhibits many of the strengths noted above.

In particular, identify:

\begin{itemize}
\tightlist
\item
  What was the limitation observed?
\item
  How do the authors discuss this limitation, and consider the effect it
  might have had?
\item
  How do the authors back up their arguments with evidence from the
  literature base?
\end{itemize}

\begin{tcolorbox}[enhanced jigsaw, breakable, bottomtitle=1mm, leftrule=.75mm, opacityback=0, colbacktitle=quarto-callout-note-color!10!white, opacitybacktitle=0.6, toprule=.15mm, colback=white, left=2mm, coltitle=black, bottomrule=.15mm, titlerule=0mm, title={Study 3: Toivo et al., 2019}, toptitle=1mm, colframe=quarto-callout-note-color-frame, arc=.35mm, rightrule=.15mm]

One concern might be that we did not explicitly control for L2
proficiency in our sample of bilinguals, which is likely to affect
cognitive effort in processing L2. However, all of our bilingual
participants (a) self-reported to be highly proficient in English, (b)
were Glasgow University students using English as their language of
study, and (c) had been living in an English-speaking country for
minimum of three months (average: two years and four months) when the
study was conducted. Together with the fact that post-trial word
recognition rates were well above 90\% across participant groups and
conditions (Fig 2), we do not believe that variation in L2 proficiency
was large enough to have a major impact on our results. That said, we
agree that proficiency should ideally be controlled for in future
research, e.g.~by adding a suitable measure as another person-specific
covariate in the analyses. Previous research had identified proficiency
as a potential mediating factor of perceived emotionality in the
Bilingualism and Emotions Questionnaire {[}50{]}, and it had also been
suggested that bilinguals with close-to-native proficiency in L2 show
similar affective processing in L1 and L2 {[}19{]}.

\end{tcolorbox}

\subsection{Anything else?}\label{anything-else}

The points above are the key components that we would expect to see in a
Discussion section, with some guidance on how to go beyond description.
However, you might also consider strengths and future research. It's
also good to include a brief conclusion. You won't have the word count
available to go into the same detail as in these papers, so try to use
the examples above to really focus on what you can incorporate within
your Discussion section to strengthen your evaluative writing.

\chapter{Qualitative report}\label{ch12_qualitative_report}

In this section, you will find information and guidance about your
qualitative report and a report writing guide. Please make sure to
familiarise yourself with the \hyperref[ch4_AIS_report]{Assessment
Information} before starting to work on the report. We have also
provided a
\href{https://gla-my.sharepoint.com/:w:/g/personal/ashley_robertson_glasgow_ac_uk/EQM2Fzsi0FJAlDKTJXgi7-8Brm7ITIBw5ta3MaXSiDZfrg?e=FMLoM4}{suggested
timeline} for the qualitative report and portfolio

For this assignment, you will be asked to write an individual
qualitative report on one of five secondary data topics. You should
conduct Thematic Analysis on your chosen dataset. A video on the
Qualitative report will be provided on Moodle - please make sure to
watch this. We provide guidelines on the report
\hyperref[qual_report_guide]{later in this chapter}

\section{Report datasets and how to use UK data
service}\label{report-datasets-and-how-to-use-uk-data-service}

Before accessing your datasets, watch
\href{https://echo360.org.uk/media/ac815f47-bb97-40cc-8c34-d82d74e4bb0e/public}{this
video} on how to use UK data service. It shows you an example of how to
download the data with one of the datasets. Please note that some of the
datasets are on sensitive topics, and contain information that may be
upsetting: make sure to familiarise yourself with the interview
schedules and have a quick scan of the data before you decide which
dataset to work with.

\href{https://beta.ukdataservice.ac.uk/myaccount/login}{The UK data
service can be accessed here}

You can choose one of the five available datasets:

\begin{enumerate}
\def\labelenumi{\arabic{enumi}.}
\tightlist
\item
  \href{https://datacatalogue.ukdataservice.ac.uk/studies/study/854556?type=reshare\#details}{Interviews
  with disabled young people and carers}
\item
  \href{https://datacatalogue.ukdataservice.ac.uk/studies/study/855322?type=reshare\#details}{LGBT+
  networks in the workplace}
\item
  \href{https://datacatalogue.ukdataservice.ac.uk/studies/study/851244?type=reshare\#details}{Long-term
  relationships}
\item
  \href{https://datacatalogue.ukdataservice.ac.uk/studies/study/855486?type=reshare\#details}{Patient
  descriptions of PMS}
\item
  \href{https://datacatalogue.ukdataservice.ac.uk/studies/study/856114?type=reshare\#details}{Student
  experiences of loneliness and social connectedness}
\end{enumerate}

Please note that some of the datasets may include pre-existing coding in
the data folder, or published analyses; you should \textbf{not} use
these as it will classify as academic misconduct; you are expected to
develop the full analysis yourself.

\section{Developing a research
question}\label{developing-a-research-question}

As part of writing up your qualitative report, we ask that you develop a
research question to help you guide your enquiry. Below, we have some
guidance on how to get started.

\subsection{Getting started with the
data}\label{getting-started-with-the-data}

For this report, you are using secondary data. This means that the data
have already been collected - before you start thinking about your
research question, have a look at the interviews to see what was asked
in the interviews. Usually, most interviews are semi-structured. This
means that there are likely to be some questions asked consistently, but
there will also be differences, depending on the responses by the
participants.

Looking at the data before you decide on your RQ might feel a bit
strange after RM1, but remember that you are not doing
hypothesis-testing with qualitative designs - it is quite normal for a
qualitative research question to change during the research process, and
this is completely fine as you are not statistically testing for a
specific effect.

Be mindful that you are not making a research question that is actually
one of the interview questions - your RQ should be broader than an
interview question.

One of the most common mistakes that students make when framing their
research question in RM2 is that it is a little too `quant' in nature.
This makes perfect sense, as most of your research methods education so
far in the programme has been focused on quant! However, for your qual
report, your RQ should be a little different to the RQs you are used to.

Look at the following examples:

\begin{enumerate}
\def\labelenumi{\arabic{enumi}.}
\tightlist
\item
  How does poor sleep impact on wellbeing?
\item
  What is the difference in wellbeing between those who belong to
  university and those who don't?
\end{enumerate}

Both questions are a bit too quantitative in nature (i.e.~1 is causal
and 2 is about a difference between groups). They would also benefit
from thinking about what is being looked at, specifically (e.g.~views,
opinions, experiences).

We'd perhaps rephrase these as follows:

\begin{enumerate}
\def\labelenumi{\arabic{enumi}.}
\tightlist
\item
  What are students' perceptions of the relationship between sleep and
  wellbeing?
\item
  What are online conversion students' experiences of belonging to
  University? How do these experiences relate to perceived wellbeing?
\end{enumerate}

These questions are now a bit more in line with qualitative and are
specific about what they are asking in each (and also what the group of
interest is). The second one covers a bit more than the first, which is
why it is split into two (a main RQ and a subsidiary one)

\subsection{Some tips for developing a research
question}\label{some-tips-for-developing-a-research-question}

\begin{itemize}
\tightlist
\item
  \textbf{Communication:} Keep your research question simple and
  specific. If you have a very broad question, it can be helpful to
  break the question down into smaller sub-questions.
\item
  \textbf{Type of question:} We want the RQ to tap in to more than
  `descriptive' experience but also offer scope to explore the `how' and
  `why' of psychological experience--Is the question asking about
  experiences, understanding, accounts of practice, or influencing
  factors?\\
\item
  \textbf{Practical constraints:} Can the question realistically be
  explored using thematic analysis and using secondary data? Try not to
  cover too much
\item
  \textbf{Guiding:} Remember that your RQ guides your enquiry but
  shouldn't be restrictive - it's not the case that you are unable to
  report themes that don't directly answer your RQ.
\item
  \textbf{Flexibility:} Research questions are more flexible than
  hypotheses-it is acceptable for the focus of the RQ to be refined
  during/after data collection.
\end{itemize}

\subsection{Formative feedback on your
RQ}\label{formative-feedback-on-your-rq}

If you would like to receive formative feedback on your RQ, you can
submit it on
\href{https://padlet.com/wilhelmiinatoivo/rm2-odl-research-questions-f844153314j3pwr2}{this
Padlet} by the end of Semester 2 Week 5. Please note that RQs submitted
later than that will not receive feedback, and you would need to drop
into office hours instead.

If you would like to discuss your RQ, you can also attend Ashley's and
Wil's office hours.

\subsection{Choosing your interviews}\label{choosing-your-interviews}

Once you have developed your research question, you should choose the
interviews you are going to analyse for the report. Some of the datasets
are really big and involve multiple components (see the qualitative
report video for further guidance on the datasets).

\begin{itemize}
\tightlist
\item
  \textbf{You should use a maximum of 2-3 interviews for your analysis}
\item
  If you are working with the couples dataset, you can do 4 interviews
  if you want to do two couples
\item
  If you are working with the couples dataset, do not use the collage
  interviews
\item
  If you are working with the LGBT dataset, only use interview data for
  your report
\item
  If you are working with the PMS dataset, only use patients' interviews
  (the expert interviews do not lend themselves to the type of
  qualitative analysis we ask you to do)
\item
  If you work with the student loneliness dataset, only use the focus
  group data
\end{itemize}

There is no one protocol for choosing your interviews. Sometimes you may
choose depending on which interviews have the most/richest data about
your specific RQ (as the datasets are suited for a large number of
different RQs), or in other cases you may choose the interviews based on
particular demographics (if this is a part of your RQ, for example,
looking specifically at carers', mothers', or men's perspectives). We
recommend scanning through multiple interviews to gauge which ones might
work best for your particular RQ (do not summarise them with AI - the
datasets are from real research participants and should not be uploaded
to any AI platforms). You will be asked to report your data selection
process in the methods of your report.

\section{Qualitative report guide}\label{qual_report_guide}

This guide is intended to help you in the write up of your qualitative
project report. When writing your report, remember to use the other
resources available to you:

\begin{itemize}
\tightlist
\item
  Your feedback for your RM1 report
\item
  Chapter 13 in \href{https://go.exlibris.link/5g3nFfGx}{Braun \& Clarke
  (2013)} is excellent support for writing up qualitative research
\item
  \href{https://go.exlibris.link/Z9qxllzH}{Braun \& Clark (2021) book}
\item
  All the resources available in this book (particularly sections for
  \hyperref[ch8_data_analysis]{8. Data analysis} and
  \hyperref[c11_writing]{11. Academic writing}
\item
  \textbf{25/26 cohort reading list:}:
  \href{https://rl.talis.com/3/glasgow/lists/30BEDB96-5F9D-9F9E-630A-205B91FA34AB.html?lang=en-GB\#B226C8C1-1926-8328-A8A3-06DC389D9809}{Published
  articles that use a Thematic Analysis}
\item
  \href{https://owl.purdue.edu/owl/research_and_citation/apa_style/apa_formatting_and_style_guide/general_format.html}{A
  helpful resource for APA formatting is OWL Purdue}
\end{itemize}

Use \hyperref[ch4_AIS_report]{the assessment criteria} as a checklist
before you submit the report, to ensure you have addressed all points.
If you have questions post on Teams and/or visit our office hours. All
the best in writing your report.

\subsection{Required Submissions}\label{required-submissions}

Two separate submissions are required for this report:

\begin{itemize}
\tightlist
\item
  \textbf{Submission 1} should contain the report itself, along with the
  specified sections and appendices.

  \begin{itemize}
  \tightlist
  \item
    Cover page with title (does not count towards the word limit)
  \item
    Abstract
  \item
    Introduction
  \item
    Method
  \item
    Analysis
  \item
    Discussion
  \item
    References (does not count towards word limit)
  \item
    Appendices (does not count towards word limit)
  \end{itemize}
\item
  \textbf{Submission 2} should contain the required information about
  the process of your analysis, using the template provided. Submission
  2 does not count towards the word limit. Please note that Phase 6 is
  not included as this is what you include in your Analysis section

  \begin{itemize}
  \tightlist
  \item
    Phase 1 (Familiarisation)
  \item
    Phase 2 (Doing Coding)
  \item
    Phase 3 (Generating Initial Themes)
  \item
    Phase 4 (Developing and Reviewing Themes)
  \item
    Phase 5 (Refining, Defining and Naming Themes)
  \end{itemize}
\end{itemize}

\subsection{Submission 1: The Qualitative
Report}\label{submission-1-the-qualitative-report}

\subsubsection{General writing style}\label{general-writing-style}

\begin{itemize}
\item
  Your writing should be clear, concise and easy to follow - remember
  that academic writing does not need to be complicated to be good. One
  tip is to
  \href{https://support.microsoft.com/en-gb/office/listen-to-your-word-documents-5a2de7f3-1ef4-4795-b24e-64fc2731b001}{use
  word to read your report back to you} -- sometimes it's easier to hear
  where a point may be overly complex or long and where you might want
  to cut it down.
\item
  Maintain a neutral, academic tone throughout. Avoid a `journalistic'
  tone, where you convince the reader based on using emotive language.
  Instead, show the reader why you think the way you do and the evidence
  base that underlies this.
\item
  Discuss one idea/argument per paragraph. An excellent resource to help
  with this is the
  \href{https://explorationsofstyle.com/2011/02/09/reverse-outlines/}{reverse
  outline}.
\item
  You should format your headings, in-text citations and references in
  APA 7th style.
\item
  APA writing style encourages the use of active voice (``I coded the
  data'' rather than ``the data were coded'') but if you feel more
  comfortable writing in passive voice, feel free to do so - just make
  sure to be consistent throughout and be careful not to make your
  writing overly complex.
\item
  Tenses are not always fixed per se (as it depends on how a sentence is
  structured). However, generally, the following applies -- past tense
  should be used when you have done something already (like your study).
  For example, we would usually see past tense in the Methods (unless
  doing a registered report) and the Analysis section. When talking
  about papers that are already published (in the Intro and Discussion),
  again these would typically be in past tense.
\item
  You can use first person singular (I) or plural (we) in the report
  itself if you want to, but in the reflexivity section make sure to
  write in first person singular. This is because this is your
  reflection of your own positionality to the RQ and the data, and it's
  often very personal.
\end{itemize}

\subsubsection{Title}\label{title}

The title should be included as part of the cover page. It should
reflect the research question(s) and define your sample (e.g., UG
students experience of\ldots). Try to be specific - your title should
include:

\begin{itemize}
\tightlist
\item
  Who
\item
  What
\item
  How
\end{itemize}

\textbf{Examples:}

\begin{itemize}
\tightlist
\item
  A Thematic Analysis \emph{(how)} of the panic attack experiences
  \emph{(what)} of primary aged children in inner city Schools in the UK
  \emph{(who)}
\item
  Barriers and enablers to modifying sleep behaviour \emph{(what)} in
  adolescents and young adults \emph{(who)}: A qualitative investigation
  \emph{(how)}
\end{itemize}

\subsubsection{Abstract (suggested word count: 100-150
words)}\label{abstract-suggested-word-count-100-150-words}

In an academic journal, the function of an abstract is to give readers a
snapshot of the whole paper. They then use this to decide whether or not
to read the full article. Therefore, you should make sure you cover all
four main sections of your report, summarising them all succinctly.

Aim to summarise all sections of the report in 150 words, including; 1)
what does the existing literature tell us about the topic 2) aim of the
study 3) brief methodology, 4) approach to analysis , 5) main findings,
and 6) key aspects of the discussion. Keep in mind that you are
summarising your research for a non-expert and you want to ``entice''
them to read more.

\begin{tcolorbox}[enhanced jigsaw, breakable, bottomtitle=1mm, leftrule=.75mm, opacityback=0, colbacktitle=quarto-callout-warning-color!10!white, opacitybacktitle=0.6, toprule=.15mm, colback=white, left=2mm, coltitle=black, bottomrule=.15mm, titlerule=0mm, title={Tips}, toptitle=1mm, colframe=quarto-callout-warning-color-frame, arc=.35mm, rightrule=.15mm]

\begin{itemize}
\tightlist
\item
  Start out by writing four sentences, each summarising one section of
  the report (introduction, methods, analysis, discussion). Now, look at
  the six sections above and see where you might need to add a bit of
  detail.
\item
  One common mistake is to ignore the discussion, saying something like
  `Findings are discussed in relation to existing research'. This
  doesn't tell us anything about what your overarching thoughts are
  about your findings and how they relate to the literature base.
\end{itemize}

\end{tcolorbox}

\subsubsection{Introduction (suggested word count 750-850
words)}\label{introduction-suggested-word-count-750-850-words}

Some activities to support you with writing your Analysis section can be
found in \hyperref[eval_intro]{Chapter 11}.

The introduction should give readers enough background to understand
your rationale for further research, whether your study is qualitative
or quantitative. Provide a clear, critical review of the most relevant
published research and highlight the key issues and debates in the
field. Use these to build a logical argument for your study.

Follow the APA-style funnel structure: begin with a broad overview of
the topic, then narrow the focus until you reach your specific rationale
and research question(s).

Because your word count is limited, focus on recent and directly
relevant studies---especially the latest qualitative research in this
specific area. Show how your project builds on the most recent work to
strengthen your rationale. Remember to cite all sources used, including
government reports, websites, and media items.

A resource for helping you to develop evaluation in your introductions
can be found at the start of \hyperref[ch11_writing]{Chapter 11}.

\begin{figure}[H]

{\centering \includegraphics[width=1\linewidth,height=\textheight,keepaspectratio]{images/intro.png}

}

\caption{Broad-to-narrow structure of an introduction}

\end{figure}%

You should end with a statement of what the aim of your study was, and
the research question.

When reviewing previous studies, don't just describe them - evaluate
them. Highlight their limitations, but only those your study can
realistically address. Consider what the findings mean, note any
conflicting results, and think about why they differ. This should
naturally build toward a clear rationale for your own study.

Then, explain your study and clearly state your research question(s),
ensuring they connect directly to the issues raised in your
introduction.

In summary:

\begin{itemize}
\tightlist
\item
  Review and critically evaluate relevant theories and research.
\item
  Justify your study by linking gaps or limitations in past work to your
  research questions.
\item
  Clearly outline your main and subsidiary research questions, ensuring
  these are phrased in a way that aligns with a qualitative approach.
\end{itemize}

Remember: just like in a quantitative report, this section shows what
has been done before and why your study is needed.

\begin{tcolorbox}[enhanced jigsaw, breakable, bottomtitle=1mm, leftrule=.75mm, opacityback=0, colbacktitle=quarto-callout-warning-color!10!white, opacitybacktitle=0.6, toprule=.15mm, colback=white, left=2mm, coltitle=black, bottomrule=.15mm, titlerule=0mm, title={Tips}, toptitle=1mm, colframe=quarto-callout-warning-color-frame, arc=.35mm, rightrule=.15mm]

\begin{itemize}
\tightlist
\item
  In your rationale, remember to tell the reader why the current study
  (i.e.~the topic and group of participants) is \textbf{important}.
\item
  If identifying limitations of existing research, make sure to pick
  something that you are able to address in your study. If you say that
  one issue is that heart rate wasn't gathered, but you don't do this in
  your study, this can be confusing for a reader.
\end{itemize}

\end{tcolorbox}

\subsubsection{Method (suggested word count 350-450
words)}\label{method-suggested-word-count-350-450-words}

The aim of the method section is to clearly report how your study was
conducted -- the reader should be able to re-run your research from the
details provided. In this case, contextualise the original study that
your secondary data are from, report details of the original data
collection, explain your process of refining and selecting the data used
for your project, and describe the qualitative analysis you used. True
`replication' is difficult due to the central role of the researcher in
influencing some aspects of the research such as analysis. This is not
an experiment.

Follow these subheadings for your method section for this report:

\textbf{Secondary data}

In this section you should describe the original study and give the
reader context about their data collection. Understanding the specific
context of each study is a key element of qualitative research, which is
why it is really important for you to explain that in order for your
secondary analysis to be reported openly.

What was the original study about? Who conducted it? Where are the data
available? You should cite your chosen dataset in APA style (do not
provide a link in text) - each dataset has this information on the UK
data service page under the Citation and copyright section. How many
interviews did they conduct, what was the methodology they used for the
interviews (you should find this information on the UK Data Service page
for your chosen data set)? Who were their participants? What were the
participants asked about?

If you have demographic information about the participants on the UK
Data Service Website, you may wish to include this, but don't worry if
your chosen dataset doesn't provide this - we don't expect you to report
any information that was not available to you.

\textbf{Data selection}

In this section you should describe your procedure for choosing your
final data set. How many interviews did you choose? How did you choose
the specific ones you decided to use for your analysis?

\textbf{Ethics}

The information you are able to provide in this section will depend on
the information made available within the dataset. Some things you might
want to comment on include: anonymisation procedures (e.g.~has
information been replaced or redacted), rapport building,
information/consent/debrief, consent as a process. You will only be able
to include information that has been made available in the dataset, so
do not worry if you do not know a lot of detail.

\textbf{Reflexivity}

In qualitative research it is important to acknowledge the role of the
researcher in interpretation of the data
\href{https://bpspsychub.onlinelibrary.wiley.com/doi/10.1348/014466599162782}{(Elliot
et al., 1999)}.

This is where you question your own motives and attitudes in doing this
project. Obviously we told you that you had to do it. But, what
motivated you to choose the dataset you did? What assumptions did you
hold prior to beginning the research? Had you considered issues around
the topic previously? Did they match what the interviews say, or did you
disagree with them? Did this affect which interviews you chose to
include in your analysis? And if so, how did this impact on your
analysis? Because different people would interpret the data differently,
it is useful for you to expose your own attitudes at this point so that
others can see how you have impacted on the analysis.

See Braun \& Clarke (2013, pp36-37, 303-304) and activities around
reflexivity. Reflexive analysis should be concise for this project,
perhaps around 3-4 sentences. However, you will also have to include a
\href{https://gla-my.sharepoint.com/:w:/g/personal/ashley_robertson_glasgow_ac_uk/IQC1PIxhBlndR74S_nfeKHNlAW5uQmGOpoRXCpLHC2fI_fI?e=aunShJ}{reflexivity
diary} as a part of your appendices.

\textbf{Data Analysis}

In this section, explain:

\begin{itemize}
\tightlist
\item
  What analysis method you used and your theoretical approach (see the
  `Thinking about theory' mini-lecture for a reminder).
\item
  The stages of your qualitative analysis. Be clear and specific -
  qualitative work is sometimes criticised for being vague. Use
  recognised guidelines (e.g., Smith for IPA; Braun \& Clarke for
  Reflexive Thematic Analysis).
\end{itemize}

When outlining the stages:

\begin{itemize}
\tightlist
\item
  Give a brief summary of the full process, linked to the relevant
  framework (e.g., RTA). Around three to four sentences total is enough
  for this length of report.
\item
  You can keep some phases short (e.g., ``familiarisation with the data
  (Phase 1); generating codes (Phase 2)\ldots{}'').
\item
  Aim for transparency - explain what happens between coding and theme
  development.
\item
  Describe the researcher as actively identifying themes, not waiting
  for themes to ``emerge,'' which incorrectly suggests they are already
  there waiting to be found.
\end{itemize}

\begin{tcolorbox}[enhanced jigsaw, breakable, bottomtitle=1mm, leftrule=.75mm, opacityback=0, colbacktitle=quarto-callout-warning-color!10!white, opacitybacktitle=0.6, toprule=.15mm, colback=white, left=2mm, coltitle=black, bottomrule=.15mm, titlerule=0mm, title={Tips}, toptitle=1mm, colframe=quarto-callout-warning-color-frame, arc=.35mm, rightrule=.15mm]

\begin{itemize}
\tightlist
\item
  In the methods section, try to be precise and succinct in your writing
  (this applies to the whole report, but is particularly important
  here). Think about what the reader needs to know and how to best
  communicate this.
\item
  Common errors include inaccurate information (e.g.~the wrong number of
  Phases in Braun \& Clarke's framework) and a superficial reflexivity
  section.
\end{itemize}

\end{tcolorbox}

\subsubsection{Analysis (suggested word count: 750-900
words)}\label{analysis-suggested-word-count-750-900-words}

Some activities to support you with writing your Analysis section can be
found in \hyperref[eval_analysis]{Chapter 11}.

The analysis section is where you present the evidence you've collected.
It works like a Results section in a quantitative report.

Your goal is to make the analysis transparent, showing clearly how your
interpretations come from the transcripts. Use extracts from the
transcripts to support all quotes and interpretations.

Please note that some of the datasets include pre-existing coding; you
should \textbf{not} use these as it will classify as academic
misconduct; you are expected to develop the full analysis yourself
through the six phases.

\textbf{Guidelines}

\begin{itemize}
\tightlist
\item
  Choose a maximum of two themes or one theme with two sub-themes to
  include in the Analysis section, even if you found more. You will not
  be able to go into enough depth if you present more themes than this,
  due to the word count restrictions.
\item
  If you want to use parts of a quote, rather than the full quote
  (e.g.~maybe you want to remove the middle sentence, which doesn't add
  anything), you can indicate you omitted something with ellipses.
\item
  Never alter any of the words or phrasing that is used, even if you
  think you know what the participant was going to say. This would be as
  unacceptable as editing numbers in a quantitative analysis!
\item
  Try to strike a balance between quotes and interpretation. We
  sometimes see lots of quotes with very little analysis or lots of
  analysis with very few quotes, both of which have problems.
\item
  \textbf{Do not} include citations in the analysis section of this
  report. Sometimes in published qualitative papers this is done, but
  due to the limited word count, we ask you to not do it for this
  assignment. Here, you should connect your research into existing
  literature in the discussion section.
\end{itemize}

\begin{tcolorbox}[enhanced jigsaw, breakable, bottomtitle=1mm, leftrule=.75mm, opacityback=0, colbacktitle=quarto-callout-important-color!10!white, opacitybacktitle=0.6, toprule=.15mm, colback=white, left=2mm, coltitle=black, bottomrule=.15mm, titlerule=0mm, title=\textcolor{quarto-callout-important-color}{\faExclamation}\hspace{0.5em}{Generative AI and data}, toptitle=1mm, colframe=quarto-callout-important-color-frame, arc=.35mm, rightrule=.15mm]

Do not, \textbf{under any circumstances}, put the transcripts or quotes
from the transcripts into GenerativeAI (this includes CoPilot). This
will not have been agreed as part of the original ethics and
participants will not have consented for this use of the data.

Failing to follow these instructions puts you at risk of academic and
research misconduct.

\end{tcolorbox}

\textbf{Structuring the Analysis section}

\begin{itemize}
\tightlist
\item
  Start with a short paragraph summarising the main themes
  \textbf{without} using any quotes. This should not be too long -
  perhaps around a sentence or two.
\item
  If you have two themes:

  \begin{itemize}
  \tightlist
  \item
    Present each theme using a sub-heading
  \item
    Give the reader a brief overview of what the theme is. It is likely
    that you will have to select a few key points that encapsulate your
    theme, rather than describe everything in your theme in detail. For
    a dissertation, where you have more available word count, you will
    be able to cover more.
  \item
    Once you've introduced the theme, tell the reader about what you
    found in your analysis of this theme. Make sure to incorporate
    quotes to support your analytical points.
  \end{itemize}
\item
  If you have one theme and two sub-themes:

  \begin{itemize}
  \tightlist
  \item
    Have sub-headings for the overall theme as well as the sub-themes.
  \item
    Briefly explain the overarching theme and then how the sub-themes
    relate to the larger theme. Do not present quotes at this point.
  \item
    Then, present each of the sub-themes in turn, under sub-headings. In
    each of these, tell the reader what you found in the analysis of
    your theme and incorporate quotes to support your analytical points.
  \end{itemize}
\end{itemize}

\textbf{Presenting Quotes}

Use \emph{long quotes} by placing them in a separate, centred indented
block. \emph{Short quotes} can be included within the sentence. Due to
the word count, you are likely to be able to do a maximum of one or two
longer quotes per theme or sub-theme.

You can find guidance on presenting participant quotes in APA format
\href{https://apastyle.apa.org/style-grammar-guidelines/citations/quoting-participants}{in
the recommended resources}. Although the word limit suggested here is 40
words, this is quite long and so we suggest taking it on a case by case
basis (e.g.~having a quote of 39 words within a sentence might be
correct but also might be confusing to read). You can run anything past
us if you are unsure.

We have also
\href{https://gla-my.sharepoint.com/:w:/g/personal/ashley_robertson_glasgow_ac_uk/EWfG8vjCAvtPjYPEJjQWxl8BPHlfzuI6weVkwIh_MVJSTQ?e=x7OT9m}{prepared
an example} of how to format quotes in APA style.

\textbf{Going Beyond Paraphrasing}

A strong analysis section will \emph{analyse} the data, going beyond
summary and paraphrasing. What do you add, as the analyst? You will have
much more knowledge of the data, and it is your job to pass on that
insight. If you paraphrase the data, the reader will be able to see this
in the quotes, but the analysis then doesn't add anything extra.

\begin{tcolorbox}[enhanced jigsaw, breakable, bottomtitle=1mm, leftrule=.75mm, opacityback=0, colbacktitle=quarto-callout-important-color!10!white, opacitybacktitle=0.6, toprule=.15mm, colback=white, left=2mm, coltitle=black, bottomrule=.15mm, titlerule=0mm, title=\textcolor{quarto-callout-important-color}{\faExclamation}\hspace{0.5em}{Important}, toptitle=1mm, colframe=quarto-callout-important-color-frame, arc=.35mm, rightrule=.15mm]

Instead of telling the reader what was talked about, explain
\textbf{what the things that were talked about reveal or suggest} about
the phenomenon you are studying.

\end{tcolorbox}

Avoid writing themes that are topic summaries, as these are at a
descriptive level.

\begin{itemize}
\tightlist
\item
  A theme of ``Stress and workload'' is mainly descriptive. It tells us
  the topics that participants mentioned, but not why they matter or how
  they relate.
\item
  A theme of ``Stress is an unavoidable requirement of `being a good
  student'\,'' is interpretative. It goes further by explaining what
  stress \emph{means} for the participants. For example, the data might
  show that students see stress as tied to their identity as students
  and that this is something that they expect, rather than something to
  avoid.
\end{itemize}

This kind of deeper interpretation goes beyond the superficial to show
the reader what the data suggests, rather than just summarising it.

\begin{tcolorbox}[enhanced jigsaw, breakable, bottomtitle=1mm, leftrule=.75mm, opacityback=0, colbacktitle=quarto-callout-warning-color!10!white, opacitybacktitle=0.6, toprule=.15mm, colback=white, left=2mm, coltitle=black, bottomrule=.15mm, titlerule=0mm, title={Tips}, toptitle=1mm, colframe=quarto-callout-warning-color-frame, arc=.35mm, rightrule=.15mm]

\begin{itemize}
\tightlist
\item
  There is no single ``correct'' interpretation, but you must show
  clearly how your interpretation is supported by the quotes you use.
\item
  Go beyond describing the data---explain the underlying meaning.
\item
  Choose quotes that directly support the point you are making; the
  connection should be obvious to the reader.
\item
  Follow the guidelines: include the recommended number of
  themes/sub-themes, and balance quotations with your interpretation.
\end{itemize}

\end{tcolorbox}

\subsubsection{Discussion (suggested word count: 750-850
words)}\label{discussion-suggested-word-count-750-850-words}

Some activities to support you with developing evaluation in your
Discussion section can be found in \hyperref[eval_discussion]{Chapter
11}.

The Discussion is where you bring together your findings and consider
these in light of the literature (and theory). It is also where you
consider the wider importance of your findings, reflect on limitations
and think about potential implications and future directions.

\textbf{Start with a brief recap}

\begin{itemize}
\tightlist
\item
  Remind the reader of the theme/s you report in the Findings and relate
  these to your Research Question.
\end{itemize}

\textbf{Connect your findings to past research}

\emph{Note: This should be the largest section within your Discussion.}

\begin{itemize}
\tightlist
\item
  Here, compare your results with the studies and theories you mentioned
  in the Introduction, as well as any new studies/theories you want to
  incorporate.

  \begin{itemize}
  \tightlist
  \item
    Say whether your findings support or contradict those from the
    literature base. However, also go beyond this and tell us what this
    \textbf{means}, more broadly. Why might this be important?
  \item
    Say whether your findings extend/challenge what is already known. If
    so, tell us what this means/why it is important, and back your
    points up with evidence.
  \end{itemize}
\end{itemize}

\begin{tcolorbox}[enhanced jigsaw, breakable, bottomtitle=1mm, leftrule=.75mm, opacityback=0, colbacktitle=quarto-callout-important-color!10!white, opacitybacktitle=0.6, toprule=.15mm, colback=white, left=2mm, coltitle=black, bottomrule=.15mm, titlerule=0mm, title=\textcolor{quarto-callout-important-color}{\faExclamation}\hspace{0.5em}{Going beyond `agreement/disagreement'}, toptitle=1mm, colframe=quarto-callout-important-color-frame, arc=.35mm, rightrule=.15mm]

It is important to go beyond simple agreement or disagreement, and
interpret the data. The following is a `not so great' example:

\begin{itemize}
\tightlist
\item
  ``Our theme showed that postgraduate students felt a lot of pressure
  to succeed, but they did not agree on what ``success'' actually meant.
  This is similar to Skye (2023) who said that postgraduates think
  success means getting good grades, although some thought it meant
  being a good group member or answering questions in class. Therefore,
  our results match Skye's findings.''
\end{itemize}

This is a stronger example:

\begin{itemize}
\tightlist
\item
  ``Our theme showed that postgraduate students feel strong pressure to
  `succeed,' yet they hold different and sometimes conflicting ideas
  about what success actually means. This supports Skye's (2023) finding
  that, while many postgraduates focus on good grades, others define
  success through group contribution or speaking up in class. Our
  findings add to this by showing that it is even more complex, as
  students tend to hold multiple ideas of success at the same time,
  creating uncertainty about how to meet expectations. This shows that
  success is not a fixed standard but something students actively
  negotiate in response to academic norms and peer pressures. This
  emphasis on negotiation echoes Isla's (2025) findings that medical
  staff similarly construct their professional identities in relation to
  shifting expectations.''
\end{itemize}

In the second, the author goes into more depth and explains to the
reader what the current findings add. They also then explain this by
linking (through the observation that negotiation was key) to a relevant
finding from a different field.

\end{tcolorbox}

\textbf{Discuss implications}

\begin{itemize}
\tightlist
\item
  Explain any practical or theoretical implications of your findings.
  Remember to make these realistic and linked to your findings.
\item
  Try to make these concrete suggestions, rather than vague statements.
\item
  Remember to back these points up with evidence from the literature
  base.
\end{itemize}

\textbf{Reflect on limitations}

\begin{itemize}
\tightlist
\item
  In this length of report, we suggest picking a maximum of one or two
  limitations, so that you have the opportunity to develop some depth.
\item
  Avoid limitations that are more suited to quantitative studies

  \begin{itemize}
  \tightlist
  \item
    sample size
  \item
    lack of generalisability
  \end{itemize}
\item
  Think about limitations that are \emph{relevant} instead

  \begin{itemize}
  \tightlist
  \item
    The kinds of experiences or voices not represented
  \item
    What the research design prioritised or obscured
  \item
    How the method of data collection might have shaped interpretation
  \end{itemize}
\item
  As with the other sections, the limitations section is stronger when
  you clearly link to the evidence base.
\end{itemize}

\textbf{Suggest future research}

\begin{itemize}
\tightlist
\item
  Show what you would do differently next time and why. Propose ideas
  for studies that naturally follow from your findings.
\item
  Link to the literature base to support your points.
\end{itemize}

\begin{figure}[H]

{\centering \includegraphics[width=1\linewidth,height=\textheight,keepaspectratio]{images/discussion.png}

}

\caption{Narrow-to-broad structure of a discussion}

\end{figure}%

\begin{tcolorbox}[enhanced jigsaw, breakable, bottomtitle=1mm, leftrule=.75mm, opacityback=0, colbacktitle=quarto-callout-warning-color!10!white, opacitybacktitle=0.6, toprule=.15mm, colback=white, left=2mm, coltitle=black, bottomrule=.15mm, titlerule=0mm, title={Tips}, toptitle=1mm, colframe=quarto-callout-warning-color-frame, arc=.35mm, rightrule=.15mm]

\begin{itemize}
\tightlist
\item
  Link to evidence in support of the points you are making throughout
  the whole Discussion section (including limitations, implications and
  future directions).
\item
  Pick limitations that are relevant to qualitative research.
\item
  It's okay to bring in new literature into the Discussion, but try to
  also bring back key studies/theories that you introduced earlier on in
  the report.
\end{itemize}

\end{tcolorbox}

\subsubsection{Recap of word count for each section of the
report}\label{recap-of-word-count-for-each-section-of-the-report}

\begin{itemize}
\tightlist
\item
  \textbf{Abstract}: 100-150 words
\item
  \textbf{Introduction}: 700-800 words
\item
  \textbf{Method}: 350-450 words
\item
  \textbf{Analysis}: 750-900 words
\item
  \textbf{Discussion}: 750-850 words
\end{itemize}

We understand this can feel like a lot to fit into the report, but it's
a good chance to practise writing concisely. It's okay to go slightly
outside these guidelines, but try not to deviate too much.

A common mistake is making the abstract or method section longer than
recommended, which then affects the quality of the other sections.

If you're struggling with the word count, feel free to drop into office
hours. We can't give general feedback, but you're welcome to ask
specific questions about how you might condense your report.

\subsubsection{References}\label{references}

Citations in the main body of the text and references in the reference
list are structured in exactly the same way as in a quantitative report,
using APA 7th style.

\subsubsection{Appendices}\label{appendices}

The appendices can often hold a lot of information. In the qualitative
report we ask you to include the following in the appendices of the
report:

\begin{itemize}
\tightlist
\item
  \textbf{Required}: A list of the interviews you have chosen and a link
  to the original dataset (do not include the transcripts themselves)
\item
  \textbf{Required}:
  \href{https://gla-my.sharepoint.com/:w:/g/personal/ashley_robertson_glasgow_ac_uk/IQC1PIxhBlndR74S_nfeKHNlAW5uQmGOpoRXCpLHC2fI_fI?e=aunShJ}{A
  reflexivity diary} that you have developed over the course of
  completing your analysis. Please use the template linked here to
  complete your diary.
\item
  \textbf{Only if available}: You might have a list of the interview
  questions (if easily obtained from your transcripts or study
  documentation on UK data service)
\item
  There may be other appendices you choose to include. If in doubt,
  please ask and we will do our best to help.
\end{itemize}

\subsection{Submission 2: Step-by-step Evidence of the Analysis
Process}\label{submission-2-step-by-step-evidence-of-the-analysis-process}

This year, we are asking that all students in RM2 complete two separate
submissions for their qualitative report.

\begin{itemize}
\tightlist
\item
  Submission 1 = full report with cover sheet and appendices
\item
  Submission 2 = step-by-step evidence of your Reflexive Thematic
  Analysis
\end{itemize}

Submission 2 will be formally assessed \hyperref[report_ILOs]{through
the fourth ILO in `Evaluation'.}

The reasons for this requirement are three-fold:

\begin{itemize}
\item
  \textbf{To support your learning of rigorous qualitative research
  skills.}

  Reflexive Thematic Analysis (RTA) relies on transparency, reflexivity,
  and a clear audit trail of decisions. By documenting each step, you
  gain a deeper understanding of how qualitative analysis works and
  build skills that are essential for doing high-quality research.
\item
  \textbf{To ensure that your submitted work reflects your own analytic
  thinking.}

  We know that generative AI tools are now widely accessible, and while
  they can be useful for support, they cannot replace the reflective,
  interpretive, and creative thinking at the heart of RTA. Taking the
  process step-by-step helps you to demonstrate your genuine engagement
  with the data and protects the integrity of your work.
\item
  \textbf{To help you experience RTA as a process, not just a product.}

  Good RTA is not about producing a set of themes; it is about the
  journey you take to get there. Keeping a portfolio helps you slow
  down, think critically, and make your analytic decisions visible. This
  will lead to richer insights and a stronger final analysis.
\end{itemize}

In short, we hope that this requirement for keeping a transparent record
of your analytic journey will help you produce stronger work and deepen
your understanding of qualitative research.

\href{https://gla-my.sharepoint.com/:w:/g/personal/ashley_robertson_glasgow_ac_uk/EcwNeR1TcCpAi1inMtQgB3cBkucx3-cH3LO-Mk2fseBkOQ?e=LEfvVw}{The
template can be found here}

\textbf{Answers to some questions}

\begin{itemize}
\tightlist
\item
  \textbf{What advice do you have for completing this process?} Our main
  piece of advice is to work on the reflexivity diary (part of the
  appendices) and the `evidencing the process' template \emph{while you
  are working on the project}. You should be completing these steps
  anyway as part of the analysis, so it should be a case of logging the
  evidence as you go. If, however, you do not complete it as you go
  along, this would make it very difficult to include properly at the
  last minute.
\item
  \textbf{Can I use NVivo?} For this assignment, no. We plan to develop
  guidance that considers how students can best use NVivo in the age of
  GenAI, but this will take some time to develop properly.
\item
  \textbf{I developed three themes, but can only report two in the qual
  report itself: What do I do?} It probably makes sense to include all
  of the information up until the end of Phase 4, as often the themes
  are intermingled together. However, for the Phase 5 part (naming and
  defining the themes), it is up to you: you can include all of your
  themes or you can include only those that are in the report.
\end{itemize}

Below, we have prepared some guidance for each section to help you.

\subsubsection{Phase 1: Familiarisation}\label{phase-1-familiarisation}

\begin{tcolorbox}[enhanced jigsaw, breakable, bottomtitle=1mm, leftrule=.75mm, opacityback=0, colbacktitle=quarto-callout-important-color!10!white, opacitybacktitle=0.6, toprule=.15mm, colback=white, left=2mm, coltitle=black, bottomrule=.15mm, titlerule=0mm, title=\textcolor{quarto-callout-important-color}{\faExclamation}\hspace{0.5em}{Phase 1: What is asked for?}, toptitle=1mm, colframe=quarto-callout-important-color-frame, arc=.35mm, rightrule=.15mm]

In the template, we ask for your familiarisation notes, when you read
through each transcript.

\end{tcolorbox}

\textbf{Guidelines for Phase 1}

Familiarity notes are brief reflections written while reading and
re-reading the data. They are often messy and very personal and are not
`right' or `wrong'.

They tend to capture:

\begin{itemize}
\tightlist
\item
  First impressions
\item
  What stands out
\item
  What seems confusing or unusual
\item
  Emotional reactions
\item
  Questions or curiosities
\item
  Early ideas about what the participant might be ``doing'' with their
  words
\item
  Noticing of tone, contradictions, or patterns
\end{itemize}

Some things you might note down:

\begin{itemize}
\tightlist
\item
  What immediately grabs your attention?
\item
  How do you feel in response to what you're reading?
\item
  What do you think the participant is trying to express? Why do you
  think this?
\item
  Are any ideas or experiences repeated?
\item
  Do they say things that don't fit together?
\item
  What are you curious about? What would you ask them if you could?
\item
  Anything about their background that shapes how you read this?
\item
  How do your assumptions/interests shape what seems important?
\end{itemize}

\begin{tcolorbox}[enhanced jigsaw, breakable, bottomtitle=1mm, leftrule=.75mm, opacityback=0, colbacktitle=quarto-callout-warning-color!10!white, opacitybacktitle=0.6, toprule=.15mm, colback=white, left=2mm, coltitle=black, bottomrule=.15mm, titlerule=0mm, title={Example}, toptitle=1mm, colframe=quarto-callout-warning-color-frame, arc=.35mm, rightrule=.15mm]

This is a short example of a familiarisation note, to give you a sense
of what might be expected.

\emph{``My first impression is that Participant 1 feels overwhelmed but
also quite capable, which creates an interesting tension. Words like
`pressure', `everyone watching', and `expectations' make me think
they're under a lot of stress. I notice I feel a bit anxious reading
this - maybe because I've had similar experiences. I'm also interested
in their comment about `putting on a front' at work. It makes me wonder
why they feel they can't be their real self in that situation.''}

This familiarity note captures initial impressions, points to specific
language that shaped those impressions, and reflects openly on personal
reactions while reading. It also highlights moments that stood out and
raises gentle questions, showing curiosity rather than early analysis.

\end{tcolorbox}

\subsubsection{Phase 2: Developing
Codes}\label{phase-2-developing-codes}

\begin{tcolorbox}[enhanced jigsaw, breakable, bottomtitle=1mm, leftrule=.75mm, opacityback=0, colbacktitle=quarto-callout-important-color!10!white, opacitybacktitle=0.6, toprule=.15mm, colback=white, left=2mm, coltitle=black, bottomrule=.15mm, titlerule=0mm, title=\textcolor{quarto-callout-important-color}{\faExclamation}\hspace{0.5em}{Phase 2: What is asked for?}, toptitle=1mm, colframe=quarto-callout-important-color-frame, arc=.35mm, rightrule=.15mm]

In the template, we ask for you to provide a table with all of your
codes in it.

\end{tcolorbox}

\textbf{Guidelines for Phase 2}

When coding, think about the following:

\begin{itemize}
\tightlist
\item
  Work through your data section by section, thinking about what feels
  meaningful, interesting or like there is a pattern
\item
  Create short labels (codes) that capture the essence of what each
  highlighted section is \emph{doing}, \emph{saying}, or
  \emph{representing}
\item
  Remember that you can be flexible with the coding process. You can
  create as many codes as you need, and you can merge or rename them
  later. This allows you to remain responsive as your coding develops
  over the dataset.
\item
  Try to keep your coding active and interpretive, focusing on meaning
  instead of counting how often something appears.
\item
  If you are unsure whether to code something or not, err on the side of
  caution and include it. This will help you avoid overlooking sections
  that seem contradictory, subtle, or hard to interpret
\item
  Keep a note of your reflections in your reflexivity diary as you code,
  which will help in later phases.
\end{itemize}

\subsubsection{Phase 3: Generating initial
themes}\label{phase-3-generating-initial-themes}

\begin{tcolorbox}[enhanced jigsaw, breakable, bottomtitle=1mm, leftrule=.75mm, opacityback=0, colbacktitle=quarto-callout-important-color!10!white, opacitybacktitle=0.6, toprule=.15mm, colback=white, left=2mm, coltitle=black, bottomrule=.15mm, titlerule=0mm, title=\textcolor{quarto-callout-important-color}{\faExclamation}\hspace{0.5em}{Phase 3: What is asked for?}, toptitle=1mm, colframe=quarto-callout-important-color-frame, arc=.35mm, rightrule=.15mm]

In the template, we ask you to a) tell us about your initial themes
and/or subthemes, b) tell us the process of how you developed these, and
c) tell us about the initial patterns you observed.

\end{tcolorbox}

In this section, you start to group together your codes into themes.
Some parts might come together nicely and others might feel more
difficult to put into groups. This is part of the process and nothing to
worry about!

\textbf{Guidelines for Phase 3}

When forming your initial themes, think about the following:

\begin{itemize}
\tightlist
\item
  Look across all your codes and identify broader meaning-based
  patterns.
\item
  Group together codes that seem related, connected, or part of a shared
  story.
\item
  Notice similarities, differences, tensions, and recurring ideas across
  your coded data. Sometimes you can have themes that look contradictory
  on the surface, but there is shared meaning underneath.
\item
  Create rough theme names that capture the central idea or organising
  concept in each cluster.
\item
  Write a short description of what each initial theme is about and why
  it matters.
\item
  Write reflections about how you developed the theme (e.g., repeated
  emotions, shared experiences, contradictions you observed).
\item
  Remember these are \textbf{initial} themes - this is not the final
  version and they can change, merge, or be dropped in later phases.
\end{itemize}

\subsubsection{Phase 4: Developing and Reviewing
Themes}\label{phase-4-developing-and-reviewing-themes}

\begin{tcolorbox}[enhanced jigsaw, breakable, bottomtitle=1mm, leftrule=.75mm, opacityback=0, colbacktitle=quarto-callout-important-color!10!white, opacitybacktitle=0.6, toprule=.15mm, colback=white, left=2mm, coltitle=black, bottomrule=.15mm, titlerule=0mm, title=\textcolor{quarto-callout-important-color}{\faExclamation}\hspace{0.5em}{Phase 4: What is asked for?}, toptitle=1mm, colframe=quarto-callout-important-color-frame, arc=.35mm, rightrule=.15mm]

In the template, we ask for you to provide details of how your themes
changed from the initial themes you identified in Phase 3.

\end{tcolorbox}

In this section, this is where you really think about what your themes
\textbf{mean}, and how they relate to each other. There is often some
back-and-forth during this part of the process, and students often end
up going back to Phase 3 (or even 2).

\textbf{Guidelines for Phase 4}

\begin{itemize}
\tightlist
\item
  Revisit the coded extracts for each theme and check whether they
  \emph{truly} fit together and reflect a coherent idea.
\item
  Look for any extracts that feel out of place, unclear, or better
  suited to a different theme --- or that suggest a new theme.
\item
  Refine each theme's focus by clarifying its central organising
  concept: what is the core meaning holding this theme together?
\item
  Check whether themes are too broad, too narrow, overlapping, or
  repetitive, and adjust by splitting, combining, or renaming.
\item
  Review themes against the \emph{full dataset} to ensure they capture
  important patterns across the data, not just isolated points.
\item
  Consider whether any parts of the dataset have been overlooked or
  require revisiting earlier phases of coding or theme development.
\item
  Make notes about how your themes changed (e.g., merged, refined,
  expanded) and why these changes improved the analysis.
\item
  Ensure each theme tells a meaningful, distinct part of the overall
  story you are constructing about the data.
\end{itemize}

\begin{tcolorbox}[enhanced jigsaw, breakable, bottomtitle=1mm, leftrule=.75mm, opacityback=0, colbacktitle=quarto-callout-warning-color!10!white, opacitybacktitle=0.6, toprule=.15mm, colback=white, left=2mm, coltitle=black, bottomrule=.15mm, titlerule=0mm, title={Example}, toptitle=1mm, colframe=quarto-callout-warning-color-frame, arc=.35mm, rightrule=.15mm]

This is a short example of how a researcher might change themes from
Phase 3 to Phase 4, to give you a sense of what might be expected.

\textbf{Initial Theme (Phase 3):} ``Feeling Under Pressure''

\begin{itemize}
\tightlist
\item
  This theme was developed from a cluster of codes including ``pressure
  to perform'', ``fear of being judged'', and ``trying to keep up''.
\item
  It seemed that participants were describing feeling stressed and
  overwhelmed by expectations.
\end{itemize}

\textbf{Refinement Process (Phase 4)}

\begin{itemize}
\tightlist
\item
  When reviewing coded extracts, I noticed that some excerpts were more
  about \emph{external expectations} (e.g., others judging them), while
  others were about \emph{internal expectations} (e.g., pushing
  themselves).
\item
  This suggested that this initial theme was too broad, and actually
  captured two different patterns, so I split the theme into two themes.
\item
  I then changed the names to better reflect the differences between
  these themes.
\end{itemize}

\textbf{Refined Themes (Phase 4):} Theme 1: ``Navigating Others'
Judgements'' and Theme 2: ``Holding Myself to High Standards''

\begin{enumerate}
\def\labelenumi{\arabic{enumi}.}
\tightlist
\item
  ``Navigating Others' Judgements''
\end{enumerate}

\begin{itemize}
\tightlist
\item
  This theme captures experiences of being watched, evaluated, or
  compared to others, reflecting a concern about how others perceive the
  participant.
\end{itemize}

\begin{enumerate}
\def\labelenumi{\arabic{enumi}.}
\setcounter{enumi}{1}
\tightlist
\item
  ``Holding Myself to High Standards''
\end{enumerate}

\begin{itemize}
\tightlist
\item
  This theme reflects participants' self-imposed pressure,
  perfectionism, and high standards.
\end{itemize}

\end{tcolorbox}

\subsubsection{Phase 5: Refining, Defining and Naming
Themes}\label{phase-5-refining-defining-and-naming-themes}

\begin{tcolorbox}[enhanced jigsaw, breakable, bottomtitle=1mm, leftrule=.75mm, opacityback=0, colbacktitle=quarto-callout-important-color!10!white, opacitybacktitle=0.6, toprule=.15mm, colback=white, left=2mm, coltitle=black, bottomrule=.15mm, titlerule=0mm, title=\textcolor{quarto-callout-important-color}{\faExclamation}\hspace{0.5em}{Phase 5: What is asked for?}, toptitle=1mm, colframe=quarto-callout-important-color-frame, arc=.35mm, rightrule=.15mm]

In the template, we ask to a) tell us how you came up with your theme
names and b) provide a definition for your theme/s.

\end{tcolorbox}

In this section, this is where you aim to clearly articulate what each
theme is about, define its boundaries, and give it a meaningful name
that captures its central organising concept.

\textbf{Guidelines}

\begin{itemize}
\tightlist
\item
  Revisit each theme and identify its central organising concept (the
  key idea that holds the theme together).
\item
  Write a clear explanation of what each theme is about, focusing on the
  shared meaning across the extracts, not just the topic.
\item
  Know the boundaries of the theme - what belongs in it and what does
  not? One suggestion is to imagine you are telling someone who doesn't
  know your research about your theme. What would you include in such a
  summary?
\item
  Check that each theme is distinct from the others and not overlapping,
  repetitive, or too similar in focus.
\item
  Refine the name of the theme so it captures its core meaning in a
  concise, engaging, and conceptually clear way. Avoid one word theme
  names or names that are very descriptive.
\item
  Write a brief description of how each theme contributes to your
  overall analysis.
\item
  Ensure each theme connects back to your research question and supports
  the analytical narrative you are building.
\end{itemize}

\subsubsection{Coded transcripts}\label{coded-transcripts}

Here, we ask that you provide all of your transcripts that have been
coded (we advise two or three, with a max of four if doing the couples
dataset). If you have coded the transcripts multiple times, please only
include one set of transcripts.

\begin{tcolorbox}[enhanced jigsaw, breakable, bottomtitle=1mm, leftrule=.75mm, opacityback=0, colbacktitle=quarto-callout-important-color!10!white, opacitybacktitle=0.6, toprule=.15mm, colback=white, left=2mm, coltitle=black, bottomrule=.15mm, titlerule=0mm, title=\textcolor{quarto-callout-important-color}{\faExclamation}\hspace{0.5em}{Coded transcripts: What is asked for?}, toptitle=1mm, colframe=quarto-callout-important-color-frame, arc=.35mm, rightrule=.15mm]

In the template, we ask for you to provide all of your coded
transcripts. These should either be commented in Word or scans of
physical copies.

\end{tcolorbox}

\chapter{Frequently Asked Questions}\label{ch13_faqs}

This page was last updated 5th March 2026

\section{General Project questions}\label{general-project-questions}

\textbf{How do the portfolio and qualitative report work together -
should we focus on one before looking at the other?} The qualitative
portfolio is designed to support you in developing skills that would be
important for actually running qualitative studies (e.g.~reflexivity,
developing questions, awareness of ethics, awareness of good/bad
practice for data collection). The report is designed to support you in
developing your data analysis skills and then writing up this
information into a report.

Re. the timescales - we would advise that you start working (or at least
thinking about) the qualitative report before your qualitative portfolio
is handed in. However, we realise everyone is different and prefer to
work in different ways. We have developed a timeline for the summative
assessment with suggested milestones, although you should revise these
deadlines to find something that works for you.

\section{Qualitative portfolio
questions}\label{qualitative-portfolio-questions}

\subsection{General portfolio
questions}\label{general-portfolio-questions}

\textbf{Where can I find information about the qualitative portfolio?}
You can find information about the portfolio in the template, the
guidelines, the AIS and the video Wil recorded. You can access the
materials in \hyperref[ch10_group_project]{chapter 10} of this book.

\textbf{What is the word count for each section?} This is in the
template, but it can be broken down as follows: Activity 1a (no word
limit for the six questions), Activity 1b (250 words), Activity 2 (250
words), Activity 3 (500 words).

\textbf{In the answers, should we use or avoid acronyms such as BPS and
GDPR?} Usually in your writing you should avoid unnecessary
abbreviations unless they are really commonly used, as this makes it
harder for the reader to follow your writing. A good general rule of
thumb is not to come up with any abbreviations yourself. These two are
very common for the context of your assignment so you can use them (this
will also help with the word count which is very limited here).

\textbf{Do we need to set out an introduction and concluding summary in
our responses to each of the sections?} You don't need this - you can
dive straight in :)

\textbf{Is it okay for us to take what a study says at face value
(i.e.~the study says x, which backs up the point)? Or do we have to be
critical about the study?} Yes, this is fine - we expect you to use the
studies to support the points you are making, rather than critically
interrogating how well the study supports the claim. It would not really
be possible to do the latter within the word count. This assessment is a
bit different to when you are writing a report, for example, where you'd
want to show how you are building on the existing literature (to help
support your rationale)

\textbf{I can't figure out how to cite the BPS code; examples online are
both ``BPS,2021'' and ``Oates et al., 2021''} Cite it as BPS.

\textbf{Is it okay to answer the questions or do we have to make a
unique contribution (as you do in an essay or a report)?} Just doing
what you are asked is fine. This is not the same type of assessment as a
report or essay as it is focusing on skills needed to run a project, so
it will be assessing different things.

\textbf{Should I delete the included instructions from the template
document so it's just the headings for each section?} Either keeping
them or deleting them is fine - whichever you prefer.

\textbf{Can I use guidelines and and strategies for qualitative research
from fields other than psychology?} Yes, perfectly fine to use these
guidelines - you don't have to stick to Psychology. Just make sure that
the guidelines apply to the type of research methodology/RQ you are
discussing.

\textbf{Is slight variability in word count between the sections allowed
as long as the final product is less than 1000?} Yes, there is slight
variability allowed - it is overall 1000 but you can change the
distribution that is suggested slightly (however, try not to change by
more than 30-40 words so you provide the right amount of detail for each
section).

\subsection{Task 1: Qualitative questions \&
reflection}\label{task-1-qualitative-questions-reflection}

\textbf{Do the qualitative questions we submit also need referencing?}
No, you don't need citations for section 1A. But make sure to support
your reflection with evidence in 1B.

\textbf{Is it acceptable to refer to `the research question' without
having to write it out in task 1B?} Yes - you don't need to repeat it as
the word count is very limited here.

\textbf{To what extent can we refine the sample? The RQ refers to
`students', but could we specify that further by limiting it to only
students who have / not sought out mental health support?} Yes, if you
want to limit it further, this is fine - we kept it somewhat flexible on
purpose so you have a bit of freedom with this.

\textbf{What sort of evidence is relevant in a reflection? Is it related
to the topic (i.e.~of mental health) or is it more about the anticipated
problems / other researchers protocol in similar research? Or a bit of
both? Should we avoid trying to justify our choices in the reflection,
or is that a good thing to do?} I think it would make sense if it's more
related to methods rather than the topic choice - would be quite good to
justify your choices in the reflection.

\textbf{Would it be a good idea to include a warm up question/ice
breaker and a closing question in our six questions?} For the purposes
of this assessment, we would suggest just sticking to the six main
questions (i.e.~not using up two of these for a warm up/exit question).
This is because we ask for a limited number, and that way we can keep
things consistent across the whole class. One thing you might want to
consider is your placement of questions (i.e.~the order in which you
would ask them).

\textbf{What is the difference between an interview and a qualitative
survey?} Interview is where the researcher is one-to-one with a
participant and asks them questions. Qualitative survey is a qualitative
questionnaire that is often sent to participants online and they fill it
out in writing.

\textbf{How does (evidence based) reflection differ (if at all) from
(evidence based) writing that we have been doing so far for our
assignments?} It doesn't really - we want you to look at the evidence
and back up your ideas with relevant sources. The portfolio materials
have some probes for your reflection. In the ethics task you should also
use the BPS code of ethics and refer to that as well.

\textbf{For activity 1, do we would need to pick a specific topic area
(from the evidence base) to explore in more depth in order to frame our
questions?} The overall topic and research question is given to you in
the materials, and then you need to come up with focus
group/interview/qualitative survey questions that you would ask the
participants to investigate the research question

\textbf{Should we mark some of the questions as priority with a star, to
note the ones that we would ask first if time was running short in the
focus group?} You can asterisk the ones you deem most important if you
want to, but there is no requirement to do so either (so basically,
whichever you prefer)

\textbf{Can we include prompts underneath some of the questions?} You
can if you want, but we ask you to stick to a maximum of one prompt per
question if you do this (and that it's really clear for the reader what
the main question is and what the prompt is)

\textbf{Is there scope to include visual prompts in our questions?}
Theoretically yes, although we've never seen this asked before so
wouldn't know what a visual prompt would look like. If you want to
include this, it might be worth contacting Ashley / Wil to show them an
example.

\textbf{What should we focus our reflection on? About half of the word
count in our reflection is currently taken up by an evidence-based
discussion of why we opted for the FG format over the interview as this
really guided how we constructed the questions thereafter. I just want
to know if such a focus is correct or if you're looking for discussion
to centre more around a specific question or two that we chose?}

There is no one correct way to structure the reflection. We are
interested in how you found the experience and why you made the
decisions that you did in constructing your list of questions. We often
see reflections on why people chose a particular approach for their
questions too (e.g.~focus group over interview/survey).

\textbf{You have advised that we choose 1 or 2 issues and go into their
depth when answering the questions on ethics and the terrible FG. What
do you recommend we do for the reflections on the questions in task 1?
Do we for instance speak about why we chose a specific approach over the
other? Why we chose these specific questions? Challenges we faced etc.
Or do we combine everything? Now, we have 250 words of a general
reflection on FGs, guideline designing on the topic under discussion,
challenges we faced, and what we couldn't get with citations of course.
I'm wondering if we should be more focused and try to answer only 1
point?} Basically, we don't specify what you can and can't reflect on,
so as long as it's relevant (which this is), then you have some
flexibility :) It is a very open question and we are anticipating lots
of different ways of answering this one in particular. The main thing is
to follow the guidelines that have been given and to look at the ILOs,
to see how well you think your response is meeting them

\textbf{Is it appropriate to briefly describe our own expertise and
experiences, and how they've informed our decisions? Would this count as
reflexivity and would you like to see evidence of reflexivity? Or would
you rather we didn't talk about ourselves directly?} Yes, this would be
completely appropriate to discuss as part of reflexivity. It's fine to
discuss yourselves directly (but it is not the only way to do this
task). You would be reflecting on why you made the choices you did and
would be demonstrating what experiences you have that went into those.
Be mindful to make this evidence-based too!

\textbf{We have included a short rationale on why we chose a focus group
which is placed before the six questions. Will this short paragraph be
included within the word count for the whole portfolio or will it not
count as part of the 1000 words?} A rationale (if you choose to include
it) should be in 1b. The only thing not included in the word count
within the activities themselves is the actual questions of Activity 1a

\textbf{Do prompts/probes have to be phrased as a question or can they
be more of a statement?} They will usually be framed as a question but
it is fine to write them as a statement for this assignment if you
prefer.

\textbf{Do we need to define terms such as open-ended questions?} No,
there isn't space for this within the word count for this particular
assignment

\textbf{If we're including our group member's experiences for
reflexivity, should we name them, or just refer to them as `group
member'?} You can either use initials or ``group member'', either one
would be fine

\textbf{Within the questions we are setting for our group project -
would it be appropriate to name ``University of Glasgow'' for instance
or would it be best to keep to just ``University'' in general?} We would
suggest just sticking to University, rather than specifying the
University of Glasgow in particular.

\textbf{is it OK to assume that the participants work in the same
university as its important for a point we've made in the moderating?}
Yes, this is fine :)

\subsection{Task 2: Ethics}\label{task-2-ethics-1}

\textbf{Should we connect our answer to Question 2 to what we did for
Q1?} No, there is no expectation that Activity 2 is linked to Activity 1
(although Activites 1a and 1b need to be linked to each other).

\textbf{Can you tell us where people tend to go most wrong in this
particular question?} Sometimes, groups will focus more on ethical
issues that are relevant for quantitative research, or research in
general, rather than specifically thinking about qualitative research.
This is something that is best avoided as it doesn't answer the question
that is set.

\subsection{Task 3: Focus group gone
wrong}\label{task-3-focus-group-gone-wrong-1}

\textbf{For focus group, do we have to back the issue with specific
observations from the video? E.g @ 4.05min X did what, which is an issue
because\ldots{}} No, this would likely take up too many words.

\textbf{Is there a preference for either moderator or facilitator?} No,
these are interchangeable and either is fine

\textbf{Are we to only quote peer reviewed literature in support of our
argument, or can we quote for example from the Braun and Clarke book?}
No, it's fine to cite Braun \& Clarke and there might also be
guides/opinion pieces etc. - these would also be fine to use in support.
Also, it's fine to use older evidence too - we do realise there's
unlikely to be lots of 2022/2023 papers out there on how to do a focus
group! We have also included some additional papers on the reading list
(under Thinking about the focus group) to give a bit of a start
(although please note we haven't had time to read these\ldots we've
skimmed through)

\textbf{We have identified the point we want to talk about. For example,
it is about the moderator's part. Should we focus on the moderator's
issue in general or pick a couple of specific examples of bad moderating
and write about them in depth? Because the word count is limited, it
isn't apparent what strategy to pursue.} Either is fine to be honest.
You could pick moderating as one issue and then another one. Or you
could focus on two separate parts of moderation. My main suggestion
would be to pick an approach that allows you to go into depth on the
issue, explaining why it is a problem, what could be improved etc.

\textbf{I'm struggling with how to approach the focus group task- do you
have any examples excellent or `A' standard work?} We don't provide
exemplars, because this assessment is specific so it would be really
hard to do that without giving out the answers. In terms of what the
focus group task is asking, please make sure you read the ILOs for this
assessment that can be found on the AIS document. The only ILO specific
for the FG task is:

Correctly identify TWO or THREE issues present in the focus group,
explain why they are problematic, and discuss possible improvements
(\#3: Focus group), which means that you should identify what is wrong
in how the focus group is conducted (this is mostly based on what the
moderator is/isn't doing), explain why these things are not good when
you are conducting a FG, and explain how they could approach this
instead.

To meet the other ILOs for this task, you should make sure you write
clearly and concisely and demonstrate evaluation in your answer (this is
an ILO for all your assessments you have done to date and means that you
need to use evidence to build your arguments and go beyond description).
For this task evaluation should include you explaining why your chosen
issues are bad practice when conducting a focus group, and citing some
evidence (this will mostly be methodological in nature) to back up your
answer.

\textbf{For the focus group task, do we need to include quotes from the
transcript when referring to the issues?} No, we would advise that you
don't do this, as it will take up a lot of the (very limited!) word
count. Instead of quoting, just say what the participant/facilitator did
that you think explains the issue you have identified.

\textbf{When I'm evaluating a study, do you want me to criticise it? Do
you want me to talk about the methodology?} For this assessment, we
aren't really looking for you to criticise the studies that you are
using - for example, for the focus group task, it's more about
explaining why it's a bad thing that the moderator is doing and then
using some evidence to back that up (e.g.~some guidance for how to run a
focus group or if there's a paper saying what you might want to do to
make your participants feel more comfortable). Here, it's not about
saying `here's this study, it had x participants' - that is more for the
report, where you are building a rationale.

\textbf{Would a good structure be: x happened, x is bad for y reason, z
should be done instead (with studies to back up)?} There are different
ways that you can structure it but the above example would be fine. The
main thing is that you are covering what we ask of you in the guidance
and the ILOs.

\textbf{Should we include timestamps of when something happened in the
video?} No, this is not necessary to include - just telling us what
you've identified will be fine

\textbf{A lot of the literature is pretty old. What should I do?} We are
aware that some of the literature is older, and it's completely fine to
use this. It's not the same type of assessment as essays/reports etc.,
and lot of the guidance is from the 90s or so.

\textbf{Do we refer to the participants in the terrible video by their
names eg: Ashley, Gaby? Or we call them participant 1, 2, 3, or the
moderator (or Wilhelmina) did this with a participant, and that with
another participant?} I think generally I would say `a participant' etc.
as it depersonalises it. However, you do have the names there as they
were on the video and - as we haven't given any guidance specifically on
this - it won't affect the marking in any way if you use our names.

\section{Qualitative report
questions}\label{qualitative-report-questions}

\subsection{General report questions}\label{general-report-questions}

\textbf{Will there be detailed guidance about each section of the qual
report?} Yes, we have a guidance in this book in the qualitative report
section in chapter 12, which will cover the key points to cover in the
report.

\textbf{Is our report written in a specific style, e.g.~IPA, thematic
analysis etc?} You will use thematic analysis to analyse your data. We
covered this in detail in the lectures in semester 1, and we will also
have activities to support this in semester 2.

\textbf{How many interviews should we select for our report?} We ask
that you choose a maximum of three interviews to analyse (less than this
is fine if you have enough data within your chosen interviews). If you
are working with the relationships dataset, you can do two couples (four
interviews) if you wish to.

\textbf{I'm not sure what all the different tasks are for the report -
where can we find this information?} We haven't shared an explicit task
list for the report and one reason for that is because RM2 is
purposefully designed to get you to work a bit more independently in
preparation for the dissertation.

Key stages are as follows:

\begin{itemize}
\tightlist
\item
  Familiarise yourself with what is expected for the report assignment
  have a look at the data and decide which dataset you want to work with
\item
  Look at the data in more detail - which interviews will you select?
\item
  Look at relevant literature and start to think about your RQ think
  about the rationale for the study
\item
  Conduct your Thematic Analysis
\item
  Write up the study as a full research report (which will be similar to
  your RM1 report).
\end{itemize}

\textbf{For the report, I am unclear on the criteria for choosing
transcripts. I have an approximate idea of a RQ and looking at the
transcripts they could be broadly grouped into 2 sets relating to how
the participant responds to the questions. Is the intention to choose
transcripts that have a consistent thread or that cover a diverse range
of views? With only 3 transcripts I'm not sure if it should be based on
commonalities or if this is kind of cherry picking?} I would say it's
about choosing transcripts that you will be able to analyse and are
relevant to the specific research question that you have. So, it would
make sense to pick transcripts that are relevant (i.e.~they cover the
issue you are interested in). They may or may not have the same
perspective - I think either would be okay to be honest. Sometimes, it
can be useful to have different views as you can then think about why
(which can be interesting in terms of evaluation of your themes) and
it's also fine to have more consistent views too.

\textbf{Can I do a different analysis and not TA?} For the RM2 project
we unfortunately only have assessment guidance and support for doing a
thematic analysis (Braun \&Clarke's version) and we currently are unable
to expand that to different type of analyses - this is because it's an
introductory qualitative course so the same way as in RM1 you were asked
to do a specific analysis on the data, we would like you to stick with
TA for the individual report.

\textbf{What voice should I write in? Should my reflexivity be written
in first person?} APA style encourages the use of active voice instead
of passive voice (``I did'' instead of ``it was done''). You can use
either but make sure to be consistent and don't make your sentences
overly complicated. For reflexivity, you should definitely use ``I'',
because it is about your own reflection in regards to the RQ and the
data

\textbf{Should I be writing my report as though I designed and conducted
the interviews etc myself?} You should not write it as if you collected
the data - you should write it as a secondary data analysis.

\textbf{I've found published papers that are using the same dataset and
I'm worried that my research is the exact same/doesn't add anything new?
I feel that I should mention these original reports when I'm setting my
study in the wider literature, but I'm not sure how then to argue that
my study is adding anything new if I acknowledge these reports - what
should I do?} The chances of your RQ and all other parameters of the
study being identical are not very high; if your RQ is worded exactly
the same as a published paper, you may wish to tweak it a little bit to
focus on a more specific aspect of the data, so you can explain in the
introduction how it adds to existing literature. Also look out for what
kind of analysis they did - then you can discuss why TA might be better
suited for your RQ/the type of data you have. What subsample of
interviews did they use and is yours different? If so, why? It doesn't
need to be an entirely different thing to what is already published;
good science is built in incremental steps, so think about what are the
things that are different in your study to the published papers, and how
that adds to what we already know

\textbf{If my report is based on an understudied population with a
limited research base, should it be introduced as exploratory research?}
It could be but it doesn't have to be. It's a bit different in qual as
you don't have confirmation vs.~exploratory in the same was as with
quant. It is likely to be exploratory by nature though, as you are
exploring people's views

\subsection{Research questions}\label{research-questions}

\textbf{Do you have any tips for how to arrive at your research question
for the qualitative report?} This book has a section on developing RQs -
work through those activities first. I would probably have a look over
some of the datasets and decide what you'd like to focus on in terms of
topic. Then I'd think about what kind of information you get from them
(e.g.~is it people's opinions/experiences/beliefs etc.) and what you'd
be interested in looking at specifically.

\textbf{Should our questions be broad enough so we can use the whole
transcripts or should we narrow them to focus on specfic areas? I think
an obvious example to illustrate what I mean would be on the Life
limiting transcripts and then having a question of something like
``experiences of animals on wellbeing etc'' but that would then heavily
focus the content of the analysis on one of the sub areas of the
transcript and a lot (but not all) might be irrelevant?} Perfectly fine
to narrow them to focus on specific areas, but with this I would say
make sure that you have enough data about your specific RQ to be able to
analyse it. You don't have to use the whole transcript by any means -
completely fine to use subsections. It's possible to revise your RQ in
thematic analysis, so - if you were unsure - you could start off broader
and then could make it a bit more specific down the line. There is the
caveat that this would affect your intro, as you build towards why
you've chosen this research question in your rationale.

\textbf{Can we post more than one RQ on the padlet if we're not sure
what route to go down?} You can post a maximum of two RQs on the padlet
to get feedback on them. If you have more, please drop into office hours
with these.

\textbf{Does the research question have to be tailored towards one of
the questions asked in the interview or can it be broader?} It's
probably better to be broader - generally, themes will tend to be
broader than a single question that is asked and it might be a little on
the specific side

\subsection{Introduction \& Discussion}\label{introduction-discussion}

\textbf{Should we include a mini paragraph at the start of the
introduction summarising the research question at the end?} We would
recommend to avoid it for this specific assignment, as this might make
it difficult for you to have a broad-to-narrow structure and could cause
unnecessary repetition.

\textbf{In the discussion section, can we refer to the data (e.g.~the
participants) or if that is too detailed for this section?}

\textbf{Two of the interviewees I have picked agree on something whilst
the other disagrees so was going to specifically refer to this. Or
should I be picking up the broader theme that came out in the discussion
section and don't need to reference the specific point I am referring
to?} In the discussion you should focus on the themes and discuss them
in relation to existing literature/future literature instead of focusing
on individual comments from your participants (these you should analyse
in the analysis section). If you think back to quant, you are not
talking about the specific numbers in the discussion but instead focus
on what the findings mean and how they build on existing literature.
Sometimes in quali literature you see papers that have a sort of
combined analysis and discussion section, in which case they might
combine individual quotes and existing literature, but for this
assignment we are asking you to separate them to two different sections
for clarity and to make it a bit simpler for you to structure.

\textbf{Do limitations/future directions need to be separate sections?}
They don't need to be two separate paragraphs - often what is a
limitation to a study can logically be addressed in the next study.

\subsection{Method}\label{method}

\textbf{How to choose my theoretical framework and where does it go?}

\textbf{I take it that we are meant to indicate if we are using a
constructionist or realist/experientialist framework for the analysis.
Are we free to choose this or are there any recommendations? I would
assume that focusing on semantic themes is more doable given the limited
word count? Also, where would we state this? In the methods section or
at the start of the analysis?} It would be good to explain which
theoretical perspective you are coming from, yes. It's important to take
the word limit into account, so it could well be the case that focusing
at a semantic level will be more doable - the key thing in your analysis
is that you try to go beyond description, and develop a narrative. If
you include details about the framework, this should be done in the
Methods section, rather than at the start of the analysis.

\textbf{Is there an easy guide to working out what theoretical approach
you are using? I understand the difference between realism, social
constructionism etc in theory but I can't work out how to apply it to my
work. I kind of think it could be any of them so I think I'm missing
something!} It's more about thinking about your own perspective (rather
than the dataset), as you are correct - any of the different approaches
could be applied to the dataset. You could focus on what type of
knowledge you are looking for, or instead, think about
epistemology/ontology (or both!). It's sometimes easier to think about
the type of knowledge you want to obtain - if it is knowing more about
people's experiences, and how they describe their world, then that would
be about seeking out phenomenological knowledge. If it is more in terms
of uncovering something concrete that happens in the world (i.e.~a
truth) then it would be more realist that you are looking for. If you
have questions about the theoretical approach specifically, feel free to
pop into any of Ashley's office hours for a chat

\textbf{Also can we draw upon our personal experiences in the report?}

\textbf{I'm looking at the LGBT+ dataset and also work in an NHS trust;
I've seen some transphobic actions happen at work and was wondering if
this is appropriate to talk about in the report if appropriate for the
RQ} Yes, there is a section for reflexivity in the methods section,
where you should reflect on your experiences and how you felt about the
topic/your experiences. You can also expand on it a bit in the
discussion if you want.

\textbf{How much of the reflexive diary can I include in the method
report section?} You can use your reflexive diary (included as an
appendix) to help you write the reflexivity section, but don't directly
copy paste it in the main body of your report - the reflexivity section
under methods is much shorter than your full diary.

\textbf{Are there good examples we can have a look at for reference in
terms of reflexivity?}

\textbf{I'm looking through the list of examples studies you have
included using thematic analysis, and I don't think they include any
form of reflexivity at least not in the method section, or am I missing
something? Or would this be done during analysis?} Here are some
resources to help, but they are not so much examples of reflexivity
sections within papers and - instead - are more about giving guidance.
Here are some links to start you off:

\begin{itemize}
\tightlist
\item
  Barrett et al (2020): How to\ldots be reflexive when conducting
  qualitative research
\item
  McLaughlan et al (2012): Phenomenological analysis of patient
  experiences of medical student teaching encounters (they embedded
  reflexivity within the `Analysis' section of the methods
\item
  Blog by Jacqui Burne:
  https://nzareblog.wordpress.com/2017/11/28/self-reflexivity/
\item
  Byrne (2022) A worked example of Braun and Clarke's approach to
  reflexive thematic analysis
\item
  Examples of reflective commentary from online resources from
  'Reflexive Thematic
\item
  Analysis' (2022) by Braun \& Clarke:
  https://study.sagepub.com/thematicanalysis/
  student-resources/chapter-3/understanding-the-process (this website
  has a lot of really useful resources - the book is available through
  the library)
\end{itemize}

\textbf{Are we meant to be using literature to support our reflexivity
like we did for the group project reflection?} You don't need to do this
- reflexivity for a specific research project is more about thinking
about your own positionality and how that might have affected the way
you approached your data so it doesn't need to be supported with
evidence.

\textbf{Do we need to report the specific questions that we used to
extract our results? I mean of course I read the full 3 interviews,
however my topic was very specific and it was located in two specific
questions in the interviews. Do I need to report that? Furthermore do I
need to report all the questions of the interview anyways? I see all the
qualitative studies report the full set of the questions.} You can add
them as an appendix but you don't need to list them all in the methods
section; there you can just describe generally what kind of questions
they were asked.

\subsection{Analysis}\label{analysis}

\textbf{According to Braun and Clarke (2013) a good analysis section
will link back to relevant research previously done on the topic (i.e
citations to link it to the wider literature) however, I would presume
that this will be better suited to the discussion section of our project
given the word count? Or should we cite previous research to link our
themes within the analysis section?} Qual reports are written in one of
two ways: 1) where the analysis and discussion are in separate sections
- in the analysis, you report your findings, in the discussion, you link
your findings to the literature, identify limitations etc. 2) you have a
joint analysis and discussion section, where you both present your
results and link them to the research. Braun \& Clarke are talking about
the second, whereas we're doing the first option. So, no, just focus on
presenting your themes in the analysis section - like with the quant
report, you will then have a separate Discussion section where you link
to the literature

\textbf{In the thematic analysis, what should we do if two participants
give opposite answers for the same questions. How do we deal with that
conflicts?} This can make for a really interesting analysis and
discussion section! You should explain the theme for the reader, and
provide context for the different views presented by the participants.
You can find more information in the Braun \& Clarke (2013) book

\textbf{In the interview transcripts came across a nice statement I
would like to quote, however the statement was made by the interviewer
not the participant. The participant agreed with the interviewers
statement. Would it be appropriate to quote the interviewer in this
instance?} It might depend on what the statement is and the context. So,
for example, is it a) the researcher summarising what the participant
said and then the participant agreeing or b) the researcher giving their
own opinion about it and the participant agreeing? I'd be more likely to
say this is okay to include if it was the first than the second, I
think. If you were to include it, it would probably be best to a) make
sure that there's enough context surrounding this (i.e.~what was there
in the lead up to it and also afterwards?), and b) only do this
sparingly.

\textbf{Can one of my themes be my RQ or an interview question?} No, you
should try to go a bit deeper in finding meanings in the data in your
analysis. If you end up with a theme that is actually your RQ or one of
the interview questions, you have \emph{summarised} the data rather than
\emph{analysing} it. This is surprisingly easy to do by accident so
don't feel bad if this has happened to you - just go back to the data
and look a bit further. What aspects of the RQ does the theme capture?
What specifically do the participants say about it?

\textbf{I started exploring the data without any previous thoughts
(inductive) and later I started reading the literature for this topic
and unintentionally I started observing some of the same conclusions in
my data (deductive). Is this a formal way of analysing the data? Can I
report a mixture of approaches?} It is possible to do this -
\href{https://journals.sagepub.com/doi/10.1177/160940690600500107\#}{here
is an example of a paper} that used a combination of inductive and
deductive coding for TA. This would not be the only way to do it, but is
an overview of one possible way in which coding and developing themes
could utilise both. My main suggestion would be to be clear with the
reader in your methods about the thinking behind your approach.

\textbf{What if in our analysis we end up developing different themes
the three interviews? Do we have to report only the common themes or is
it ok if we also analyse themes that are evident only in one or two
interviews?} Ideally at the themes stage you should consider the dataset
as a whole so most of your themes should come from more than one
interview Try to find quotes which illustrate the theme not just from
one participant. Sometimes themes may be more strongly present in parts
of the data and that's ok but try to avoid having one theme that's from
one participant only and another one that's from another participant -
the themes should represent the data a bit more holistically

\textbf{What if I develop more than two themes - can I mention the other
themes in my analysis section?} You can briefly say that you developed
more themes but don't waste your word count on explaining them in more
detail or talking about the codes related to them (you can talk about
them in your submission 2 document where you evidence all the stages. In
terms of the discussion, you should focus on the two themes you have
chosen and developed a full narrative for (the same way in quant you
would not discuss an analysis you did not actually finish) and connect
those to existing literature. You can talk about the other stuff as
future directions if you would like.

\textbf{Should the themes be presented in a hierarchical way related
with each other? }

\textbf{What if we detect themes that are not related with each other?
Do we include data that don't fit anywhere else or if they are trivial
we can exclude them from the final} It depends on how you identify your
themes - they may be strongly related to each other and hierarchical and
they may not. There is no one correct way. In the report writing guide,
we ask that you report a maximum of two themes (or one theme with 2-3
subthemes) so I would choose those where you would be able to go into
the most depth in terms of interpretation report?

\#\#\#Appendices

\textbf{What should I include in the reflexive diary?} Please see a
detailed template with probing questions here:
\href{https://gla-my.sharepoint.com/:w:/g/personal/ashley_robertson_glasgow_ac_uk/IQC1PIxhBlndR74S_nfeKHNlAW5uQmGOpoRXCpLHC2fI_fI?e=aunShJ}{reflexivity
diary}.

\cleardoublepage
\phantomsection
\addcontentsline{toc}{part}{Appendices}
\appendix

\chapter{Qualitative Research Showcases}\label{appendix_a_showcase}

Throughout RM2, we will showcase some exciting qualitative research
which will go beyond what you are taught on the course. For each
showcase, we'll provide a couple of research papers for you to look at.
You don't need to read these in full (unless you find it interesting and
want to!) - the purpose is to give you a little taster of the
multifaceted and creative things that can be achieved with qualitative
methods. We hope that you find these papers interesting!

\section{Qualitative research showcase 1 - Alternative data collection
methods}\label{qualitative-research-showcase-1---alternative-data-collection-methods}

On RM2, we mainly focus on ``standard'' data collection methodologies
(interviews, focus groups, qualitative surveys), but much of published
qualitative research uses more creative methods, such as art-based
approaches (music or photo elicitation, storytelling, collages, etc.),
participant-provided diaries, or data scraping. Arts-based approaches
allow for data collection which is less formal than an interview/focus
group situation might be and they give the participant space to express
themselves more freely. Diary studies allow for data collection over
time without overburdening the participants, while data scraping is very
sueful when working with sensitive topics and/or minority populations
that are hard to reach or overburdened by research participation
requests.

\begin{itemize}
\tightlist
\item
  Dr Gail Prasad's study explores children's views of plurilingualism
  through collage-creation.
  \href{https://www.tandfonline.com/doi/full/10.1080/13670050.2017.1420033}{Access
  the study}
\item
  Dr December-Maxwell and colleagues investigate involuntary celibacy
  through Reddit data.
  \href{https://link.springer.com/article/10.1007/s12119-020-09724-6}{Access
  the study}
\end{itemize}

\textbf{Dissertation hint:} Data scraping from Reddit (this can be done
with R) has been used in previous dissertation projects where primary
data collection would be too difficult due to the topic/population, or
where a student has wanted to work with secondary data but the topic
doesn't have open data available in research databases - keep this in
mind for next year!

\section{Qualitative Research Showcase 2 - Participatory
Research}\label{qualitative-research-showcase-2---participatory-research}

Participatory research refers to the involvement of the groups that are
affected by the issue that is being studied. It is becoming increasingly
more utilised in qualitative research, and it is considered good
practice particularly in the areas where the studied group is
disadvantaged or marginalised. The underlying philosophy of
participatory research is to see the people studied not as ``subjects''
to research, but actively participating stakeholders.

\begin{itemize}
\tightlist
\item
  Drs Vaughn and Jacquez have written an excellent review of the
  different types and core principles of participatory research.
  \href{https://jprm.scholasticahq.com/article/13244-participatory-research-methods-choice-points-in-the-research-process}{Access
  the article}
\item
  Drs Yaros and Schueller conducted a study on positive computing
  technologies for children with children as co-creators of the
  research. \href{https://www.jmir.org/2017/1/e14/}{Access the article}
\item
  Dr Crompton and colleagues studied peer support of autistic young
  adults in schools with a research team that directly involves autistic
  stakeholders. Pay particular attention to the Community Involvement
  section of the article.
  \href{https://journals.sagepub.com/doi/full/10.1177/13623613221081189}{Access
  the article}
\end{itemize}

\textbf{Dissertation hint:} If you want to work with marginalised or
disadvantaged populations for your dissertation, have a conversation
early on with your supervisor about how you might want to incorporate
participatory elements into your study. Due to the limited timescale,
for dissertations this often looks less fully fledged than what you
would see in published papers. For example, for a dissertation, the
researcher student might have lived experience themselves, or they might
reach out to University networks to find a person with lived experience
to review the research question and/or analysis.

\section{Qualitative Research Showcase 3 - Interesting Analysis
Methods}\label{qualitative-research-showcase-3---interesting-analysis-methods}

In this qualitative research showcase, we highlight a couple of
different types of analyses to give a (very small!) snapshot of the
variety of qualitative techniques. Please bear in mind these are beyond
what is expected from you on RM2 - don't try to apply these analyses to
your report.

\begin{itemize}
\tightlist
\item
  Dr Wilson and colleagues conducted a meta-synthesis of qualitative
  studies on older adulthood. This is a way to bring together findings
  from multiple qualitative studies - have a look and see what
  differences and similarities between this and quantitative
  meta-analysis.
  \href{https://www.cambridge.org/core/journals/ageing-and-society/article/abs/developing-a-model-of-resilience-in-older-adulthood-a-qualitative-metasynthesis/AE8457482A9371FF74F391E97E43800F}{Access
  the article}
\item
  Drs Goodman and Burke conducted a discourse analysis on attitudes to
  immigrants. Compare and contrast how the analysis is presented and how
  the quotes are used as opposed to a Thematic Analysis.
  \href{https://journals.sagepub.com/doi/abs/10.1177/0957926509360743}{Access
  the paper}
\end{itemize}

\textbf{Dissertation hint:} If you are doing a qualitative dissertation,
most supervisors will encourage you to do Thematic Analysis because it's
highly structured and there is a lot of guidance available. If you are
interested in more advanced qualitative analyses, make sure that you
have a discussion about this when talking to potential supervisors to
identify whether they support this.




\end{document}
